\documentclass[11pt]{book}

\usepackage[a4paper,margin=30mm]{geometry}
\usepackage{amsmath,amssymb,amsthm}
\usepackage[T1]{fontenc}
\usepackage{lmodern}
\usepackage{microtype}
\usepackage[hidelinks]{hyperref}

\newtheorem{theorem}{Theorem}[section]
\newtheorem{lemma}[theorem]{Lemma}
\newtheorem{proposition}[theorem]{Proposition}
\newtheorem{corollary}[theorem]{Corollary}
\theoremstyle{definition}
\newtheorem{algorithm}[theorem]{Algorithm}
\newtheorem{definition}[theorem]{Definition}
\newtheorem{remark}[theorem]{Remark}
\newtheorem{example}[theorem]{Example}

\newcommand{\ZZ}{\mathbb Z}
\newcommand{\len}{\operatorname{len}}
\newcommand{\bits}{\operatorname{bits}}
\newcommand{\MZ}{\mathsf M_{\ZZ}}

\title{Algorithms}
\author{}
\date{}

\begin{document}
\maketitle
\tableofcontents

\chapter{Integer polynomials}

\section{Ordinary full product}

Let
\[
 f(x)=\sum_{i=0}^{m-1} a_i x^i,
 \qquad
 g(x)=\sum_{j=0}^{n-1} b_j x^j
\]
be polynomials in $\ZZ[x]$.  Their \emph{lengths} are $m=\len(f)$ and
$n=\len(g)$, with the zero polynomial assigned length zero.  Put
\[
 N=\max(m,n), \qquad s=\min(m,n).
\]
When only a length parameter is displayed in a complexity bound, the operands
are assumed to be balanced, so that $m,n=\Theta(N)$.  Unequal operands may
always be padded with leading zero coefficients, although an implementation
need not actually perform this padding.

For $z\in\ZZ$, define
\[
 \bits(z)=
 \begin{cases}
  0, & z=0,\\
  \lfloor\log_2|z|\rfloor+1, & z\ne0,
 \end{cases}
\]
and suppose initially that $\bits(a_i),\bits(b_j)\leq\tau$.  Thus every
input coefficient has absolute value strictly less than $2^\tau$.  We write
$\MZ(t)$ for the bit complexity of multiplying two integers of at most $t$
bits.  We assume only the standard monotonicity and $\MZ(t)=\Omega(t)$.
Additions and subtractions of $t$-bit integers cost $O(t)$ bit operations.

The coefficient $c_k$ of $x^k$ in $fg$ is a sum of at most $s$ products
$a_i b_j$.  Consequently
\begin{equation}\label{eq:product-coefficient-bound}
 |c_k| < s2^{2\tau},
 \qquad
 \bits(c_k)\leq 2\tau+\lceil\log_2 s\rceil+1.
\end{equation}
This elementary bound will be used repeatedly.

The algorithms below represent four rather different ways to obtain the full
product.  Classical multiplication computes the convolution directly.
Karatsuba and Toom--Cook reduce one product to several shorter polynomial
products.  Kronecker substitution instead encodes the entire polynomial as one
large integer and delegates the multiplication to an integer algorithm.  The
historical origins of the subquadratic methods are Karatsuba--Ofman
\cite{KaratsubaOfman1962} and Toom \cite{Toom1963}; Cook's thesis
\cite{Cook1966} developed the interpolation viewpoint further.  Brent and
Zimmermann \cite{BrentZimmermann2010} give a modern account of Karatsuba and
Toom--Cook multiplication, while Bodrato \cite{Bodrato2007} treats Toom--Cook
specifically for polynomial rings.

\subsection{Classical multiplication}

The identity
\begin{equation}\label{eq:convolution}
 fg=\sum_{k=0}^{m+n-2}
 \left(\sum_{\substack{0\leq i<m,\ 0\leq j<n\\i+j=k}} a_i b_j\right)x^k
\end{equation}
immediately gives the classical, or schoolbook, algorithm.

\begin{algorithm}[Classical full product]\label{alg:classical-product}
Input $f=\sum_{i=0}^{m-1}a_ix^i$ and
$g=\sum_{j=0}^{n-1}b_jx^j$.
\begin{enumerate}
 \item Set $c_0,\ldots,c_{m+n-2}\leftarrow0$.
 \item For $i=0,\ldots,m-1$ and $j=0,\ldots,n-1$, set
 \[
    c_{i+j}\leftarrow c_{i+j}+a_i b_j.
 \]
 \item Return $\sum_k c_kx^k$.
\end{enumerate}
\end{algorithm}

\begin{theorem}\label{thm:classical-correctness-complexity}
Algorithm~\ref{alg:classical-product} returns $fg$.  It uses exactly $mn$
coefficient multiplications and $mn$ coefficient accumulations.  For balanced
inputs this is $\Theta(N^2)$ coefficient operations.  If the input
coefficients have at most $\tau$ bits, its bit complexity is
\[
 O\!\left(mn\bigl(\MZ(\tau)+\tau+\log s\bigr)\right),
\]
and hence
\[
 O\!\left(N^2\bigl(\MZ(\tau)+\tau+\log N\bigr)\right)
\]
for balanced operands.
\end{theorem}

\begin{proof}
For fixed $i,j$, the product $a_i b_j$ contributes to exactly the coefficient
of $x^{i+j}$.  Thus after all pairs have been visited, $c_k$ is precisely the
inner sum in \eqref{eq:convolution}, proving correctness.

There are $mn$ pairs $(i,j)$, which proves the coefficient-operation count.
Each product $a_i b_j$ costs $O(\MZ(\tau))$ bit operations.  During an
accumulation, the destination has at most $O(\tau+\log s)$ bits by
\eqref{eq:product-coefficient-bound}, so the addition costs
$O(\tau+\log s)$.  Summing these costs over all $mn$ pairs gives the stated
bound.
\end{proof}

\subsection{Karatsuba multiplication}

Karatsuba multiplication replaces four half-size products by three.  Choose a
split point $h$, normally $h=\lceil N/2\rceil$, and write
\[
 f=f_0+x^h f_1,
 \qquad
 g=g_0+x^h g_1,
\]
where $f_0,g_0$ have length at most $h$.  Define
\[
 z_0=f_0g_0,
 \qquad
 z_2=f_1g_1,
 \qquad
 z_1=(f_0+f_1)(g_0+g_1).
\]
Then
\begin{equation}\label{eq:karatsuba-identity}
 fg=z_0+x^h(z_1-z_0-z_2)+x^{2h}z_2.
\end{equation}
Only the three products $z_0,z_1,z_2$ are recursive multiplications.

\begin{algorithm}[Karatsuba full product]\label{alg:karatsuba-product}
Input $f,g\in\ZZ[x]$ of length at most $N$.
\begin{enumerate}
 \item If $N$ is below a chosen base-case threshold, use classical
 multiplication.
 \item Put $h\leftarrow\lceil N/2\rceil$ and split
 \[
   f=f_0+x^hf_1,\qquad g=g_0+x^hg_1.
 \]
 \item Recursively compute
 \[
   z_0=f_0g_0,\qquad z_2=f_1g_1,
   \qquad z_1=(f_0+f_1)(g_0+g_1).
 \]
 \item Return the right-hand side of \eqref{eq:karatsuba-identity}.
\end{enumerate}
\end{algorithm}

\begin{theorem}\label{thm:karatsuba-correctness}
Algorithm~\ref{alg:karatsuba-product} returns $fg$.
\end{theorem}

\begin{proof}
Expanding $z_1$ gives
\[
 z_1=z_0+f_0g_1+f_1g_0+z_2.
\]
Hence $z_1-z_0-z_2=f_0g_1+f_1g_0$.  Substituting this into
\eqref{eq:karatsuba-identity} gives
\[
 f_0g_0+x^h(f_0g_1+f_1g_0)+x^{2h}f_1g_1
 =(f_0+x^hf_1)(g_0+x^hg_1)=fg.
\]
The recursive calls are correct by induction on the input length, with the
classical base case providing the induction base.
\end{proof}

\begin{theorem}\label{thm:karatsuba-complexity}
For balanced length-$N$ inputs, recursive Karatsuba multiplication uses
\[
 \Theta\!\left(N^{\log_2 3}\right)
\]
coefficient operations.  If the input coefficients have at most $\tau$ bits,
a bit-complexity bound is
\[
 O\!\left(
 N^{\log_2 3}\,\MZ(\tau+\log N)
 \right).
\]
\end{theorem}

\begin{proof}
Ignoring rounding at the split point, the three recursive products have length
$N/2$, while splitting, additions, subtractions and recombination require
$O(N)$ coefficient operations.  Thus
\[
 T(N)=3T(N/2)+O(N),
\]
which gives $T(N)=\Theta(N^{\log_2 3})$ by the Master Theorem.  This
analyzes recursive use of Karatsuba itself; replacing recursive products by an
asymptotically no slower multiplication method preserves the upper bound.

For the bit bound, consider a path of depth $d$ in the recursion tree.
Whenever the middle product is chosen, an input coefficient may be replaced by
a sum of two coefficients from the preceding level.  Thus a coefficient at
depth $d$ is a signed sum of at most $2^d$ original coefficients, and has at
most $\tau+d+O(1)$ bits.  Since the depth is $O(\log N)$, every integer
multiplication at a leaf has operands of $O(\tau+\log N)$ bits.

At level $d$ there are $3^d$ subproblems of length $O(N/2^d)$.  Hence the
total number of coefficient additions at that level is
$O(N(3/2)^d)$, and summing over all levels gives
$O(N^{\log_2 3})$.  Each such addition involves
$O(\tau+\log N)$ bits.  Since $\MZ(t)=\Omega(t)$, the additions are absorbed by
$O(N^{\log_2 3}\MZ(\tau+\log N))$, as are the leaf multiplications.
\end{proof}

\subsection{Three-way Toom--Cook multiplication}

Three-way Toom--Cook, or Toom-3, splits each operand into three blocks.  Let
$h=\lceil N/3\rceil$, put $X=x^h$, and write
\[
 f=f_0+f_1X+f_2X^2,
 \qquad
 g=g_0+g_1X+g_2X^2,
\]
where each block has length at most $h$.  Introduce an auxiliary variable $t$
and define
\[
 F(t)=f_0+f_1t+f_2t^2,
 \qquad
 G(t)=g_0+g_1t+g_2t^2.
\]
Their product has the form
\[
 H(t)=F(t)G(t)=c_0+c_1t+c_2t^2+c_3t^3+c_4t^4.
\]
The desired polynomial is
\begin{equation}\label{eq:toom-recompose}
 fg=H(X)=c_0+c_1X+c_2X^2+c_3X^3+c_4X^4.
\end{equation}
Thus it is enough to recover the five coefficient-polynomials $c_i$.

The evaluation points used here are
\[
 0,\quad 1,\quad -1,\quad 2,\quad \infty,
\]
where evaluation at infinity means taking the leading coefficient in $t$.
Set
\[
 v_0=H(0),\quad v_1=H(1),\quad v_{-1}=H(-1),
 \quad v_2=H(2),\quad v_\infty=c_4.
\]
Each $v$ is obtained by one recursive product, so there are five products of
roughly one third the original length.

\begin{algorithm}[Three-way Toom--Cook full product]\label{alg:toom3-product}
Input $f,g\in\ZZ[x]$ of length at most $N$.
\begin{enumerate}
 \item If $N$ is below a chosen base-case threshold, use a simpler
 multiplication algorithm.
 \item Put $h\leftarrow\lceil N/3\rceil$, $X=x^h$, and split
 \[
   f=f_0+f_1X+f_2X^2,\qquad
   g=g_0+g_1X+g_2X^2.
 \]
 \item Evaluate $F$ and $G$ at $0,1,-1,2,\infty$ and recursively multiply
 corresponding values to obtain
 $v_0,v_1,v_{-1},v_2,v_\infty$.
 \item Interpolate by the sequence
 \begin{align*}
   r_3 &\leftarrow v_2-v_1-15v_\infty,\\
   r_1 &\leftarrow (v_1-v_{-1})/2,\\
   r_2 &\leftarrow v_{-1}-v_0-v_\infty+r_1,\\
   r_3 &\leftarrow (r_3-r_1-3r_2)/6,\\
   r_1 &\leftarrow r_1-r_3.
 \end{align*}
 \item Set
 \[
    c_0=v_0,\quad c_1=r_1,\quad c_2=r_2,
    \quad c_3=r_3,\quad c_4=v_\infty,
 \]
 and return \eqref{eq:toom-recompose}.
\end{enumerate}
\end{algorithm}

\begin{lemma}\label{lem:toom3-interpolation}
The interpolation step in Algorithm~\ref{alg:toom3-product} recovers the
coefficients $c_0,\ldots,c_4$ exactly.  In particular, the indicated divisions
by $2$ and $6$ are exact in $\ZZ[x]$.
\end{lemma}

\begin{proof}
Expanding the five evaluations gives
\begin{align*}
 v_0&=c_0,\\
 v_1&=c_0+c_1+c_2+c_3+c_4,\\
 v_{-1}&=c_0-c_1+c_2-c_3+c_4,\\
 v_2&=c_0+2c_1+4c_2+8c_3+16c_4,\\
 v_\infty&=c_4.
\end{align*}
After the first assignment,
\[
 r_3=c_1+3c_2+7c_3.
\]
The second assignment is exact because
$v_1-v_{-1}=2(c_1+c_3)$, and it yields
$r_1=c_1+c_3$.  Therefore
\[
 r_2=v_{-1}-v_0-v_\infty+r_1=c_2.
\]
Next,
\[
 r_3-r_1-3r_2
 =(c_1+3c_2+7c_3)-(c_1+c_3)-3c_2=6c_3.
\]
so the division by $6$ is exact and leaves $r_3=c_3$.  Finally
$r_1\leftarrow r_1-r_3$ gives $r_1=c_1$.  The remaining two coefficients were
already $c_0=v_0$ and $c_4=v_\infty$.
\end{proof}


\subsubsection{Optimality of the interpolation denominator}

The exact divisions in Toom--Cook interpolation are not merely an artefact of
one choice of formulas.  For integer evaluation points there is an intrinsic
arithmetic obstruction.  It is convenient to make this precise in terms of the
interpolation matrix.

For four distinct finite integers $a_0,a_1,a_2,a_3$, together with the point
$\infty$, let
\[
 E(a_0,a_1,a_2,a_3)=
 \begin{pmatrix}
  1&a_0&a_0^2&a_0^3&a_0^4\\
  1&a_1&a_1^2&a_1^3&a_1^4\\
  1&a_2&a_2^2&a_2^3&a_2^4\\
  1&a_3&a_3^2&a_3^3&a_3^4\\
  0&0&0&0&1
 \end{pmatrix}.
\]
Thus $E(c_0,\ldots,c_4)^T$ is the vector of values of
$c_0+c_1t+\cdots+c_4t^4$ at these five points.  Define the
\emph{interpolation denominator} of $E$ to be the least positive integer
$D(E)$ such that $D(E)E^{-1}$ has integer entries.

\begin{theorem}\label{thm:toom3-denominator-lower-bound}
For every choice of four distinct finite integer evaluation points,
\[
 6\mid D(E).
\]
For the points $0,1,-1,2,\infty$ used in
Algorithm~\ref{alg:toom3-product}, one has $D(E)=6$.  Hence this choice is
optimal with respect to the interpolation denominator among all Toom-3
schemes using four finite integer points and $\infty$.
\end{theorem}

\begin{proof}
Let $p$ be either $2$ or $3$.  There are four finite evaluation points but
only $p$ residue classes modulo $p$.  Hence two of the $a_i$ are congruent
modulo $p$.  The corresponding two rows of $E$ are then equal modulo $p$, so
$E$ is singular over $\mathbb F_p$.

Suppose that $p\nmid D(E)$.  Since $C=D(E)E^{-1}$ is an integer matrix,
\[
 CE=D(E)I.
\]
Reducing modulo $p$, the scalar $D(E)$ is invertible.  Thus
$D(E)^{-1}C$ would be a left inverse of $E$ over $\mathbb F_p$, contradicting
the singularity of $E$.  Therefore $p\mid D(E)$ for both $p=2$ and $p=3$,
and hence $6\mid D(E)$.

For the particular ordering $0,1,-1,2,\infty$, the evaluation matrix is
\[
 E=
 \begin{pmatrix}
 1&0&0&0&0\\
 1&1&1&1&1\\
 1&-1&1&-1&1\\
 1&2&4&8&16\\
 0&0&0&0&1
 \end{pmatrix},
\]
and direct inversion gives
\[
 E^{-1}=
 \begin{pmatrix}
 1&0&0&0&0\\
 -\frac12&1&-\frac13&-\frac16&2\\
 -1&\frac12&\frac12&0&-1\\
 \frac12&-\frac12&-\frac16&\frac16&-2\\
 0&0&0&0&1
 \end{pmatrix}.
\]
All denominators divide $6$, so $D(E)\leq6$.  The first part gives the
reverse divisibility, and therefore $D(E)=6$.
\end{proof}

\begin{corollary}\label{cor:toom3-no-shifts-only}
No linear interpolation scheme for Toom-3 using four finite integer points and
$\infty$ can recover the five product blocks using only additions,
subtractions, multiplication by integers, and exact divisions by powers of
$2$.
\end{corollary}

\begin{proof}
Any straight-line expression built from those operations is a linear form in
the five evaluated products whose rational coefficients have denominators
that are powers of $2$.  Consequently a complete shifts-only interpolation
would give an inverse evaluation matrix with no factor $3$ in its
interpolation denominator.  This contradicts
Theorem~\ref{thm:toom3-denominator-lower-bound}.
\end{proof}

\begin{remark}
Theorem~\ref{thm:toom3-denominator-lower-bound} concerns the common
denominator of the inverse interpolation map, rather than the cost of any
particular straight-line program.  For the present evaluation points the
factor $6$ is therefore intrinsic, although different interpolation sequences
may distribute its factors among several exact divisions.  Stronger general
lower bounds on the number of prime divisions required by integer-point
Toom--Cook interpolation are proved by Nissim, Schwartz and Spiizer
\cite{NissimSchwartzSpiizer2026}.
\end{remark}

\begin{theorem}\label{thm:toom3-correctness-complexity}
Algorithm~\ref{alg:toom3-product} returns $fg$.  For balanced length-$N$
inputs it uses
\[
 \Theta\!\left(N^{\log_3 5}\right)
\]
coefficient operations.  If all input coefficients have at most $\tau$ bits,
its bit complexity is
\[
 O\!\left(
 N^{\log_3 5}\,\MZ(\tau+\log N)
 \right).
\]
\end{theorem}

\begin{proof}
Lemma~\ref{lem:toom3-interpolation} recovers the coefficient-polynomials of
$H(t)=F(t)G(t)$.  Substituting $t=X=x^h$ then gives $fg$ by
\eqref{eq:toom-recompose}, proving correctness inductively from the base case.

Evaluation and interpolation use only a constant number of coefficientwise
additions, subtractions, multiplications by small integers, and exact divisions
by the constants $2$ and $6$.  Their cost is therefore $O(N)$ coefficient
operations at each node.  There are five recursive products of length at most
$\lceil N/3\rceil$, so
\[
 T(N)=5T(N/3)+O(N),
\]
and hence $T(N)=\Theta(N^{\log_3 5})$.  As with Karatsuba, recursive
products may instead be delegated to any asymptotically no slower method
without weakening this upper bound.

For bit complexity, every evaluation or interpolation coefficient is obtained
from coefficients at the preceding recursion level using a fixed finite set of
small integer linear combinations.  Thus one recursion level increases the
coefficient bit size by only $O(1)$.  After $O(\log N)$ levels all operands of
base-case integer multiplications have $O(\tau+\log N)$ bits.
At level $d$ the total number of coefficient positions processed is
$O(N(5/3)^d)$; summing over the recursion gives
$O(N^{\log_3 5})$.  Linear-time coefficient operations are again absorbed by
$\MZ(\tau+\log N)$, giving the claimed bit bound.
\end{proof}

\begin{remark}
The exponent
\[
 \log_3 5\approx1.465
\]
is smaller than Karatsuba's $\log_2 3\approx1.585$.  This does not imply that
Toom-3 is preferable for all finite sizes: it performs more evaluation,
interpolation and temporary arithmetic, so practical multiplication routines
normally switch between several algorithms.
\end{remark}

\subsection{Kronecker substitution}

Kronecker substitution converts multiplication in $\ZZ[x]$ into one integer
multiplication.  We assume here that both operands are nonzero; a zero operand is
handled immediately.  For signed coefficients it is convenient to use a balanced
base expansion.  Let
\[
 \tau_f=\max_i\bits(a_i),
 \qquad
 \tau_g=\max_j\bits(b_j),
\]
and let $s=\min(m,n)$.  Choose
\begin{equation}\label{eq:ks-width}
 k=\tau_f+\tau_g+\lceil\log_2 s\rceil+1,
 \qquad B=2^k.
\end{equation}
Then every coefficient of $fg$ has absolute value strictly less than $B/2$.

Evaluate the operands at $B$:
\[
 A=f(B)=\sum_i a_iB^i,
 \qquad
 D=g(B)=\sum_j b_jB^j,
\]
and compute the integer product $C=AD$.  Algebraically,
\begin{equation}\label{eq:ks-packed-product}
 C=(fg)(B)=\sum_{r=0}^{m+n-2}c_rB^r.
\end{equation}
Because the true coefficients lie in the centered interval
$(-B/2,B/2)$, they can be recovered one at a time as balanced base-$B$ digits.

For an integer $z$, write $\operatorname{cres}_B(z)$ for the unique integer
$r$ satisfying
\[
 r\equiv z\pmod B,
 \qquad -B/2\leq r<B/2.
\]
The strict coefficient bound ensures that no coefficient is an ambiguous
endpoint.

\begin{algorithm}[Balanced Kronecker substitution]\label{alg:kronecker-product}
Input $f,g\in\ZZ[x]$.
\begin{enumerate}
 \item Compute $k$ from \eqref{eq:ks-width} and put $B=2^k$.
 \item Pack
 \[
    A\leftarrow\sum_i a_iB^i,
    \qquad D\leftarrow\sum_j b_jB^j.
 \]
 \item Compute the single integer product $C\leftarrow AD$.
 \item For $r=0,\ldots,m+n-2$ do:
 \begin{enumerate}
  \item $c_r\leftarrow\operatorname{cres}_B(C)$;
  \item $C\leftarrow(C-c_r)/B$.
 \end{enumerate}
 \item Return $\sum_r c_rx^r$.
\end{enumerate}
\end{algorithm}

The use of $B=2^k$ makes packing, centered residue extraction and division by
$B$ bit-shift operations.  Standard Kronecker substitution and refinements are
discussed by Harvey \cite{Harvey2009}; in particular, that paper emphasizes
the reduction of polynomial multiplication to highly optimized integer
multiplication.

\begin{lemma}\label{lem:ks-bound}
With $k$ chosen by \eqref{eq:ks-width}, every coefficient $c_r$ of $fg$
satisfies $|c_r|<B/2$.
\end{lemma}

\begin{proof}
Each $c_r$ is a sum of at most $s$ products.  Since
$|a_i|<2^{\tau_f}$ and $|b_j|<2^{\tau_g}$,
\[
 |c_r|<s2^{\tau_f+\tau_g}
 \leq 2^{\tau_f+\tau_g+\lceil\log_2s\rceil}
 =2^{k-1}=B/2.
\]
\end{proof}

\begin{theorem}\label{thm:ks-correctness}
Algorithm~\ref{alg:kronecker-product} returns $fg$.
\end{theorem}

\begin{proof}
Equation~\eqref{eq:ks-packed-product} holds because evaluation at the integer
$B$ is a ring homomorphism $\ZZ[x]\to\ZZ$.  By
Lemma~\ref{lem:ks-bound}, $c_0$ lies in $(-B/2,B/2)$, so it is the unique
centered residue of $C$ modulo $B$.  Hence the first unpacking step recovers
$c_0$ exactly.  Subtracting $c_0$ and dividing by $B$ gives
\[
 \frac{C-c_0}{B}=c_1+c_2B+\cdots.
\]
The same argument now recovers $c_1$.  Induction recovers every coefficient in
turn.
\end{proof}

\begin{theorem}\label{thm:ks-complexity}
Let the input lengths be at most $N$ and the coefficient bit sizes at most
$\tau$.  Balanced Kronecker substitution uses one integer multiplication of
operands having
\[
 O\bigl(N(\tau+\log N)\bigr)
\]
bits.  With linear-time power-of-two packing and unpacking, its bit complexity
is
\[
 O\!\left(
 \MZ\bigl(N(\tau+\log N)\bigr)
 \right).
\]
In particular, for coefficients of bounded bit size this becomes
\[
 O\bigl(\MZ(N\log N)\bigr),
\]
while if $\tau=\Omega(\log N)$ it is
$O(\MZ(N\tau))$.
\end{theorem}

\begin{proof}
From \eqref{eq:ks-width}, $k=O(\tau+\log N)$.  A packed operand contains at
most $N$ base-$2^k$ digits, so its bit length is $O(Nk)$, namely
$O(N(\tau+\log N))$.  The integer product therefore has cost
\[
 O\!\left(\MZ\bigl(N(\tau+\log N)\bigr)\right).
\]
Because the radix is a power of two, packing and unpacking can be performed by
shifts, masks and copies touching only $O(Nk)$ bits in total.  This linear term
is absorbed by $\MZ(Nk)$ under the standing assumption
$\MZ(t)=\Omega(t)$.  Substituting $\tau=O(1)$ or
$\tau=\Omega(\log N)$ gives the two special cases.
\end{proof}

\begin{remark}
The four algorithms optimize different resources.  Classical multiplication
has the smallest setup cost.  Karatsuba and Toom-3 reduce the number of
coefficient multiplications while remaining purely polynomial algorithms.
Kronecker substitution can be superior when the coefficient sizes and lengths
make one packed integer multiplication attractive.  A practical dispatcher may
therefore use all four rather than regard any one of them as universally best.
\end{remark}

\subsection{Polynomial powering}

For $f\in\ZZ[x]$ and a nonnegative integer $e$, polynomial powering asks for
$f^e$.  We use the convention $f^0=1$, including $0^0=1$.  Repeated
multiplication needs $e-1$ products, whereas binary powering reduces this to
$O(\log e)$ products.  The method is the ordinary left-to-right binary
exponentiation algorithm; see, for example, \cite{BrentZimmermann2010}.

\begin{algorithm}[Binary polynomial powering]\label{alg:binary-poly-power}
Input $f\in\ZZ[x]$ and $e\ge0$.
\begin{enumerate}
 \item If $e=0$, return $1$.
 \item Write
 \[
   e=(1b_{\ell-2}\cdots b_1b_0)_2,
   \qquad \ell=\lfloor\log_2 e\rfloor+1,
 \]
 and put $r\leftarrow f$.
 \item For $j=\ell-2,\ldots,0$, set $r\leftarrow r^2$ and, if $b_j=1$,
 set $r\leftarrow rf$.
 \item Return $r$.
\end{enumerate}
\end{algorithm}

\begin{theorem}\label{thm:binary-poly-power}
Algorithm~\ref{alg:binary-poly-power} returns $f^e$.  For $e\ge1$, if
$w(e)$ denotes the number of $1$-bits of $e$, it uses exactly
\[
  \ell-1\quad\hbox{squarings},\qquad w(e)-1\quad\hbox{other products},
\]
and hence at most $2\lfloor\log_2 e\rfloor$ polynomial multiplications.
If $f\ne0$ has degree $d$, then
\[
 \deg(f^e)=ed.
\]
\end{theorem}

\begin{proof}
After processing a prefix of the binary expansion of $e$, let $q$ be the
integer represented by that prefix.  The loop invariant is $r=f^q$.
Appending a zero bit replaces $q$ by $2q$ and the squaring gives
$r=f^{2q}$.  Appending a one bit replaces $q$ by $2q+1$ and the following
multiplication by $f$ gives $r=f^{2q+1}$.  At termination the processed
prefix is all of $e$, proving correctness.  Every bit after the leading bit
causes one squaring, and precisely the remaining $1$-bits cause an extra
multiplication.  The degree formula follows from $\deg(uv)=\deg u+\deg v$
for nonzero polynomials.
\end{proof}

The coefficient growth is conveniently bounded using the $1$-norm
$\|f\|_1=\sum_i|a_i|$.  Since $\|fg\|_1\le\|f\|_1\|g\|_1$,
\[
 \|f^e\|_\infty\le\|f^e\|_1\le\|f\|_1^e.
\]

\begin{corollary}\label{cor:binary-poly-power-growth}
Let $f$ have length $n$ and coefficient bit size at most $\tau$.  Every
coefficient occurring in Algorithm~\ref{alg:binary-poly-power}, and every
coefficient of $f^e$, has
\[
 O\bigl(e(\tau+\log(n+1))\bigr)
\]
bits.  If $N=ed+1$ bounds the final length and $\mathsf M(N)$ denotes the
coefficient-operation cost of multiplication at length $N$, then the safe
coefficient-operation bound is
\[
 O\bigl(\log(e+1)\,\mathsf M(N)\bigr).
\]
Under the usual regularity hypothesis that multiplication costs decrease
geometrically when the operand lengths are halved, the sum over the binary
chain is $O(\mathsf M(N))$.
\end{corollary}

\begin{proof}
We have $\|f\|_1<n2^\tau$.  Every intermediate is $f^q$ for some
$0\le q\le e$, so
\[
 \bits([x^j]f^q)
 \le O\bigl(q(\tau+\log(n+1))\bigr),
\]
which gives the first claim.  The safe complexity bound follows by charging
each of the $O(\log(e+1))$ products the final multiplication cost.  Along the
binary chain the exponents, and hence the polynomial lengths, grow
geometrically.  Summing multiplication costs over these geometrically growing
lengths gives $O(\mathsf M(N))$ under the stated regularity assumption.
\end{proof}

\section{Evaluation, composition and Taylor shift}

Let
\[
 f(x)=\sum_{i=0}^{n-1}a_ix^i\in\ZZ[x].
\]
Three closely related operations will be considered.  For $c\in\ZZ$,
\emph{evaluation} computes $f(c)\in\ZZ$.  For $g\in\ZZ[x]$,
\emph{composition} computes
\[
 (f\circ g)(x)=f(g(x)).
\]
Finally, for $c\in\ZZ$, the \emph{Taylor shift} is the special composition
$f(x+c)$.  Horner's rule gives the simplest algorithm for all three.  A
balanced divide-and-conquer form, closely related to Estrin's evaluation
scheme \cite{Estrin1960}, exposes independent subexpressions and replaces a
long serial chain by products whose sizes grow geometrically.

The divide-and-conquer composition described below is deliberately elementary.
Asymptotically faster general composition algorithms based on baby-step/
giant-step decompositions and matrix multiplication are known; see
Brent--Kung \cite{BrentKung1978}.  Those methods require additional linear
algebra infrastructure and are not needed for the algorithms considered here.

\subsection{Integer evaluation by Horner's rule}

The nested identity
\begin{equation}\label{eq:horner-evaluation}
 f(c)=a_0+c\bigl(a_1+c(a_2+\cdots+c(a_{n-2}+ca_{n-1})\cdots)\bigr)
\end{equation}
gives Horner's rule.

\begin{algorithm}[Horner evaluation]\label{alg:horner-evaluation}
Input $f=\sum_{i=0}^{n-1}a_ix^i\in\ZZ[x]$ and $c\in\ZZ$.
\begin{enumerate}
 \item If $n=0$, return $0$.
 \item Put $r\leftarrow a_{n-1}$.
 \item For $i=n-2,\ldots,0$, set
 \[
   r\leftarrow rc+a_i.
 \]
 \item Return $r$.
\end{enumerate}
\end{algorithm}

\begin{theorem}\label{thm:horner-evaluation}
Algorithm~\ref{alg:horner-evaluation} returns $f(c)$.  For $n>0$ it uses
exactly $n-1$ integer multiplications and $n-1$ additions.
\end{theorem}

\begin{proof}
After coefficient $a_i$ has been processed, the invariant is
\[
 r=a_i+c\bigl(a_{i+1}+c(\cdots+ca_{n-1})\bigr).
\]
The next iteration establishes the same formula with $i$ replaced by $i-1$.
At $i=0$ this is exactly \eqref{eq:horner-evaluation}.  The operation count is
immediate from the loop.
\end{proof}

Put
\[
 \xi=\max(1,\bits(c)),
 \qquad
 B=\tau+\lceil\log_2(n+1)\rceil+(n-1)\xi,
\]
where the coefficients of $f$ have at most $\tau$ bits.  Since
\[
 |f(c)|\le n2^\tau\max(1,|c|)^{n-1},
\]
the output has $O(B)$ bits, and the same bound applies to every Horner
intermediate.

\begin{corollary}\label{cor:horner-evaluation-bit}
With the notation above, Horner evaluation has the direct bit-complexity bound
\[
 O\bigl(n\MZ(B)\bigr).
\]
This deliberately treats every multiplication as a general $B$-bit integer
product; multiplication by a small fixed $c$ can be substantially cheaper.
\end{corollary}

\subsection{Divide-and-conquer integer evaluation}

Pair adjacent coefficients rather than extending one nested expression from
right to left.  Initially put $p=c$.  The first level forms
\[
 a_0+pa_1,\quad a_2+pa_3,\quad a_4+pa_5,\ldots,
\]
then replaces $p$ by $p^2$.  The next level pairs these blocks, and so on.

\begin{algorithm}[Divide-and-conquer evaluation]\label{alg:dc-evaluation}
Input $f=\sum_{i=0}^{n-1}a_ix^i$ and $c\in\ZZ$.
\begin{enumerate}
 \item If $n=0$, return $0$.  Set $w_i\leftarrow a_i$ and $p\leftarrow c$.
 \item While more than one $w_i$ remains, replace each adjacent pair by
 \[
    w_{2i}+p w_{2i+1},
 \]
 carrying an unmatched final entry unchanged.  Then put $p\leftarrow p^2$
 unless the reduction is complete.
 \item Return the sole remaining entry.
\end{enumerate}
\end{algorithm}

\begin{lemma}\label{lem:dc-evaluation-invariant}
After $j$ pairing levels, each surviving complete block beginning at index
$i2^j$ equals
\[
 \sum_{r=0}^{2^j-1} a_{i2^j+r}c^r,
\]
with the evident truncation for the last incomplete block, and the current
power is $p=c^{2^j}$.
\end{lemma}

\begin{proof}
The assertion is trivial for $j=0$.  Suppose it holds at level $j$.  Combining
two adjacent blocks gives
\[
 \sum_{r=0}^{2^j-1}a_{i2^{j+1}+r}c^r
 +c^{2^j}\sum_{r=0}^{2^j-1}
 a_{i2^{j+1}+2^j+r}c^r,
\]
which is the required block of length $2^{j+1}$.  Squaring the current power
changes $c^{2^j}$ into $c^{2^{j+1}}$.
\end{proof}

\begin{theorem}\label{thm:dc-evaluation}
Algorithm~\ref{alg:dc-evaluation} returns $f(c)$.  It performs exactly
$n-1$ block multiplications and $n-1$ block additions, together with
\[
 \max(0,\lceil\log_2 n\rceil-1)
\]
squarings used to construct the powers $c,c^2,c^4,\ldots$.
\end{theorem}

\begin{proof}
At the final level Lemma~\ref{lem:dc-evaluation-invariant} gives one block
containing all coefficients, namely $f(c)$.  Every pairing reduces the number
of live blocks by one, so exactly $n-1$ pairings are required in total.  If
$L=\lceil\log_2 n\rceil$, there are $L$ nonempty pairing levels when $n>1$;
the power is squared after every level except the last, giving $L-1$
squarings.
\end{proof}

Although Horner and divide-and-conquer evaluation use essentially the same
number of arithmetic operations, the operand sizes occur in a different
order.  At level $j$ of Algorithm~\ref{alg:dc-evaluation}, the block size is
at most
\[
 B_j=O\bigl(\tau+2^j\xi+\log(n+1)\bigr),
\]
and there are $O(n/2^j)$ such products.  Hence a direct bit bound is
\begin{equation}\label{eq:dc-evaluation-bit}
 O\left(
   \sum_{j=0}^{\lceil\log_2 n\rceil-1}
   \frac{n}{2^j}\MZ(B_j)
   +\sum_{j=0}^{\lceil\log_2 n\rceil-2}\MZ(2^j\xi)
 \right).
\end{equation}
For a power-law integer multiplication cost
$\MZ(t)=\Theta(t^\alpha)$ with $\alpha>1$, this is
\[
 O\bigl(n\MZ(\tau+\xi)+\MZ(\tau+n\xi)\bigr).
\]
For quasi-linear integer multiplication, the level sum may carry one extra
logarithmic factor.  Thus the balanced method becomes attractive when $c$ and
the accumulated values are large, whereas Horner's rule has lower setup cost
and is preferable for short inputs or very small $c$.

\subsection{Polynomial composition by Horner's rule}

The same nested identity works with a polynomial in place of the integer
$c$:
\begin{equation}\label{eq:horner-composition}
 f(g)=a_0+g\bigl(a_1+g(a_2+\cdots+g(a_{n-2}+ga_{n-1})\cdots)\bigr).
\end{equation}

\begin{algorithm}[Horner composition]\label{alg:horner-composition}
Input $f=\sum_{i=0}^{n-1}a_ix^i$ and $g\in\ZZ[x]$.
\begin{enumerate}
 \item If $n=0$, return $0$.  Put $r\leftarrow a_{n-1}$.
 \item For $i=n-2,\ldots,0$, set
 \[
     r\leftarrow rg+a_i.
 \]
 \item Return $r$.
\end{enumerate}
\end{algorithm}

\begin{theorem}\label{thm:horner-composition}
Algorithm~\ref{alg:horner-composition} returns $f\circ g$.  It uses exactly
$n-1$ polynomial multiplications by $g$ and $n-1$ additions to the constant
coefficient.  If $d=\deg f$ and $e=\deg g>0$, then
\[
 \deg(f\circ g)=de.
\]
\end{theorem}

\begin{proof}
The proof of the nested invariant is identical to that for
Theorem~\ref{thm:horner-evaluation}, now in the ring $\ZZ[x]$.  If $g$ is
nonconstant, the leading term of $a_dg^d$ has degree $de$, whereas every
$a_i g^i$ with $i<d$ has smaller degree, so cancellation at degree $de$ is
impossible.
\end{proof}

If $g$ has length $m$, a useful generic coefficient-operation bound is
\begin{equation}\label{eq:horner-composition-cost}
 O\left(\sum_{k=1}^{n-1}\mathsf M(km,m)+n\right),
\end{equation}
where $\mathsf M(u,v)$ denotes the cost of multiplying lengths $u$ and $v$.
Padding both operands to the larger length gives the simpler, but weaker,
bound $O(n\mathsf M(nm))$.

\subsection{Divide-and-conquer polynomial composition}

Algorithm~\ref{alg:dc-evaluation} works verbatim in the coefficient ring
$\ZZ[x]$, with the scalar power $c^{2^j}$ replaced by the polynomial power
$g^{2^j}$.  This gives a balanced composition algorithm.

\begin{algorithm}[Divide-and-conquer composition]\label{alg:dc-composition}
Input $f=\sum_{i=0}^{n-1}a_ix^i$ and $g\in\ZZ[x]$.
\begin{enumerate}
 \item If $n\le1$, return $f$.  Set $w_i\leftarrow a_i$ as constant
 polynomials and $p\leftarrow g$.
 \item While more than one $w_i$ remains, replace each adjacent pair by
 \[
    w_{2i}+p w_{2i+1},
 \]
 carrying an unmatched final entry unchanged.  Then put $p\leftarrow p^2$
 unless the reduction is complete.
 \item Return the sole remaining polynomial.
\end{enumerate}
\end{algorithm}

\begin{theorem}\label{thm:dc-composition}
Algorithm~\ref{alg:dc-composition} returns $f\circ g$.
At pairing level $j$, $p=g^{2^j}$, and a complete live block beginning at
coefficient $i2^j$ represents
\[
 \sum_{r=0}^{2^j-1}a_{i2^j+r}g^r.
\]
Consequently the algorithm performs $n-1$ block multiplications and
$O(\log n)$ squarings of powers of $g$.
\end{theorem}

\begin{proof}
The proof of Lemma~\ref{lem:dc-evaluation-invariant} uses only ring operations,
so it remains valid with $c$ replaced by $g$.  The final block is therefore
$f(g)$.  The operation count follows as in
Theorem~\ref{thm:dc-evaluation}.
\end{proof}

Let $m=\len(g)$ and let $\mathsf M(r)$ denote the coefficient-operation cost
of multiplying polynomials of length $O(r)$.  At level $j$, both the current
power and a complete block have length $O(2^jm)$, while there are
$O(n/2^j)$ block products.  Hence
\begin{equation}\label{eq:dc-composition-cost}
 C(n,m)=O\left(
   \sum_{j<\lceil\log_2 n\rceil}
   \left(\frac{n}{2^j}+1\right)\mathsf M(2^jm)
 \right).
\end{equation}

\begin{corollary}\label{cor:dc-composition-cost}
If $\mathsf M(r)=\Theta(r^\alpha)$ with $\alpha>1$, then
\[
 C(n,m)=O(\mathsf M(nm)).
\]
If $\mathsf M(r)=O(r\log^\kappa r)$ for fixed $\kappa\ge0$, then
\[
 C(n,m)=O(\mathsf M(nm)\log n).
\]
\end{corollary}

\begin{proof}
For a power-law multiplication cost, the level-$j$ contribution in
\eqref{eq:dc-composition-cost} is
\[
 O\bigl(nm^\alpha2^{j(\alpha-1)}+m^\alpha2^{j\alpha}\bigr),
\]
so the geometrically largest final levels dominate and give
$O((nm)^\alpha)$.  In the quasi-linear case each nonempty level is at most a
constant multiple of $\mathsf M(nm)$, and there are $O(\log n)$ levels.
\end{proof}

For bit complexity it is useful to bound all coefficients at once.  Suppose
$f$ has coefficient bit size at most $\tau_f$ and $g$ has length $m$ and
coefficient bit size at most $\tau_g$.  The $1$-norm satisfies
\[
 \|g\|_1<m2^{\tau_g},
\]
and therefore
\begin{equation}\label{eq:composition-height}
 \|f(g)\|_\infty
 \le \|f(g)\|_1
 \le n2^{\tau_f}\max(1,m2^{\tau_g})^{n-1}.
\end{equation}
The same estimate, with a smaller exponent, controls every block and power
formed by Algorithms~\ref{alg:horner-composition} and
\ref{alg:dc-composition}.  Thus all coefficients have
\begin{equation}\label{eq:composition-bits}
 W=O\bigl(\tau_f+\log(n+1)+n(\tau_g+\log(m+1))\bigr)
\end{equation}
bits.  A direct bit bound is obtained from the coefficient-operation bounds
above by charging coefficient multiplications $O(\MZ(W))$ each and additions
$O(W)$ each.  This is conservative but makes the coefficient explosion from
composition explicit.

\subsection{Taylor shift}

For $c\in\ZZ$, the Taylor shift $f(x+c)$ is composition by the linear
polynomial $x+c$.  Its coefficients satisfy the binomial formula
\begin{equation}\label{eq:taylor-binomial}
 [x^j]f(x+c)=\sum_{i=j}^{n-1}a_i\binom{i}{j}c^{i-j}.
\end{equation}
The companion implementation provides both a specialized Horner shift and the
divide-and-conquer composition of Algorithm~\ref{alg:dc-composition} with
$g=x+c$.

\begin{algorithm}[In-place Horner Taylor shift]\label{alg:horner-taylor}
Input $f=\sum_{i=0}^{n-1}a_ix^i$ and $c\in\ZZ$.
\begin{enumerate}
 \item If $n=0$ or $c=0$, return $f$.
 \item Put $r\leftarrow a_{n-1}$.
 \item For $i=n-2,\ldots,0$, replace
 \[
    r(x)\leftarrow (x+c)r(x)+a_i,
 \]
 updating the coefficients of $r$ from high degree to low degree so that the
 update uses only one coefficient array.
 \item Return $r$.
\end{enumerate}
\end{algorithm}

\begin{theorem}\label{thm:horner-taylor}
Algorithm~\ref{alg:horner-taylor} returns $f(x+c)$.  It uses exactly
\[
 \frac{n(n-1)}2
\]
multiplications by $c$ and the same number of coefficient additions, apart
from coefficient copies used to create the new leading term.
\end{theorem}

\begin{proof}
After processing $a_i$, Horner's invariant is
\[
 r(x)=a_i+(x+c)\bigl(a_{i+1}+(x+c)(\cdots+(x+c)a_{n-1})\bigr),
\]
so the final result is $f(x+c)$.  When the current polynomial has degree $d$,
multiplication by $x+c$ requires $d+1$ scalar multiplications by $c$ and
$d+1$ additions.  The successive values are $d=0,1,\ldots,n-2$, whose sum
of $d+1$ is $n(n-1)/2$.
\end{proof}

If $\bits(a_i)\le\tau$ and $\xi=\max(1,\bits(c))$, then
\eqref{eq:taylor-binomial}, or the $1$-norm bound
$\|f(x+c)\|_1\le\|f\|_1(1+|c|)^{n-1}$, gives
\begin{equation}\label{eq:taylor-growth}
 \bits([x^j]f(x+c))=O\bigl(\tau+\log(n+1)+n(\xi+1)\bigr).
\end{equation}
The same bound controls the Horner intermediates.  Consequently a direct bit
bound for the specialized algorithm is
\[
 O\left(n^2\MZ\bigl(\tau+n(\xi+1)+\log(n+1)\bigr)\right),
\]
with the usual improvement when multiplication by the fixed integer $c$ is
cheaper than a general integer product.

\begin{theorem}\label{thm:dc-taylor}
Computing the Taylor shift as divide-and-conquer composition with $g=x+c$ has
coefficient-operation cost
\[
 O\left(
  \sum_{j<\lceil\log_2n\rceil}
  \left(\frac{n}{2^j}+1\right)\mathsf M(2^j)
 \right).
\]
Hence the cost is $O(\mathsf M(n))$ when
$\mathsf M(n)=\Theta(n^\alpha)$ with $\alpha>1$, and
$O(\mathsf M(n)\log n)$ for quasi-linear multiplication.  Its coefficient
bit size is bounded by \eqref{eq:taylor-growth}.
\end{theorem}

\begin{proof}
This is Corollary~\ref{cor:dc-composition-cost} specialized to a polynomial
$g$ of length two.  The coefficient-growth statement follows from
\eqref{eq:composition-height} with $\|x+c\|_1=1+|c|$.
\end{proof}

\begin{remark}
The specialized Horner shift has quadratic coefficient complexity but very
small setup cost and uses only one coefficient array.  The balanced method
reduces the asymptotic polynomial-arithmetic cost, but it constructs and
multiplies powers of $x+c$.  A practical dispatcher should therefore retain
both algorithms rather than choose solely from their asymptotic bounds.
\end{remark}



\section{Formal differentiation}\label{sec:formal-differentiation}

For
\[
 f(x)=\sum_{i=0}^{d}a_i x^i\in\ZZ[x],
\]
the \emph{formal derivative} is
\[
 f'(x)=\sum_{i=1}^{d} i a_i x^{i-1}.
\]
No analytic notion of a limit is involved: differentiation is an algebraic
operation on coefficients.  Over $\ZZ$ it is a derivation, so it is additive
and satisfies $(fg)'=f'g+fg'$.  The first derivative is coefficientwise and
therefore elementary; the only algorithmic point of substance is to compute a
higher derivative without constructing all lower derivatives first.

For $k\ge0$, write
\[
 (i)_k=i(i-1)\cdots(i-k+1),\qquad (i)_0=1
\]
for the falling factorial.  Repeated application of the definition gives
\begin{equation}\label{eq:nth-derivative-formula}
 f^{(k)}(x)=\sum_{i=k}^{d}(i)_k a_i x^{i-k}
 =\sum_{i=k}^{d}k!\binom{i}{k}a_i x^{i-k}.
\end{equation}
Thus $f^{(k)}=0$ when $k>d$, while for $0\le k\le d$,
\[
 \deg f^{(k)}=d-k.
\]

\subsection{First derivative}

\begin{algorithm}[Formal derivative]\label{alg:first-derivative}
Input $f=\sum_{i=0}^{n-1}a_i x^i\in\ZZ[x]$.
\begin{enumerate}
 \item If $n\le1$, return $0$.
 \item For $i=1,\ldots,n-1$, set
 \[
    b_{i-1}\leftarrow i a_i.
 \]
 \item Return $b=\sum_{i=0}^{n-2}b_i x^i$.
\end{enumerate}
\end{algorithm}

\begin{theorem}\label{thm:first-derivative}
Algorithm~\ref{alg:first-derivative} returns $f'$.  For $n>1$ it performs
exactly $n-1$ multiplications of an integer coefficient by a nonnegative
machine integer.  If the input coefficients have at most $\tau$ bits, every
output coefficient has
\[
 O\bigl(\tau+\log(n+1)\bigr)
\]
bits.
\end{theorem}

\begin{proof}
The coefficient of $x^{i-1}$ in the formal derivative is $ia_i$, which is
exactly the value written in step 2.  Since $1\le i<n$,
\[
 \bits(ia_i)\le \bits(a_i)+\lceil\log_2 n\rceil+O(1),
\]
proving the size bound.
\end{proof}

Consequently, if multiplication by an $O(\log n)$-bit nonnegative integer is
charged by the general integer multiplication bound, a safe bit-complexity
estimate is
\begin{equation}\label{eq:first-derivative-bit}
 O\left(n\MZ\bigl(\tau+\log(n+1)\bigr)\right).
\end{equation}
For small indices this deliberately overestimates the cost of scalar
multiplication.

\subsection{Direct higher derivatives}

Formula~\eqref{eq:nth-derivative-formula} already avoids the construction of
$k-1$ intermediate derivative polynomials.  Computing a binomial coefficient
from scratch for every $i$ is unnecessary, however.  The falling factorials
satisfy the exact recurrence
\begin{equation}\label{eq:falling-factorial-recurrence}
 (k)_k=k!,\qquad
 (i+1)_k=(i)_k\frac{i+1}{i+1-k}\quad(i\ge k).
\end{equation}
The quotient is an integer because both sides are falling factorials.

\begin{algorithm}[Direct $k$-th derivative]\label{alg:nth-derivative}
Input $f=\sum_{i=0}^{n-1}a_i x^i\in\ZZ[x]$ and $k\ge0$.
\begin{enumerate}
 \item If $k=0$, return $f$.  If $k\ge n$, return $0$.
 \item Set $u\leftarrow k!$.
 \item For $i=k,\ldots,n-1$:
 \begin{enumerate}
  \item set $b_{i-k}\leftarrow u a_i$;
  \item if $i+1<n$, replace
  \[
      u\leftarrow u(i+1)/(i+1-k),
  \]
  where the division is exact.
 \end{enumerate}
 \item Return $b=\sum_{j=0}^{n-k-1}b_jx^j$.
\end{enumerate}
\end{algorithm}

\begin{theorem}\label{thm:nth-derivative-correct}
Algorithm~\ref{alg:nth-derivative} returns $f^{(k)}$.  During the iteration
with index $i$, the invariant is
\[
 u=(i)_k.
\]
Hence every division in step 3(b) is exact.
\end{theorem}

\begin{proof}
At the first iteration, $u=k!=(k)_k$.  If $u=(i)_k$, then
\eqref{eq:falling-factorial-recurrence} changes it to $(i+1)_k$, proving the
invariant by induction.  Step 3(a) therefore writes
$(i)_k a_i$ into the coefficient of $x^{i-k}$.  Equation
\eqref{eq:nth-derivative-formula} proves that the returned polynomial is
$f^{(k)}$.
\end{proof}

Let $m=n-k>0$.  Apart from the one factorial computation, the algorithm uses
exactly
\[
 m
\]
general coefficient multiplications, together with $m-1$ multiplications and
$m-1$ exact divisions by machine integers to advance the falling factorial.
A naive computation by applying the first-derivative algorithm $k$ times
instead visits
\begin{equation}\label{eq:repeated-derivative-count}
 \sum_{j=0}^{k-1}(n-j-1)
 =\frac{k(2n-k-1)}{2}
\end{equation}
coefficients.  Thus direct higher differentiation reduces the number of
coefficient positions processed from $\Theta(kn)$ in the balanced range to
$O(n)$, and allocates only the final polynomial rather than $k$ successive
intermediate polynomials.  For very small fixed $k$, repeated differentiation
can nevertheless be competitive because it multiplies only by small
integers.

\begin{lemma}[Coefficient growth]\label{lem:nth-derivative-growth}
Suppose $\deg f=d$, $0\le k\le d$, and $H(f)=\max_i|a_i|$.  Then
\[
 H\bigl(f^{(k)}\bigr)\le (d)_k H(f).
\]
If the coefficients of $f$ have at most $\tau$ bits, every coefficient of
$f^{(k)}$ therefore has
\begin{equation}\label{eq:nth-derivative-bits}
 O\bigl(\tau+k\log(d+1)\bigr)
\end{equation}
bits.
\end{lemma}

\begin{proof}
By \eqref{eq:nth-derivative-formula}, the coefficient descending from $a_i$
is $(i)_k a_i$.  For $k\le i\le d$ one has $(i)_k\le(d)_k\le d^k$.
Taking absolute values and then binary logarithms proves both assertions.
\end{proof}

\begin{theorem}[Bit complexity]\label{thm:nth-derivative-bit}
Let $m=n-k>0$ and put
\[
 W=O\bigl(\tau+k\log(n+1)\bigr).
\]
If $\mathsf F(k)$ denotes the bit cost of computing $k!$, then
Algorithm~\ref{alg:nth-derivative} has the safe bound
\[
 O\bigl(\mathsf F(k)+m\MZ(W)\bigr).
\]
Using the elementary sequential factorial computation gives, for example,
\[
 \mathsf F(k)=O\bigl(k\MZ(k\log(k+1))\bigr).
\]
The actual factorial routine may use a faster product tree or specialised
factorial algorithm.
\end{theorem}

\begin{proof}
Lemma~\ref{lem:nth-derivative-growth} bounds both the final coefficient sizes
and the falling factorials by $O(W)$ bits.  Each of the $m$ coefficient
products and each of the $O(m)$ updates of $u$ is therefore bounded safely by
$O(\MZ(W))$ bit operations; division by a machine integer is no more expensive
under this coarse bound.  Adding the factorial setup cost proves the first
claim.  Multiplying successively by $1,2,\ldots,k$ and charging every product
at the final factorial size $O(k\log(k+1))$ gives the second.
\end{proof}

The higher-derivative routine is therefore not merely a convenience wrapper
around first differentiation: it is a separate direct algorithm whose work is
proportional to the number of input and output coefficient positions, while
the unavoidable growth of the falling-factorial multipliers is accounted for
explicitly in the bit bound.

\section{Low and high truncated products}

Suppose again that
\[
 f(x)=\sum_{i=0}^{m-1}a_ix^i,
 \qquad
 g(x)=\sum_{j=0}^{n-1}b_jx^j.
\]
For a positive integer $r$, define the \emph{low product}
\[
 \operatorname{low}_r(fg)=(fg)\bmod x^r.
\]
Thus only coefficients of degrees $0,\ldots,r-1$ are retained.  When
$r\geq m+n-1$ this is simply the full product.  The terminology \emph{short
product} is also common; systematic recursive short-product algorithms were
studied by Mulders \cite{Mulders2000}.

If
\[
 h(x)=\sum_{i=0}^{L-1}c_ix^i
\]
has length $L$ and $1\leq r\leq L$, define its \emph{high $r$-term part},
shifted down to degree zero, by
\[
 \operatorname{high}_r(h)=\sum_{i=0}^{r-1}c_{L-r+i}x^i.
\]
The low product is the more natural primitive.  High products will be reduced
to low products by reversal.

Two elementary observations are useful throughout.  First,
\begin{equation}\label{eq:low-truncate-inputs}
 \operatorname{low}_r(fg)
 =\operatorname{low}_r\bigl(\operatorname{low}_r(f)
                            \operatorname{low}_r(g)\bigr),
\end{equation}
so coefficients of either input of degree at least $r$ may be discarded.
Second, every retained coefficient is a sum of at most $r$ products of input
coefficients.  Hence, if the input coefficients have at most $\tau$ bits,
then every coefficient occurring in a low product has
$O(\tau+\log r)$ bits, by the same argument as
\eqref{eq:product-coefficient-bound}.

\subsection{Classical low product}

The classical algorithm simply visits the portion of the convolution triangle
which lies below degree $r$.

\begin{algorithm}[Classical low product]\label{alg:classical-low-product}
Input $f,g\in\ZZ[x]$ and a positive integer $r$.
\begin{enumerate}
 \item Replace each input by its first $r$ coefficients if necessary.
 \item Set $c_0,\ldots,c_{r-1}\leftarrow0$.
 \item For every pair $(i,j)$ with
 \[
    0\leq i<m,\qquad 0\leq j<n,\qquad i+j<r,
 \]
 set
 \[
    c_{i+j}\leftarrow c_{i+j}+a_ib_j.
 \]
 \item Return $\sum_{k=0}^{r-1}c_kx^k$, omitting trailing zero
 coefficients in the usual polynomial representation.
\end{enumerate}
\end{algorithm}

\begin{theorem}\label{thm:classical-low-product}
Algorithm~\ref{alg:classical-low-product} returns
$\operatorname{low}_r(fg)$.  It uses exactly
\begin{equation}\label{eq:classical-low-count}
 P(m,n,r)=
 \sum_{i=0}^{\min(m,r)-1}\min(n,r-i)
\end{equation}
coefficient multiplications and the same number of coefficient
accumulations.  In particular, if $m,n\geq r$, then
\[
 P(m,n,r)=\frac{r(r+1)}2.
\]
If all input coefficients have at most $\tau$ bits, the bit complexity is
\[
 O\!\left(
 P(m,n,r)\bigl(\MZ(\tau)+\tau+\log r\bigr)
 \right).
\]
For $m,n\geq r$ this is
\[
 O\!\left(r^2\bigl(\MZ(\tau)+\tau+\log r\bigr)\right).
\]
\end{theorem}

\begin{proof}
By the convolution identity \eqref{eq:convolution}, a pair $(i,j)$ contributes
precisely to the coefficient of degree $i+j$.  A pair is therefore relevant
exactly when $i+j<r$, which proves correctness.  For fixed $i$, the admissible
indices $j$ are
\[
 0\leq j<\min(n,r-i),
\]
and summing their number gives \eqref{eq:classical-low-count}.  If both
inputs have at least $r$ coefficients, this becomes
\[
 r+(r-1)+\cdots+1=\frac{r(r+1)}2.
\]

Each coefficient multiplication costs $O(\MZ(\tau))$.  A retained output
coefficient is a sum of at most $r$ such products and consequently has
$O(\tau+\log r)$ bits, so each accumulation costs
$O(\tau+\log r)$.  Multiplying these costs by the number of contributing
pairs proves the bit bound.
\end{proof}

For balanced length-$r$ operands, the leading term in the classical operation
count is therefore one half of the leading term for a full classical product.
This is the geometric reason that a truncated algorithm is worthwhile even
before any subquadratic multiplication method is used.

\subsection{Recursive low product}

A low product can be reduced to one full product and two smaller low products.
After applying \eqref{eq:low-truncate-inputs}, pad the operands with leading
zeros if necessary so that both have formal length $r$.  Put
\[
 k=\left\lceil\frac r2\right\rceil,
 \qquad
 h=r-k=\left\lfloor\frac r2\right\rfloor,
\]
and split
\[
 f=f_0+x^k f_1,
 \qquad
 g=g_0+x^k g_1,
\]
where $f_0,g_0$ have length at most $k$ and $f_1,g_1$ have length at most
$h$.  Then
\begin{align}
 fg={}&f_0g_0+x^k(f_0g_1+f_1g_0)+x^{2k}f_1g_1.
\end{align}
Since $2k\geq r$, the final term cannot contribute below degree $r$.
Furthermore, after multiplication by $x^k$, only the lowest $h=r-k$
coefficients of each cross product can contribute.

\begin{algorithm}[Recursive low product]\label{alg:recursive-low-product}
Input $f,g\in\ZZ[x]$ and a positive integer $r$.
\begin{enumerate}
 \item If $r$ is below a chosen base-case threshold, use
 Algorithm~\ref{alg:classical-low-product}.
 \item Truncate both inputs to length at most $r$.  Put
 \[
 k\leftarrow\lceil r/2\rceil,
 \qquad h\leftarrow r-k,
 \]
 and write
 \[
 f=f_0+x^kf_1,
 \qquad g=g_0+x^kg_1.
 \]
 \item Compute the full product
 \[
 z_0\leftarrow f_0g_0.
 \]
 \item Recursively compute
 \[
 z_1\leftarrow\operatorname{low}_h(f_0g_1),
 \qquad
 z_2\leftarrow\operatorname{low}_h(f_1g_0).
 \]
 \item Return
 \[
 \operatorname{low}_r\bigl(z_0+x^k(z_1+z_2)\bigr).
 \]
\end{enumerate}
\end{algorithm}

\begin{theorem}\label{thm:recursive-low-correctness}
Algorithm~\ref{alg:recursive-low-product} returns
$\operatorname{low}_r(fg)$.
\end{theorem}

\begin{proof}
With the decomposition above,
\[
 fg=f_0g_0+x^k(f_0g_1+f_1g_0)+x^{2k}f_1g_1.
\]
Because $2k\geq r$,
\[
 \operatorname{low}_r(x^{2k}f_1g_1)=0.
\]
For any polynomial $u$,
\[
 \operatorname{low}_r(x^ku)
 =x^k\operatorname{low}_{r-k}(u)
 =x^k\operatorname{low}_h(u).
\]
Thus
\[
 \operatorname{low}_r(fg)
 =\operatorname{low}_r\!\left(
 f_0g_0+x^k\bigl(
 \operatorname{low}_h(f_0g_1)
 +\operatorname{low}_h(f_1g_0)
 \bigr)\right),
\]
which is exactly the quantity returned by the algorithm.  Induction on $r$
completes the proof.
\end{proof}

For the complexity analysis, let $\mathsf M(r)$ denote the number of
coefficient operations needed by the chosen full-product algorithm on operands
of length at most $r$, and let $L(r)$ denote the corresponding cost of
Algorithm~\ref{alg:recursive-low-product}.  Apart from the recursive products,
the algorithm performs $O(r)$ coefficient copies and additions.  Therefore
\begin{equation}\label{eq:low-product-recurrence}
 L(r)\leq
 \mathsf M\!\left(\left\lceil\frac r2\right\rceil\right)
 +2L\!\left(\left\lfloor\frac r2\right\rfloor\right)
 +O(r).
\end{equation}

\begin{theorem}\label{thm:recursive-low-complexity}
Suppose
\[
 \mathsf M(r)=\Theta(r^\alpha)
\]
for some fixed $\alpha>1$.  Then
\[
 L(r)=\Theta(r^\alpha)=\Theta(\mathsf M(r)).
\]
More precisely, if
\[
 \mathsf M(r)\sim A r^\alpha
\]
and the linear recombination cost is lower order, then the balanced recurrence
\eqref{eq:low-product-recurrence} has leading term
\begin{equation}\label{eq:balanced-low-leading-constant}
 L(r)\sim \frac{A}{2^\alpha-2}\,r^\alpha.
\end{equation}
\end{theorem}

\begin{proof}
Ignoring ceilings and floors changes only lower-order terms.  Substitution into
\eqref{eq:low-product-recurrence} gives the standard recurrence
\[
 L(r)=2L(r/2)+\Theta(r^\alpha).
\]
Since $\alpha>1$, the work per level decreases geometrically by the factor
$2^{1-\alpha}<1$, so the recursion tree sums to $\Theta(r^\alpha)$.

For the sharper statement, first take $r=2^t$.  The contribution of the
full products at recursion depth $j$ is asymptotic to
\[
 2^j A\left(\frac{r}{2^{j+1}}\right)^\alpha
 =\frac{Ar^\alpha}{2^\alpha}
   \left(2^{1-\alpha}\right)^j.
\]
The leaf and linear recombination costs are $o(r^\alpha)$, while the displayed
terms form a geometric series.  Hence
\[
 L(r)\sim
 \frac{Ar^\alpha}{2^\alpha}
 \sum_{j\geq0}\left(2^{1-\alpha}\right)^j
 =\frac{A}{2^\alpha-2}r^\alpha.
\]
For general $r$, replacing the subproblem sizes by the adjacent powers of two,
or equivalently retaining the ceilings and floors in the same recursion-tree
sum, changes the result by only lower-order terms.  This proves
\eqref{eq:balanced-low-leading-constant}.
\end{proof}

For example, the leading constant in
\eqref{eq:balanced-low-leading-constant} is $1/2$ for quadratic full
multiplication and $1$ for Karatsuba multiplication.  The equal split is a
particularly simple short-product decomposition; unequal block sizes can give
better constants for subquadratic full multiplication.  Mulders introduced
this refinement, and Hanrot and Zimmermann determined the optimal cutoff in
the Karatsuba model \cite{Mulders2000,HanrotZimmermann2004}.

A corresponding bit-complexity statement can be made without committing to a
particular full-product backend.  Let $\mathsf B(r,\tau)$ be an upper bound for
multiplying two length-$r$ polynomials whose coefficients have at most $\tau$
bits, and let $L_{\rm bit}(r,\tau)$ denote the cost of the recursive low
product.  No evaluation or coefficient transformation is performed in the
recursion, so recursive inputs still have at most $\tau$ bits.  Recombination
uses integers of $O(\tau+\log r)$ bits.  Consequently
\begin{equation}\label{eq:low-bit-recurrence}
 L_{\rm bit}(r,\tau)
 \leq
 \mathsf B(\lceil r/2\rceil,\tau)
 +2L_{\rm bit}(\lfloor r/2\rfloor,\tau)
 +O\bigl(r(\tau+\log r)\bigr).
\end{equation}
In particular, whenever for some fixed $\alpha>1$ one has
\[
 \mathsf B(r,\tau)
 =O\!\left(r^\alpha\MZ(\tau+\log r)\right),
\]
the same recursion-tree argument gives
\[
 L_{\rm bit}(r,\tau)
 =O\!\left(r^\alpha\MZ(\tau+\log r)\right).
\]
Thus the recursive truncation preserves the asymptotic bit complexity of the
Karatsuba- and Toom-type full-product bounds proved in the preceding section.

\subsection{High product by reversal}

For a polynomial of formal length $d$, define the fixed-length reversal
\[
 \operatorname{rev}_d\!\left(\sum_{i=0}^{d-1}u_ix^i\right)
 =\sum_{i=0}^{d-1}u_i x^{d-1-i}.
\]
Leading zero coefficients are retained for the purpose of this operation.
If $f$ and $g$ have lengths $m$ and $n$, respectively, then
\begin{equation}\label{eq:reversal-product-identity}
 \operatorname{rev}_{m+n-1}(fg)
 =\operatorname{rev}_m(f)\operatorname{rev}_n(g).
\end{equation}
The highest coefficients of $fg$ therefore become the lowest coefficients of
the product of the reversed operands.

\begin{algorithm}[High product by reversal]\label{alg:high-by-reversal}
Input nonzero $f,g\in\ZZ[x]$ of lengths $m,n$ and
$1\leq r\leq m+n-1$.
\begin{enumerate}
 \item Put $r_f\leftarrow\min(r,m)$ and $r_g\leftarrow\min(r,n)$.
 Extract the highest $r_f$ coefficients of $f$ and the highest $r_g$
 coefficients of $g$, shifting each tail down to degree zero; call the
 resulting polynomials $f_{\rm tail}$ and $g_{\rm tail}$.
 \item Form only the needed reversed tails,
 \[
    f^*\leftarrow\operatorname{rev}_{r_f}(f_{\rm tail}),
    \qquad
    g^*\leftarrow\operatorname{rev}_{r_g}(g_{\rm tail}).
 \]
 \item Compute
 \[
    q\leftarrow\operatorname{low}_r(f^*g^*).
 \]
 Regard $q$ as having formal length $r$, padding by leading zeros if needed.
 \item Return $\operatorname{rev}_r(q)$.
\end{enumerate}
\end{algorithm}

\begin{theorem}\label{thm:high-by-reversal}
Algorithm~\ref{alg:high-by-reversal} returns
$\operatorname{high}_r(fg)$.  If a low $r$-term product costs $L(r)$
coefficient operations, then a high $r$-term product costs
\[
 L(r)+O(r)
\]
coefficient operations when only the relevant input tails are materialized.
The same reduction adds only $O(r\tau)$ bit operations for coefficients of
$O(\tau)$ bits, apart from the cost of the low product itself.
\end{theorem}

\begin{proof}
Write
\[
 fg=\sum_{i=0}^{m+n-2}c_ix^i.
\]
By \eqref{eq:reversal-product-identity},
\[
 \operatorname{rev}_m(f)\operatorname{rev}_n(g)
 =c_{m+n-2}+c_{m+n-3}x+\cdots+c_0x^{m+n-2}.
\]
Its lowest $r$ coefficients are therefore
\[
 c_{m+n-2}+c_{m+n-3}x+\cdots+c_{m+n-r-1}x^{r-1}.
\]
Reversal of formal length $r$ gives
\[
 c_{m+n-r-1}+c_{m+n-r}x+\cdots+c_{m+n-2}x^{r-1},
\]
which is exactly $\operatorname{high}_r(fg)$.

Only the lowest $r$ coefficients of each full reversed operand can influence
its low product.  These are exactly the reversals of the highest
$\min(r,m)$ and $\min(r,n)$ coefficients of the original operands.  Thus the
full reversals need never be materialized.  Forming the two reversed tails and
reversing the output require $O(r)$ coefficient copies.  For
$O(\tau)$-bit coefficients this is $O(r\tau)$ bit movement.  All remaining
work is the low product.
\end{proof}

The reduction is useful beyond multiplication itself.  Any algorithm expressed
in terms of a low product immediately yields a symmetric high-product variant
without a second multiplication theory; this is why low products are usually
taken as the primary truncated operation.


\section{Middle products}

The middle product is the truncated operation naturally dual to ordinary
polynomial multiplication.  It is particularly useful when a recursive
algorithm needs coefficients from an interior window rather than from either
end of a product.  Hanrot, Quercia and Zimmermann introduced the operation in
this form and emphasized its algorithmic role \cite{HanrotQuerciaZimmermann2004};
Bostan, Lecerf and Schost give a systematic treatment of the transposition
principle behind it \cite{BostanLecerfSchost2003}.  Harvey gives the
Karatsuba middle-product formulas in essentially the form used below
\cite{Harvey2012Middle}.

For $n\geq1$, let
\[
 a(x)=\sum_{i=0}^{2n-2}a_i x^i,
 \qquad
 b(x)=\sum_{j=0}^{n-1}b_j x^j,
\]
where missing coefficients are interpreted as zero.  The \emph{balanced
middle product} is
\begin{equation}\label{eq:middle-definition}
 \operatorname{MP}_n(a,b)
 =\sum_{i=0}^{n-1}
   \left(\sum_{j=0}^{n-1}a_{n-1+i-j}b_j\right)x^i.
\end{equation}
Thus $\operatorname{MP}_n(a,b)$ consists of coefficients $n-1$ through
$2n-2$ of $ab$, shifted down to degrees $0$ through $n-1$.  More generally,
one may ask for any consecutive window of a product; the balanced operation
\eqref{eq:middle-definition} is the canonical interior case from which the
fast algorithms below are built.

If $\bits(a_i),\bits(b_j)\leq\tau$, every coefficient in
\eqref{eq:middle-definition} is a sum of exactly $n$ products after formal
zero-padding, and hence
\begin{equation}\label{eq:middle-coeff-bound}
 \bits\bigl([x^i]\operatorname{MP}_n(a,b)\bigr)
 \leq 2\tau+\lceil\log_2 n\rceil+1.
\end{equation}

\subsection{Classical middle product}

Equation~\eqref{eq:middle-definition} gives the base case directly.

\begin{algorithm}[Classical middle product]\label{alg:middle-classical}
Input $a$ of formal length $2n-1$ and $b$ of formal length $n$.
\begin{enumerate}
 \item Set $c_0,\ldots,c_{n-1}\leftarrow0$.
 \item For $i=0,\ldots,n-1$ and $j=0,\ldots,n-1$, set
 \[
    c_i\leftarrow c_i+a_{n-1+i-j}b_j.
 \]
 \item Return $\sum_{i=0}^{n-1}c_ix^i$.
\end{enumerate}
\end{algorithm}

\begin{theorem}\label{thm:middle-classical}
Algorithm~\ref{alg:middle-classical} returns
$\operatorname{MP}_n(a,b)$.  It uses exactly $n^2$ coefficient
multiplications and $n^2$ coefficient accumulations.  For coefficients of at
most $\tau$ bits, its bit complexity is
\[
 O\!\left(n^2\bigl(\MZ(\tau)+\tau+\log n\bigr)\right).
\]
\end{theorem}

\begin{proof}
The formula accumulated into $c_i$ is precisely the inner sum in
\eqref{eq:middle-definition}, proving correctness.  There are $n^2$ pairs
$(i,j)$.  Each coefficient product costs $O(\MZ(\tau))$, while by
\eqref{eq:middle-coeff-bound} each accumulation involves
$O(\tau+\log n)$ bits.  Summing over all pairs gives the stated bound.
\end{proof}

The classical middle product should be compared with the complete product of a
$(2n-1)$-term polynomial by an $n$-term polynomial, which uses
$n(2n-1)$ classical coefficient products.  The middle window therefore removes
about half of the quadratic work, while retaining all $n^2$ products that can
contribute to the requested coefficients.

\subsection{Transposition}

The structural reason fast middle products exist is that the middle product is
the transpose of polynomial multiplication, after reversing the short
operand.  Write
\[
 b^*=\operatorname{rev}_n(b)
      =\sum_{j=0}^{n-1}b_{n-1-j}x^j.
\]
For fixed $b^*$, multiplication defines a linear map
\[
 \mu_{b^*}:\ZZ^n\longrightarrow\ZZ^{2n-1},
 \qquad u\longmapsto u b^*.
\]
We use the usual coefficient-vector scalar product.

\begin{theorem}[Middle product as a transpose]\label{thm:middle-transpose}
For fixed $b$, the transpose of $\mu_{b^*}$ is the map
\[
 a\longmapsto\operatorname{MP}_n(a,b).
\]
Consequently, a bilinear multiplication algorithm of rank $R$ can be
transposed to a middle-product algorithm using the same $R$ recursive
multiplications, with only linear pre- and postprocessing changed.
\end{theorem}

\begin{proof}
Let $M$ be the $(2n-1)\times n$ matrix of multiplication by $b^*$.
Its entry in row $k$ and column $j$ is
$b^*_{k-j}$ when $0\leq k-j<n$, and zero otherwise.  Hence the $i$th
coordinate of $M^{\mathsf T}a$ is
\[
 \sum_{\ell=0}^{n-1}a_{i+\ell}b^*_{\ell}
 =\sum_{\ell=0}^{n-1}a_{i+\ell}b_{n-1-\ell}.
\]
Putting $j=n-1-\ell$ gives
\[
 \sum_{j=0}^{n-1}a_{n-1+i-j}b_j,
\]
which is exactly coefficient $i$ of \eqref{eq:middle-definition}.

For the complexity statement, write a rank-$R$ bilinear multiplication
algorithm abstractly as
\[
 \mu(u,v)=\sum_{r=1}^{R} w_r\,L_r(u)K_r(v),
\]
where $L_r,K_r$ are linear forms and the $w_r$ are fixed output vectors.
For fixed $v$, transposition replaces the linear map
$u\mapsto\mu(u,v)$ by
\[
 h\longmapsto
 \sum_{r=1}^{R}L_r^{\mathsf T}
 \bigl(K_r(v)\langle w_r,h\rangle\bigr).
\]
There are still exactly $R$ multiplicative nodes; only the surrounding linear
maps have been transposed.  This is the bilinear form of the transposition
principle.  A general straight-line formulation, including precise linear
operation counts, is given in \cite{BostanLecerfSchost2003}.
\end{proof}

Thus middle products can have the same multiplication exponent as ordinary
products.  In particular, transposing Karatsuba gives three half-size middle
products, while transposing Toom-3 gives five one-third-size middle products.
The explicit algorithms below are integral schedules for these transposed
bilinear algorithms.

\subsection{Block form of the transposition}

A block identity makes the Toom-4/2 and Toom-6/3 formulas transparent.  Let
$n=qk$.  Split the short operand into $q$ blocks
\[
 G_j=\sum_{\ell=0}^{k-1}b_{jk+\ell}x^\ell,
 \qquad 0\leq j<q,
\]
and form the $2q-1$ overlapping windows of the long operand
\[
 F_t=\sum_{u=0}^{2k-2}a_{tk+u}x^u,
 \qquad 0\leq t\leq2q-2.
\]
Write the output as
\[
 \operatorname{MP}_{qk}(a,b)=C_0+x^kC_1+\cdots+x^{(q-1)k}C_{q-1},
 \qquad \len(C_s)\leq k.
\]

\begin{lemma}[Block middle-product identity]\label{lem:middle-block}
For $0\leq s<q$,
\begin{equation}\label{eq:middle-block}
 C_s=\sum_{j=0}^{q-1}
 \operatorname{MP}_k(F_{q-1+s-j},G_j).
\end{equation}
\end{lemma}

\begin{proof}
Let $0\leq h<k$.  The coefficient of $x^{sk+h}$ in the shifted middle
product is the coefficient of degree
$qk-1+sk+h$ in $ab$.  A term from block $G_j$ has short-operand index
$jk+\ell$, so the corresponding long-operand index is
\[
 qk-1+sk+h-jk-\ell
 =(q-1+s-j)k+(k-1+h-\ell).
\]
This is exactly the coefficient selected by the $h$th coefficient of
$\operatorname{MP}_k(F_{q-1+s-j},G_j)$.  Summing over $j$ proves
\eqref{eq:middle-block} coefficientwise.
\end{proof}

For $q=2$, this is the transpose of the outer two-block convolution of
Karatsuba multiplication.  For $q=3$, it is the transpose of the outer
three-block convolution used by Toom-3.

\subsection{Toom-4/2: transposed Karatsuba}

Let $n=2k$.  Lemma~\ref{lem:middle-block} gives
\begin{align}
 C_0&=\operatorname{MP}_k(F_0,G_1)
      +\operatorname{MP}_k(F_1,G_0),\label{eq:t42-c0}\\
 C_1&=\operatorname{MP}_k(F_1,G_1)
      +\operatorname{MP}_k(F_2,G_0).\label{eq:t42-c1}
\end{align}
Computing these four terms independently would not improve on the outer
quadratic decomposition.  Transposed Karatsuba reduces them to three recursive
middle products.

\begin{algorithm}[Toom-4/2 middle product]\label{alg:middle-toom42}
Input $a$ of formal length $4k-1$ and $b$ of formal length $2k$.
Form $F_0,F_1,F_2$ and $G_0,G_1$ as above, then compute
\begin{align*}
 P_0&\leftarrow\operatorname{MP}_k(F_0+F_1,G_1),\\
 P_1&\leftarrow\operatorname{MP}_k(F_1,G_0-G_1),\\
 P_2&\leftarrow\operatorname{MP}_k(F_1+F_2,G_0).
\end{align*}
Return
\[
 (P_0+P_1)+x^k(P_2-P_1).
\]
\end{algorithm}

\begin{theorem}\label{thm:middle-toom42}
Algorithm~\ref{alg:middle-toom42} returns
$\operatorname{MP}_{2k}(a,b)$.  If used recursively, its coefficient-operation
complexity satisfies
\[
 T_{42}(n)=3T_{42}(n/2)+O(n),
\]
and therefore
\[
 T_{42}(n)=\Theta\!\left(n^{\log_2 3}\right).
\]
For input coefficients of at most $\tau$ bits, a corresponding bit-complexity
bound is
\[
 O\!\left(n^{\log_2 3}\MZ(\tau+\log n)\right).
\]
\end{theorem}

\begin{proof}
By bilinearity,
\begin{align*}
 P_0+P_1
 &=\operatorname{MP}_k(F_0,G_1)
   +\operatorname{MP}_k(F_1,G_1)
   +\operatorname{MP}_k(F_1,G_0)
   -\operatorname{MP}_k(F_1,G_1)\\
 &=C_0
\end{align*}
by \eqref{eq:t42-c0}.  Similarly,
\begin{align*}
 P_2-P_1
 &=\operatorname{MP}_k(F_1,G_0)
   +\operatorname{MP}_k(F_2,G_0)
   -\operatorname{MP}_k(F_1,G_0)
   +\operatorname{MP}_k(F_1,G_1)\\
 &=C_1
\end{align*}
by \eqref{eq:t42-c1}.  This proves correctness.

There are three recursive middle products of size $n/2$.  Forming the three
linear combinations and recombining the two output blocks costs $O(n)$
coefficient additions and subtractions, giving the stated recurrence and
exponent.

Along a recursion path, each Toom-4/2 evaluation replaces a coefficient by a
sum or difference of at most two coefficients from the preceding level.
After $O(\log n)$ levels, coefficient sizes have therefore grown by only
$O(\log n)$ bits.  The recursion tree has
$\Theta(n^{\log_2 3})$ multiplicative leaves and the same order of linear
work.  Since $\MZ(t)=\Omega(t)$, the linear work is absorbed by
$O(n^{\log_2 3}\MZ(\tau+\log n))$.
\end{proof}

The name Toom-4/2 refers to the approximate operand lengths: the long operand
has $4k-1$ coefficients and the short operand $2k$, while the recursive middle
products have shape $(2k-1)\times k\to k$.

\subsection{Odd sizes in Toom-4/2}

A simple residue correction avoids padding when $n$ is odd.  Let
$n=2k+1$ and put $m=n-1=2k$.  Define
\[
 \widetilde a(x)=\sum_{i=0}^{2m-2}a_{i+1}x^i,
 \qquad
 \widetilde b(x)=\sum_{j=0}^{m-1}b_jx^j.
\]
Thus $\widetilde a$ drops the lowest coefficient of $a$ and shifts the
remaining coefficients down, while $\widetilde b$ drops the last coefficient
of $b$.

\begin{algorithm}[Odd-size reduction for Toom-4/2]\label{alg:middle-toom42-odd}
Input the balanced middle-product problem of size $n=2k+1$.
\begin{enumerate}
 \item Compute
 \[
    z\leftarrow\operatorname{MP}_{m}(\widetilde a,\widetilde b),
    \qquad m=n-1.
 \]
 \item For $0\leq i<m$, set
 \[
    z_i\leftarrow z_i+a_i b_m.
 \]
 \item Compute the last output coefficient directly as
 \[
    z_m\leftarrow\sum_{j=0}^{m}a_{2m-j}b_j.
 \]
 \item Return $\sum_{i=0}^{m}z_i x^i$.
\end{enumerate}
\end{algorithm}

If a formally present coefficient is outside the actual operand length, it is
simply zero.

\begin{theorem}\label{thm:middle-toom42-odd}
Algorithm~\ref{alg:middle-toom42-odd} returns
$\operatorname{MP}_n(a,b)$.  Beyond the size-$m$ recursive middle product, it
uses at most $2n-1$ coefficient multiplications and $O(n)$ additions.
Consequently it preserves the $\Theta(n^{\log_2 3})$ coefficient-operation
complexity of recursive Toom-4/2.  Its extra bit cost is
\[
 O\!\left(n\bigl(\MZ(\tau)+\tau+\log n\bigr)\right),
\]
which is lower order than the recursive Toom-4/2 bound.
\end{theorem}

\begin{proof}
For $0\leq i<m$, coefficient $i$ of the core is
\[
 \sum_{j=0}^{m-1}a_{m+i-j}b_j,
\]
because shifting $a$ down by one converts local middle degree $m-1+i$ into
original degree $m+i=n-1+i$.  The only missing contribution to the desired
coefficient is the term with $j=m$, namely $a_i b_m$.  This proves the
correction in step 2.

The last desired coefficient corresponds to original product degree
$n-1+m=2m$ and is therefore
\[
 \sum_{j=0}^{m}a_{2m-j}b_j,
\]
which is step 3.  There are $m$ correction products in step 2 and $n$ in step
3, for $m+n=2n-1$ in total.  The complexity statements follow immediately
from Theorem~\ref{thm:middle-toom42} and the coefficient-size bound
\eqref{eq:middle-coeff-bound}.
\end{proof}

\subsection{Toom-6/3: transposed Toom-3}

Let $n=3k$.  With $F_0,\ldots,F_4$ and $G_0,G_1,G_2$ as in
Lemma~\ref{lem:middle-block}, write
\[
 \operatorname{MP}_{3k}(a,b)=C_0+x^kC_1+x^{2k}C_2.
\]
The block identity gives
\begin{align}
 C_0&=\operatorname{MP}_k(F_0,G_2)
      +\operatorname{MP}_k(F_1,G_1)
      +\operatorname{MP}_k(F_2,G_0),\label{eq:t63-c0}\\
 C_1&=\operatorname{MP}_k(F_1,G_2)
      +\operatorname{MP}_k(F_2,G_1)
      +\operatorname{MP}_k(F_3,G_0),\label{eq:t63-c1}\\
 C_2&=\operatorname{MP}_k(F_2,G_2)
      +\operatorname{MP}_k(F_3,G_1)
      +\operatorname{MP}_k(F_4,G_0).\label{eq:t63-c2}
\end{align}
A direct evaluation would use nine recursive products.  Transposed Toom-3
reduces this to five.

\begin{algorithm}[Toom-6/3 middle product]\label{alg:middle-toom63}
Input $a$ of formal length $6k-1$ and $b$ of formal length $3k$.
Form the five recursive products
\begin{align*}
 P_0&\leftarrow\operatorname{MP}_k
 (2F_0-F_1-2F_2+F_3,\ G_2),\\
 P_1&\leftarrow\operatorname{MP}_k
 (F_3-2F_1-F_2,\ G_0+G_1+G_2),\\
 P_2&\leftarrow\operatorname{MP}_k
 (F_3+2F_1-3F_2,\ G_0-G_1+G_2),\\
 P_3&\leftarrow\operatorname{MP}_k
 (F_3-F_1,\ 4G_0+2G_1+G_2),\\
 P_4&\leftarrow\operatorname{MP}_k
 (F_4-F_2-2F_3+2F_1,\ G_0).
\end{align*}
Interpolate by
\begin{align*}
 Q&\leftarrow(P_3-P_2)/3,\\
 H&\leftarrow(P_2-P_1)/2,\\
 C_1&\leftarrow H+Q,\\
 C_0&\leftarrow(P_0-P_1+Q)/2,\\
 C_2&\leftarrow Q+C_1+P_4.
\end{align*}
Return $C_0+x^kC_1+x^{2k}C_2$.
\end{algorithm}

The divisions by $2$ are exact halvings and may be implemented as shifts over
$\ZZ$.  The division by $3$ is also exact; no division by $6$ is required in
this transposed schedule.

\begin{theorem}\label{thm:middle-toom63-correct}
Algorithm~\ref{alg:middle-toom63} returns
$\operatorname{MP}_{3k}(a,b)$.  All three indicated exact divisions are valid
in $\ZZ[x]$.
\end{theorem}

\begin{proof}
For brevity write
\[
 D_{ij}=\operatorname{MP}_k(F_i,G_j).
\]
Then \eqref{eq:t63-c0}--\eqref{eq:t63-c2} say
\[
 C_0=D_{02}+D_{11}+D_{20},\quad
 C_1=D_{12}+D_{21}+D_{30},\quad
 C_2=D_{22}+D_{31}+D_{40}.
\]
Expanding the five products by bilinearity gives
\begin{align*}
 P_0={}&2D_{02}-D_{12}-2D_{22}+D_{32},\\
 P_1={}&-2D_{10}-2D_{11}-2D_{12}-D_{20}-D_{21}-D_{22}
          +D_{30}+D_{31}+D_{32},\\
 P_2={}&2D_{10}-2D_{11}+2D_{12}-3D_{20}+3D_{21}-3D_{22}
          +D_{30}-D_{31}+D_{32},\\
 P_3={}&-4D_{10}-2D_{11}-D_{12}+4D_{30}+2D_{31}+D_{32},\\
 P_4={}&2D_{10}-D_{20}-2D_{30}+D_{40}.
\end{align*}
Consequently
\begin{align*}
 Q={}&-2D_{10}-D_{12}+D_{20}-D_{21}+D_{22}+D_{30}+D_{31},\\
 H={}&2D_{10}+2D_{12}-D_{20}+2D_{21}-D_{22}-D_{31}.
\end{align*}
These identities prove at once that $P_3-P_2$ is divisible by $3$ and
$P_2-P_1$ by $2$.  Adding gives
\[
 H+Q=D_{12}+D_{21}+D_{30}=C_1.
\]
A further direct simplification gives
\[
 P_0-P_1+Q=2(D_{02}+D_{11}+D_{20})=2C_0,
\]
so the second halving is exact and recovers $C_0$.  Finally,
\[
 Q+C_1+P_4=D_{22}+D_{31}+D_{40}=C_2.
\]
Thus all three output blocks are recovered exactly.
\end{proof}

\begin{theorem}\label{thm:middle-toom63-complexity}
Recursive Toom-6/3 middle multiplication satisfies
\[
 T_{63}(n)=5T_{63}(n/3)+O(n),
\]
and hence
\[
 T_{63}(n)=\Theta\!\left(n^{\log_3 5}\right),
 \qquad \log_3 5\approx1.465.
\]
For input coefficients of at most $\tau$ bits, its bit complexity is
\[
 O\!\left(n^{\log_3 5}\MZ(\tau+\log n)\right).
\]
\end{theorem}

\begin{proof}
The algorithm makes five recursive middle products of size $n/3$.  Every
other operation is coefficientwise addition, subtraction, multiplication by a
fixed small integer, or exact division by $2$ or $3$, and therefore costs
$O(n)$ coefficient operations at one level.  The recurrence and exponent
follow from the Master Theorem.

At each recursive level, every evaluated coefficient is an integer linear
combination with coefficients bounded by an absolute constant; the largest
sum of absolute multipliers in the displayed transforms is also bounded by an
absolute constant.  Thus each level increases coefficient bit size by only
$O(1)$, and a recursion path of depth $O(\log n)$ produces
$O(\tau+\log n)$-bit coefficients.  The exact divisions and all linear
transforms are linear-time in these coefficient sizes and are absorbed by the
$\Theta(n^{\log_3 5})$ recursive work because $\MZ(t)=\Omega(t)$.
\end{proof}

The name Toom-6/3 again records the approximate operand lengths: the long
operand has $6k-1$ coefficients and the short operand $3k$, while each
recursive call has shape $(2k-1)\times k\to k$.  The five recursive calls are
the transposed counterpart of the five products in ordinary Toom-3.

\subsection{Sizes not divisible by three}

Toom-6/3 can be extended to every size with only linear correction work.  Let
\[
 n=m+r,\qquad 3\mid m,\qquad r\in\{1,2\}.
\]
Define
\[
 \widetilde a(x)=\sum_{i=0}^{2m-2}a_{i+r}x^i,
 \qquad
 \widetilde b(x)=\sum_{j=0}^{m-1}b_jx^j.
\]
The size-$m$ core supplies the first $m$ output coefficients arising from the
first $m$ coefficients of the short operand.  Only $r$ omitted short-operand
rows and $r$ final output diagonals remain.

\begin{algorithm}[Residue correction for Toom-6/3]\label{alg:middle-toom63-tail}
Input a balanced size-$n$ middle product with $n=m+r$, $3\mid m$ and
$r\in\{1,2\}$.
\begin{enumerate}
 \item Compute
 \[
    z\leftarrow\operatorname{MP}_m(\widetilde a,\widetilde b).
 \]
 \item For $0\leq i<m$, add the omitted short-operand tail:
 \[
    z_i\leftarrow z_i+
       \sum_{j=m}^{n-1}a_{n-1+i-j}b_j.
 \]
 \item For $m\leq i<n$, compute the final output diagonals directly:
 \[
    z_i\leftarrow\sum_{j=0}^{n-1}a_{n-1+i-j}b_j.
 \]
 \item Return $\sum_{i=0}^{n-1}z_ix^i$.
\end{enumerate}
\end{algorithm}

\begin{theorem}\label{thm:middle-toom63-tail}
Algorithm~\ref{alg:middle-toom63-tail} returns
$\operatorname{MP}_n(a,b)$.  Its correction phase uses exactly
\[
 rm+rn=r(2n-r)
\]
coefficient products when all formal coefficients are present.  Since
$r\leq2$, this is $O(n)$ work.  Therefore the correction preserves both the
$\Theta(n^{\log_3 5})$ coefficient-operation complexity and the bit-complexity
bound of Theorem~\ref{thm:middle-toom63-complexity}.
\end{theorem}

\begin{proof}
For $0\leq i<m$, coefficient $i$ of the size-$m$ core is
\[
 \sum_{j=0}^{m-1}a_{r+m-1+i-j}b_j
 =\sum_{j=0}^{m-1}a_{n-1+i-j}b_j.
\]
Thus it is exactly the desired size-$n$ middle coefficient with the terms
$j=m,\ldots,n-1$ omitted, proving step 2.

The core has only $m$ outputs, so the last $r$ desired coefficients have not
been computed at all.  Their defining formula is precisely the sum in step 3.
The first correction rectangle contains $rm$ coefficient products and the
$r$ final diagonals contain $rn$, giving $r(m+n)=r(2n-r)$.  This is linear in
$n$ because $r$ is at most two.  Its bit cost is
$O(n(\MZ(\tau)+\tau+\log n))$, which is lower order than the recursive
Toom-6/3 cost.
\end{proof}

\subsection{Summary of middle-product complexities}

Let $n$ denote the number of requested balanced middle coefficients and let
$\tau$ bound the input coefficient bit size.  The algorithms above have the
following asymptotic costs:
\[
\begin{array}{c|c|c}
\text{algorithm} & \text{coefficient operations} & \text{bit complexity}\\
\hline
\text{classical} & \Theta(n^2) &
 O(n^2(\MZ(\tau)+\tau+\log n))\\
\text{Toom-4/2} & \Theta(n^{\log_2 3}) &
 O(n^{\log_2 3}\MZ(\tau+\log n))\\
\text{Toom-6/3} & \Theta(n^{\log_3 5}) &
 O(n^{\log_3 5}\MZ(\tau+\log n)).
\end{array}
\]
The odd Toom-4/2 reduction and the two nonzero residue classes for Toom-6/3
add only $O(n)$ coefficient work and do not alter these exponents.  More
generally, the transposition principle shows that whenever an ordinary
multiplication algorithm is expressed as a suitable bilinear straight-line
program, a middle product with the same multiplicative rank and asymptotic
multiplication complexity exists.  The explicit integral schedules above are
the versions needed here over $\ZZ[x]$.



\section{Classical and divide-and-conquer division}\label{sec:classical-division}

Polynomial division over $\ZZ[x]$ differs from division over a field because a
quotient need not exist in $\ZZ[x]$.  This section treats only divisions whose
required coefficient quotients are exact.  Pseudo-division, which removes this
restriction by scaling the dividend, is treated later.

Let
\[
 a(x)=\sum_{i=0}^{m-1}a_i x^i,
 \qquad
 b(x)=\sum_{j=0}^{d-1}b_j x^j,
 \qquad b_{d-1}\ne0,
\]
with $m\geq d$, and put
\[
 n=m-d+1.
\]
Thus $n$ is the maximum possible length of the quotient.  A division with
remainder in $\ZZ[x]$ is an identity
\begin{equation}\label{eq:division-identity}
 a=bq+r,
 \qquad
 q,r\in\ZZ[x],
 \qquad
 \deg r<\deg b=d-1.
\end{equation}
If such $q,r$ exist, they are unique because $\ZZ[x]$ is an integral domain.
When $m<d$, the quotient is zero and the remainder is $a$, so only $m\geq d$
is interesting below.

We shall also use truncated formal power-series division.  If $b_0\ne0$, an
$n$-term series quotient is a polynomial $q$ of length at most $n$ satisfying
\begin{equation}\label{eq:series-division-definition}
 bq\equiv a\pmod{x^n}.
\end{equation}
Over $\ZZ$, such a quotient need not exist.  Every division by $b_0$ appearing
below is therefore required to be exact.  This convention is useful even when
$b_0$ is not a unit: the algorithms succeed precisely on inputs for which the
successive triangular equations have integral solutions.

For a polynomial $f$ of length at most $s$, define its reversal at formal
length $s$ by
\[
 \operatorname{rev}_s(f)=x^{s-1}f(x^{-1}).
\]
Leading zeroes are significant in this notation.  Reversal will turn ordinary
polynomial quotient computation into truncated series division.

For length-only complexity, let $\mathsf M(n)$ denote the coefficient-operation
cost of multiplying two length-$n$ polynomials and let $\mathsf P(n)$ denote
the cost of a balanced middle product of size $n$.  The transposition results
of the preceding section give $\mathsf P(n)=O(\mathsf M(n))$ for the
multiplication families used here.  For bit complexity, let $\tau$ bound the
bit size of the input coefficients and let $\rho$ bound the bit size of the
coefficients of the quotient being computed.  Intermediate product
coefficients then have
\[
 O(\tau+\rho+\log n)
\]
bits.  A more explicit worst-case bound for $\rho$ is proved below.

Classical polynomial and power-series division are standard; see, for example,
Brent and Zimmermann \cite{BrentZimmermann2010} and von zur Gathen and Gerhard
\cite{vonZurGathenGerhard2013}.  The use of middle products to accelerate power
series division is developed systematically by Hanrot, Quercia and Zimmermann
\cite{HanrotQuerciaZimmermann2004}.

\subsection{Classical polynomial division with remainder}

The classical algorithm removes the leading term of the current remainder at
each step.  Since the leading coefficient of the divisor need not be a unit,
the cancellation is possible in $\ZZ[x]$ only when the corresponding scalar
division is exact.

\begin{algorithm}[Classical polynomial division]\label{alg:division-classical}
Input $a,b\in\ZZ[x]$ with $b\ne0$ and $m\geq d$.
\begin{enumerate}
 \item Set $r\leftarrow a$, $q\leftarrow0$ and $n\leftarrow m-d+1$.
 \item For $k=n-1,n-2,\ldots,0$:
 \begin{enumerate}
  \item Set
  \[
   t\leftarrow \frac{[x^{d-1+k}]r}{b_{d-1}}.
  \]
  If this quotient is not integral, report that the division does not exist in
  $\ZZ[x]$.
  \item Set $q_k\leftarrow t$ and
  \[
   r\leftarrow r-tx^k b.
  \]
 \end{enumerate}
 \item Return $(q,r)$.
\end{enumerate}
\end{algorithm}

\begin{theorem}\label{thm:division-classical-correctness}
Algorithm~\ref{alg:division-classical} succeeds if and only if there exist
$q,r\in\ZZ[x]$ satisfying \eqref{eq:division-identity}.  When it succeeds, it
returns that unique pair.
\end{theorem}

\begin{proof}
Before the iteration with index $k$, all coefficients of the current remainder
above degree $d-1+k$ have already been cancelled.  Any quotient satisfying
\eqref{eq:division-identity} must therefore have
\[
 q_k=\frac{[x^{d-1+k}]r}{b_{d-1}},
\]
because no lower quotient coefficient can affect degree $d-1+k$.  Hence an
inexact scalar quotient proves that no quotient in $\ZZ[x]$ can exist.

If the scalar quotient is exact, subtracting $q_kx^kb$ cancels the coefficient
of degree $d-1+k$ and changes no higher coefficient.  Induction on descending
$k$ therefore proves that after the final iteration every coefficient of $r$
of degree at least $d-1$ is zero and
\[
 a=bq+r.
\]
Thus the returned pair satisfies \eqref{eq:division-identity}.  The preceding
argument also shows that every $q_k$ is forced, proving uniqueness and the
``if and only if'' statement.
\end{proof}

\begin{theorem}\label{thm:division-classical-complexity}
Algorithm~\ref{alg:division-classical} performs $n$ exact scalar divisions and,
in its direct form, $nd$ coefficient multiplications and subtractions.  Hence
its coefficient-operation complexity is $\Theta(nd)$.

If all input coefficients have at most $\tau$ bits and all quotient
coefficients have at most $\rho$ bits, its bit complexity is
\[
 O\!\left(nd\,\MZ(\tau+\rho+\log n)\right).
\]
Without an a priori bound on $\rho$, one may use the worst-case estimate
\[
 \rho=O\!\left(n(\tau+1)\right),
\]
which gives the rigorous but deliberately coarse bound
\[
 O\!\left(nd\,\MZ(n(\tau+1)+\log n)\right).
\]
\end{theorem}

\begin{proof}
There are $n$ quotient coefficients.  Each iteration subtracts a scalar
multiple of all $d$ coefficients of $b$, giving the coefficient-operation
count.

For the bit bound, a product $q_kb_j$ has at most
$O(\tau+\rho)$ bits, and a coefficient may accumulate at most $n$ such terms,
so $O(\tau+\rho+\log n)$ bits suffice for the arithmetic.  Standard fast exact
integer division of operands of this size costs
$O(\MZ(\tau+\rho+\log n))$, and the additions are absorbed because
$\MZ(t)=\Omega(t)$.

For the final estimate, read the quotient equations from high degree to low
degree, or equivalently reverse them as in
Theorem~\ref{thm:division-reversal} below.  If $A=2^\tau$ bounds the absolute
values of the relevant dividend coefficients and $B=2^\tau$ bounds those of
the divisor, exact division by a nonzero integer cannot increase absolute
value, and the triangular recurrence gives
\[
 |q_i|\leq A+B\sum_{j<i}|q_j|.
\]
Induction yields
\[
 |q_i|\leq A(1+B)^i.
\]
Therefore
\[
 \bits(q_i)\leq \tau+i\log_2(1+2^\tau)+O(1)
              =O((i+1)(\tau+1)),
\]
and $i<n$ gives the stated worst-case estimate.
\end{proof}

If only the quotient is wanted, it is unnecessary to update coefficients which
can only belong to the final remainder.  The companion implementation includes
this triangular quotient-only specialization.  Equivalently, it is classical
series division applied after reversal, as shown below.  If the quotient has
length $n$ and the divisor length is $d$, it uses
\begin{equation}\label{eq:triangular-quotient-count}
 \sum_{i=0}^{n-1}\min(i,d-1)
 =
 \begin{cases}
  n(n-1)/2, & n\leq d,\\[2mm]
  (d-1)(n-d/2), & n>d,
 \end{cases}
\end{equation}
coefficient multiply-subtract operations, in addition to $n$ exact scalar
divisions.

\subsection{Classical truncated series division}

Equation~\eqref{eq:series-division-definition} is triangular in the
coefficients of $q$.  Writing
\[
 q=\sum_{i=0}^{n-1}q_ix^i,
\]
the coefficient of $x^i$ gives
\begin{equation}\label{eq:series-division-recurrence}
 q_i=\frac{a_i-\displaystyle\sum_{j=1}^{\min(i,d-1)}b_jq_{i-j}}{b_0}.
\end{equation}
This yields the classical power-series quotient directly.

\begin{algorithm}[Classical truncated series quotient]\label{alg:series-quo-classical}
Input $a,b\in\ZZ[x]$, $b_0\ne0$, and a precision $n\geq0$.
\begin{enumerate}
 \item For $i=0,1,\ldots,n-1$, form
 \[
  u_i\leftarrow a_i-
  \sum_{j=1}^{\min(i,d-1)}b_jq_{i-j},
 \]
 where missing coefficients of $a$ are zero.
 \item If $u_i/b_0$ is not integral, report that no quotient satisfying
 \eqref{eq:series-division-definition} exists in $\ZZ[x]$; otherwise set
 \[
  q_i\leftarrow u_i/b_0.
 \]
 \item Return $q=\sum_{i=0}^{n-1}q_ix^i$.
\end{enumerate}
\end{algorithm}

\begin{theorem}\label{thm:series-quo-classical}
Algorithm~\ref{alg:series-quo-classical} succeeds if and only if there is a
$q\in\ZZ[x]$ of length at most $n$ satisfying
$bq\equiv a\pmod{x^n}$.  The quotient is then unique.  The algorithm uses $n$
exact scalar divisions and
\[
 C(n,d)=\sum_{i=0}^{n-1}\min(i,d-1)
\]
coefficient multiplications and accumulations.  In particular,
\[
 C(n,d)=n(n-1)/2
\]
when $d\geq n$.
\end{theorem}

\begin{proof}
Assume $q_0,\ldots,q_{i-1}$ have already been determined.  The coefficient of
$x^i$ in $bq$ is
\[
 b_0q_i+\sum_{j=1}^{\min(i,d-1)}b_jq_{i-j}.
\]
Thus equality with $a_i$ forces exactly
\eqref{eq:series-division-recurrence}.  Hence an inexact division proves that
no integral series quotient exists, while an exact division determines $q_i$
uniquely.  Induction on $i$ proves correctness and uniqueness.  At stage $i$
there are $\min(i,d-1)$ coefficient products, giving the operation count.
\end{proof}

\begin{corollary}\label{cor:series-classical-bit}
If the input coefficients have at most $\tau$ bits and the quotient
coefficients have at most $\rho$ bits, classical truncated series division has
bit complexity
\[
 O\!\left((C(n,d)+n)\MZ(\tau+\rho+\log n)\right).
\]
In the balanced case $d\geq n$ this is
\[
 O\!\left(n^2\MZ(\tau+\rho+\log n)\right).
\]
Moreover every exact quotient satisfies the worst-case estimate
\[
 \rho=O\!\left(n(\tau+1)\right).
\]
\end{corollary}

\begin{proof}
The arithmetic-size estimate is the same as in
Theorem~\ref{thm:division-classical-complexity}.  For completeness, if
$A=B=2^\tau$, recurrence \eqref{eq:series-division-recurrence} and exactness
give
\[
 |q_i|\leq A+B\sum_{j<i}|q_j|.
\]
The induction
\[
 |q_i|\leq A(1+B)^i
\]
proves the stated coefficient-growth bound.  Each coefficient product,
accumulation and exact division therefore fits in the displayed bit-complexity
bound.
\end{proof}

The linear growth of quotient bit size in this worst-case estimate is real, not
an artefact of the proof.  For example, inversion of $1-2x$ produces the
series $1+2x+4x^2+\cdots$.  Length-only and bit-complexity estimates should
therefore be kept conceptually separate.

\subsection{Divide-and-conquer series division}

The classical recurrence determines quotient coefficients one at a time.
Divide-and-conquer division determines a low block of quotient coefficients,
uses one middle product to compute the corresponding error block, and then
recurses on the remaining high block.

Let $n=k+h$, where
\[
 k=\lfloor n/2\rfloor,
 \qquad
 h=n-k.
\]
Write the desired quotient as
\[
 q=q_0+x^kq_1,
 \qquad
 \len(q_0)\leq k,
 \qquad
 \len(q_1)\leq h.
\]
First compute
\[
 bq_0\equiv a\pmod{x^k}.
\]
Hence $a-bq_0$ has zero coefficients below degree $k$.  Define the next error
block by
\begin{equation}\label{eq:series-dc-error}
 e=\left(\frac{a-bq_0}{x^k}\right)\bmod x^h.
\end{equation}
Only coefficients $k,\ldots,n-1$ of $bq_0$ are required to form $e$; this is a
middle-product window, not a full product.  The remaining quotient block is
then characterized by
\[
 bq_1\equiv e\pmod{x^h}.
\]

\begin{algorithm}[Divide-and-conquer truncated series quotient]\label{alg:series-quo-dc}
Input $a,b\in\ZZ[x]$, $b_0\ne0$, and precision $n$.
\begin{enumerate}
 \item If $n$ is a base-case size, use
 Algorithm~\ref{alg:series-quo-classical}.
 \item Put $k\leftarrow\lfloor n/2\rfloor$ and $h\leftarrow n-k$.
 \item Recursively compute
 \[
  q_0\leftarrow \operatorname{SeriesQuo}(a,b,k).
 \]
 \item Compute only coefficients $k,\ldots,k+h-1$ of $bq_0$, and use them to
 form $e$ as in \eqref{eq:series-dc-error}.
 \item Recursively compute
 \[
  q_1\leftarrow \operatorname{SeriesQuo}(e,b,h).
 \]
 \item Return $q_0+x^kq_1$.
\end{enumerate}
\end{algorithm}

\begin{theorem}\label{thm:series-quo-dc-correctness}
Algorithm~\ref{alg:series-quo-dc} returns the unique integral quotient $q$
satisfying
\[
 bq\equiv a\pmod{x^n},
\]
whenever that quotient exists, and otherwise detects an inexact scalar division
in a base case.
\end{theorem}

\begin{proof}
By induction, the first recursive call returns $q_0$ satisfying
\[
 bq_0\equiv a\pmod{x^k}.
\]
Consequently the coefficients below degree $k$ of $a-bq_0$ vanish, so the
truncated quotient in \eqref{eq:series-dc-error} is well defined.  The second
recursive call gives
\[
 bq_1\equiv e\pmod{x^h}.
\]
Multiplying by $x^k$ and using the definition of $e$ gives
\[
 x^kbq_1\equiv a-bq_0\pmod{x^{k+h}}.
\]
Therefore
\[
 b(q_0+x^kq_1)\equiv a\pmod{x^n},
\]
as required.  Uniqueness follows from the same triangular argument used for
classical series division.  If an integral quotient does not exist, some
forced scalar quotient in that triangular system is nonintegral, and recursion
eventually reaches the corresponding base-case equation.
\end{proof}

\begin{theorem}\label{thm:series-quo-dc-complexity}
Let $D(n)$ be the coefficient-operation cost of
Algorithm~\ref{alg:series-quo-dc}, and let $\mathsf P(n)$ be the cost of a
balanced middle product of size $n$.  Up to rounding,
\begin{equation}\label{eq:series-dc-recurrence}
 D(n)=2D(n/2)+\mathsf P(n/2)+O(n).
\end{equation}
If
\[
 \mathsf P(n)\sim A n^\alpha
 \qquad (\alpha>1),
\]
then
\begin{equation}\label{eq:series-dc-leading}
 D(n)\sim \frac{A}{2^\alpha-2}n^\alpha.
\end{equation}
In particular, if $\mathsf P(n)=\Theta(\mathsf M(n))$ and
$\mathsf M(n)=\Theta(n^\alpha)$ with $\alpha>1$, series division has
$\Theta(\mathsf M(n))$ coefficient complexity.
\end{theorem}

\begin{proof}
There are two recursive quotient computations of half size.  The only
nonrecursive multiplication is the window of $bq_0$ used in
\eqref{eq:series-dc-error}.  When $n=2k$, its $i$th requested coefficient is
\[
 [x^{k+i}](bq_0)
 =\sum_{j=0}^{k-1} b_{k+i-j}(q_0)_j,
 \qquad 0\leq i<k.
\]
If $\widetilde b_t=b_{t+1}$, this is exactly coefficient $i$ of
$\operatorname{MP}_k(\widetilde b,q_0)$.  Thus the correction is one balanced
middle product of size $k=n/2$.  If $n=2k+1$, then $h=k+1$ and the requested
window $k,\ldots,2k$ is directly a balanced middle product of size $h$ after
regarding $q_0$ as having one leading zero coefficient.  Thus rounding changes
only the argument from $n/2$ to $\lceil n/2\rceil$.  Forming $e$ and
concatenating the two quotient blocks cost $O(n)$ coefficient operations.  This
proves \eqref{eq:series-dc-recurrence}.

Suppose $D(n)\sim Kn^\alpha$ and $\mathsf P(n)\sim An^\alpha$.  Substitution
into the leading terms of \eqref{eq:series-dc-recurrence} gives
\[
 K=\frac{2K+A}{2^\alpha},
\]
so
\[
 K=\frac{A}{2^\alpha-2}.
\]
The same result follows by summing the recursion tree.  Because $\alpha>1$,
the $O(n)$ work at each node contributes only lower-order work.  Finally, the
transposition principle gives $\mathsf P(n)=O(\mathsf M(n))$ for the
multiplication algorithms under consideration, proving the last statement.
\end{proof}

The leading-constant formula gives useful checks on the recurrence.  With a
classical middle product, $\alpha=2$ and $A=1$, so
\[
 D(n)\sim \frac12n^2,
\]
which is exactly the leading multiplication count of the classical triangular
series quotient.  With transposed Karatsuba,
$\alpha=\log_2 3$ and $2^\alpha-2=1$, so the leading constant equals that of
one full-size Karatsuba middle product.  With transposed Toom-3,
$\alpha=\log_3 5$, the factor is
\[
 \frac1{2^{\log_3 5}-2}\approx1.31478.
\]
These constants concern the idealized smooth recurrence; practical dispatch
between algorithms is a separate implementation question.

For multiplication costs which are only quasi-linear, the recurrence need not
collapse to $O(\mathsf M(n))$.  Directly unrolling
\eqref{eq:series-dc-recurrence} gives the general bound
\[
 D(n)=O\!\left(
  \sum_{j\geq0}2^j\mathsf P\!\left(\frac{n}{2^{j+1}}\right)
  +n\log n
 \right).
\]
For example, if $\mathsf P(n)=O(n\log n)$, this gives
$D(n)=O(n\log^2 n)$.  The polynomial-exponent middle products used in the
companion implementation lie in the simpler $\alpha>1$ regime above.

\begin{theorem}\label{thm:series-quo-dc-bit}
Suppose input coefficients have at most $\tau$ bits and quotient coefficients
have at most $\rho$ bits.  If a size-$n$ middle product on coefficients of at
most $\lambda$ bits costs
\[
 O\!\left(n^\alpha\MZ(\lambda+\log n)\right)
 \qquad (\alpha>1),
\]
then divide-and-conquer series division has bit complexity
\[
 O\!\left(
 n^\alpha\MZ(\tau+\rho+\log n)
 \right).
\]
Using the worst-case coefficient-growth estimate
$\rho=O(n(\tau+1))$ gives
\[
 O\!\left(
 n^\alpha\MZ(n(\tau+1)+\log n)
 \right).
\]
\end{theorem}

\begin{proof}
Every recursively generated dividend block is a coefficient block of $b$ times
the corresponding remaining quotient block.  Hence its coefficients, and the
coefficients entering every correction middle product, have
$O(\tau+\rho+\log n)$ bits.  Using this global bound at every recursion node,
the bit recurrence has the same shape as
\eqref{eq:series-dc-recurrence}, with the middle-product term replaced by its
stated bit cost.  Since $\alpha>1$, the recursion-tree sum is dominated by the
same $n^\alpha$ order.  The second display follows from
Corollary~\ref{cor:series-classical-bit}.
\end{proof}

\subsection{Ordinary quotient by reversal}

Ordinary polynomial quotient computation is a series-division problem at the
high end of the operands.  Reversal moves that triangular system to the low
end.

\begin{theorem}[Reversal identity for quotient]\label{thm:division-reversal}
Assume
\[
 a=bq+r,
 \qquad
 \deg r<d-1,
\]
where $m=\len(a)$, $d=\len(b)$ and $n=m-d+1$.  Put
\[
 A=\operatorname{rev}_m(a)\bmod x^n,
 \qquad
 B=\operatorname{rev}_d(b),
 \qquad
 Q=\operatorname{rev}_n(q).
\]
Then
\begin{equation}\label{eq:division-reversal-congruence}
 BQ\equiv A\pmod{x^n}.
\end{equation}
Conversely, an integral $Q$ satisfying
\eqref{eq:division-reversal-congruence} determines the unique ordinary
quotient by $q=\operatorname{rev}_n(Q)$.
\end{theorem}

\begin{proof}
Because
\[
 m-1=(d-1)+(n-1),
\]
reversal turns the product $bq$ into an exact product:
\[
 \operatorname{rev}_m(bq)
 =\operatorname{rev}_d(b)\operatorname{rev}_n(q)=BQ.
\]
On the other hand,
\[
 x^{m-1}r(x^{-1})
\]
is divisible by $x^n$, because every term of $r$ has degree at most $d-2$ and
therefore becomes a term of degree at least
\[
 m-1-(d-2)=n.
\]
Reversing $a=bq+r$ and reducing modulo $x^n$ gives
\eqref{eq:division-reversal-congruence}.

Conversely, the congruence determines $Q$ uniquely by the triangular
series-division equations because the constant coefficient of $B$ is
$b_{d-1}\ne0$.  Reversing $Q$ therefore gives the only possible ordinary
quotient.  Subtracting $bq$ from $a$ leaves no coefficient of degree at least
$d-1$, so the resulting remainder has degree below $d-1$.
\end{proof}

\begin{algorithm}[Ordinary quotient by reversed series division]\label{alg:quotient-reversal}
Input $a,b\in\ZZ[x]$ with $m\geq d$ and $n=m-d+1$.
\begin{enumerate}
 \item Take the highest $n$ coefficients of $a$ and reverse them to obtain
 \[
  A=\operatorname{rev}_m(a)\bmod x^n.
 \]
 \item Take the highest $\min(d,n)$ coefficients of $b$ and reverse them.
 This gives exactly $B\bmod x^n$, where $B=\operatorname{rev}_d(b)$.
 \item Compute the truncated series quotient
 \[
  Q\leftarrow A/B\pmod{x^n}
 \]
 by Algorithm~\ref{alg:series-quo-dc} or its classical base case.
 \item Return $q=\operatorname{rev}_n(Q)$.
\end{enumerate}
\end{algorithm}

\begin{corollary}\label{cor:quotient-reversal-complexity}
If reversing and copying $n$ coefficients costs $O(n)$ and the relevant
series quotient costs $D(n)$, ordinary quotient computation by reversal costs
\[
 D(n)+O(n)
\]
coefficient operations.  Thus, with
$\mathsf P(n)=\Theta(n^\alpha)$ for $\alpha>1$, the quotient has
$\Theta(n^\alpha)$ coefficient complexity, independent of the total divisor
length except for the highest $\min(d,n)$ coefficients which can enter the
triangular system.
\end{corollary}

\begin{proof}
Theorem~\ref{thm:division-reversal} proves correctness.  Constructing the
truncated reversed dividend, the required reversed divisor tail, and the final
reversal of $Q$ touches only $O(n)$ coefficients.  The remaining work is the
series quotient.
\end{proof}

\subsection{Recovering the remainder}

Once a quotient $q$ is known, the remainder is mathematically trivial:
\begin{equation}\label{eq:remainder-recovery}
 r=a-bq.
\end{equation}
The point is computational: a straightforward quotient-and-remainder routine
forms the complete product $bq$, whereas later algorithms will show that only
the low remainder window need be formed.

\begin{algorithm}[Remainder recovery by a full product]\label{alg:remainder-full-product}
Input $a,b$ and a quotient $q$ known to satisfy the high quotient equations.
\begin{enumerate}
 \item Compute $p\leftarrow bq$ by an ordinary full-product algorithm.
 \item Return $r\leftarrow a-p$.
\end{enumerate}
\end{algorithm}

\begin{theorem}\label{thm:remainder-full-product}
If $q$ is the quotient obtained from
Algorithm~\ref{alg:quotient-reversal}, then
Algorithm~\ref{alg:remainder-full-product} returns the unique remainder in
\eqref{eq:division-identity}.  Its additional coefficient-operation cost is
one product of lengths $d$ and $n$, plus $O(m)$ additions or subtractions.
For balanced $d,n=\Theta(N)$ this is $O(\mathsf M(N))$.
\end{theorem}

\begin{proof}
The reversed series quotient solves exactly the $n$ highest coefficient
equations of $a=bq+r$.  Hence all coefficients of $a-bq$ in degrees
$d-1,d,\ldots,m-1$ vanish.  Therefore $r=a-bq$ has degree at most $d-2$, so it
is the required remainder.  The complexity statement is immediate from the
two steps of the algorithm.
\end{proof}

The full-product recovery is adequate when divisor and quotient lengths are
comparable.  If the quotient is short compared with a long divisor, it may do
substantially more work than quotient computation itself.  A later section on
Newton division will replace it by a low truncated product of exactly the
remainder length.

\subsection{Complexity summary}

Let $n$ denote the requested quotient length and suppose the relevant middle
product satisfies
\[
 \mathsf P(n)=\Theta(n^\alpha),\qquad \alpha>1.
\]
Ignoring coefficient growth, the main results of this section are
\[
\begin{array}{c|c}
\text{operation} & \text{coefficient-operation complexity}\\
\hline
\text{classical balanced series quotient} & \Theta(n^2)\\
\text{D\&C series quotient} & \Theta(n^\alpha)\\
\text{ordinary quotient by reversal} & \Theta(n^\alpha)\\
\text{quotient and remainder by full recovery} &
 \Theta(n^\alpha)+O(\mathsf M(d,n)).
\end{array}
\]
Here $\mathsf M(d,n)$ denotes the cost of multiplying lengths $d$ and $n$; in
the balanced case it is simply $\mathsf M(n)$ up to constant rescaling.  With
the middle-product algorithms already established,
\[
 \alpha=2,\qquad \log_2 3,\qquad \log_3 5
\]
for the classical, Toom-4/2 and Toom-6/3 regimes respectively.

For bit complexity, if $\rho$ bounds quotient coefficient size, replace a
coefficient multiplication at the relevant scale by
\[
 \MZ(\tau+\rho+\log n).
\]
This gives, for example,
\[
 O\!\left(n^\alpha\MZ(\tau+\rho+\log n)\right)
\]
for divide-and-conquer series or ordinary quotient division.  The general
worst-case estimate $\rho=O(n(\tau+1))$ is always available, but retaining
$\rho$ separately is usually more informative because exact polynomial
quotients in applications often have much smaller coefficients than that
worst case.

\section{Mulders short division and bidirectional exact division}

The algorithms of the preceding section compute an ordinary quotient by solving
all of its triangular equations from one end.  Two further situations admit
more specialized organization.  First, when the quotient is no longer than the
divisor, Mulders' \emph{short division} computes only the quotient and avoids
forming most of the final remainder.  Second, when the division is known to be
exact, the quotient may be determined independently from its low and high ends;
a single middle-product window then verifies the equations not used by either
half computation.

Mulders introduced short multiplication and short division in
\cite{Mulders2000}.  The polynomial division variant and practical refinements
are discussed by Hart and Novocin \cite{HartNovocin2011}.  The bidirectional
idea has a close analogue in exact integer division due to Krandick and
Jebelean \cite{KrandickJebelean1996}; the polynomial version below is most
naturally expressed in terms of the series-division and middle-product
operations already developed in this chapter.

Throughout the short-division subsection, the balanced problem has a dividend
$A$ of length at most $2n-1$ and a divisor $B$ of length exactly $n$.  Hence the
quotient has length at most $n$.  As before, every scalar division required by
the quotient equations must be exact in $\ZZ$.

\subsection{Balanced Mulders short division}

The recursive split used here is deliberately the simple balanced one.  Put
\[
 n_1=\lceil n/2\rceil,
 \qquad
 n_2=\lfloor n/2\rfloor,
 \qquad
 \delta=n_1-n_2\in\{0,1\}.
\]
Write the divisor and quotient as
\begin{equation}\label{eq:mulders-block-splits}
 B=B_0+x^{n_2}B_1,
 \qquad
 Q=Q_0+x^{n_2}Q_1,
\end{equation}
where
\[
 \len(B_0),\len(Q_0)\leq n_2,
 \qquad
 \len(B_1),\len(Q_1)\leq n_1.
\]
The high divisor block $B_1$ will itself be viewed in a second way.  Define
\[
 B_3=\left\lfloor B/x^{n_1}\right\rfloor,
\]
so that $\len(B_3)=n_2$.  There is a polynomial $C$ of length at most
$\delta$ such that
\begin{equation}\label{eq:mulders-b1-b3}
 B_1=C+x^\delta B_3.
\end{equation}
For even $n$ one has $\delta=0$ and simply $B_1=B_3$.  For odd $n$,
$\delta=1$ and $C$ is the single coefficient of $B$ lying between the low
$n_2$ coefficients and the high $n_2$ coefficients.

Expanding \eqref{eq:mulders-block-splits} gives
\begin{equation}\label{eq:mulders-product-split}
 BQ
 =B_0Q_0
 +x^{n_2}(B_0Q_1+B_1Q_0)
 +x^{2n_2}B_1Q_1.
\end{equation}
The algorithm first obtains $Q_1$ from the very highest coefficients, then uses
one middle product and the remainder of that first division to manufacture a
second short-division problem for $Q_0$.

Define the synthetic polynomial
\begin{equation}\label{eq:mulders-U}
 U=
 \sum_{i=0}^{n_1-1} A_{n+n_2-1+i}x^{n_1-1+i},
\end{equation}
where coefficients beyond the length of $A$ are understood to be zero.  Thus
$U$ has $n_1-1$ initial zero coefficients followed by the highest $n_1$
coefficients of $A$.

\begin{lemma}\label{lem:mulders-high-block}
Suppose $A=BQ+R$ with $\deg R<n-1$.  If
\[
 U=B_1Q_1+R_1,
 \qquad \deg R_1<n_1-1,
\]
is the ordinary quotient-and-remainder division of $U$ by $B_1$, then the
quotient is precisely the high block $Q_1$ in
\eqref{eq:mulders-block-splits}.
\end{lemma}

\begin{proof}
The three terms in \eqref{eq:mulders-product-split} other than
$x^{2n_2}B_1Q_1$ have degree at most
\[
 n+n_2-2.
\]
The remainder $R$ has still smaller degree, at most $n-2$.  Consequently the
coefficients of $A$ in degrees
\[
 n+n_2-1,\ldots,2n-2
\]
come only from $x^{2n_2}B_1Q_1$.  After removing the shift $2n_2$, these are
exactly the coefficients of $B_1Q_1$ in degrees
\[
 n_1-1,\ldots,2n_1-2.
\]
By construction, $U$ has the same coefficients in those positions.  Ordinary
polynomial quotient computation depends only on this triangular set of top
coefficients, so dividing $U$ by $B_1$ yields $Q_1$.  The lower coefficients
which were replaced by zero affect only the remainder $R_1$.
\end{proof}

The remainder $R_1$ is not discarded.  It records exactly the low part of
$B_1Q_1$ which was lost when $U$ was formed from high coefficients only.  Let
\[
 H_i=[x^{n_1-1+i}]B_0Q_1,
 \qquad 0\leq i<n_2,
\]
so the vector $(H_0,\ldots,H_{n_2-1})$ is one middle-product window.  Extend
the coefficient sequence of $R_1$ by zero outside its natural range, and form
a second synthetic polynomial $T$ whose low $n_2-1$ coefficients are zero and
whose high coefficients are
\begin{equation}\label{eq:mulders-T-def}
 [x^{n_2-1+i}]T
 =A_{n-1+i}-H_i+[x^{\delta-1+i}]R_1,
 \qquad 0\leq i<n_2.
\end{equation}

\begin{lemma}\label{lem:mulders-low-block}
Under the hypotheses of Lemma~\ref{lem:mulders-high-block}, the highest
$n_2$ coefficients of $T$ agree with the highest $n_2$ coefficients of
$B_3Q_0$.  Hence the ordinary quotient of $T$ by $B_3$ is $Q_0$.
\end{lemma}

\begin{proof}
Fix $0\leq i<n_2$ and consider degree $n-1+i$ of
\eqref{eq:mulders-product-split}.  The term $B_0Q_0$ has degree at most
$2n_2-2<n-1$, so it contributes nothing.  The term $x^{n_2}B_0Q_1$
contributes $H_i$.

For the other cross term, the relevant coefficient is
\[
 [x^{n_1-1+i}]B_1Q_0.
\]
Using \eqref{eq:mulders-b1-b3}, the term $CQ_0$ has degree at most
$n_1-2$, and therefore contributes nothing at this position.  Hence
\[
 [x^{n_1-1+i}]B_1Q_0
 =[x^{n_2-1+i}]B_3Q_0.
\]
Finally, the term $x^{2n_2}B_1Q_1$ contributes
\[
 [x^{\delta-1+i}]B_1Q_1.
\]
The index $\delta-1+i$ is at most $n_1-2$, where $U$ is zero.  Since
$R_1=U-B_1Q_1$, this contribution equals
\[
 -[x^{\delta-1+i}]R_1.
\]
Because the ordinary remainder $R$ contributes nothing in degree $n-1+i$, we
obtain
\[
 A_{n-1+i}
 =H_i+[x^{n_2-1+i}]B_3Q_0-[x^{\delta-1+i}]R_1.
\]
Rearranging gives exactly \eqref{eq:mulders-T-def}.  Thus the top $n_2$
coefficients of $T$ and $B_3Q_0$ agree.  The triangular high quotient
equations therefore force the quotient of $T$ by $B_3$ to be $Q_0$.
\end{proof}

The two lemmas give the recursive algorithm.

\begin{algorithm}[Balanced Mulders short quotient]\label{alg:mulders-balanced}
Input $A,B\in\ZZ[x]$ with $\len(A)\leq2n-1$ and $\len(B)=n$.
\begin{enumerate}
 \item If $n$ is below a base-case threshold, use classical quotient
 computation.
 \item Put
 \[
 n_1\leftarrow\lceil n/2\rceil,
 \quad n_2\leftarrow\lfloor n/2\rfloor,
 \quad \delta\leftarrow n_1-n_2,
 \]
 and form $B_0,B_1,B_3$ as above.
 \item Form $U$ from the highest $n_1$ coefficients of $A$ by
 \eqref{eq:mulders-U}.
 \item Compute the full division
 \[
 U=B_1Q_1+R_1.
 \]
 \item Compute the middle-product window
 \[
 H_i=[x^{n_1-1+i}]B_0Q_1,
 \qquad 0\leq i<n_2.
 \]
 \item Form $T$ by \eqref{eq:mulders-T-def}.
 \item Recursively compute the short quotient
 \[
 Q_0\leftarrow T\mathbin{\mathrm{quo}}B_3.
 \]
 \item Return
 \[
 Q=Q_0+x^{n_2}Q_1.
 \]
\end{enumerate}
\end{algorithm}

\begin{theorem}\label{thm:mulders-balanced-correctness}
Algorithm~\ref{alg:mulders-balanced} succeeds if and only if the ordinary
quotient of $A$ by $B$ exists in $\ZZ[x]$, and on success it returns that
quotient.
\end{theorem}

\begin{proof}
If the quotient exists, Lemma~\ref{lem:mulders-high-block} proves that the
first division returns the correct $Q_1$, and
Lemma~\ref{lem:mulders-low-block} proves that the recursive division returns
the correct $Q_0$.  Induction on $n$, with the classical quotient as base
case, proves that the returned polynomial is the ordinary quotient.

Conversely, every full or recursive quotient step solves a subset of the
ordinary triangular high coefficient equations.  If one of its required
scalar divisions is not exact, no quotient in $\ZZ[x]$ can satisfy those
equations.  Thus failure is conclusive.  If all recursive steps succeed, the
two lemmas show that the assembled $Q$ satisfies all $n$ high quotient
equations, hence $A-BQ$ has degree below $n-1=\deg B$ and $Q$ is the ordinary
quotient.
\end{proof}

The formula \eqref{eq:mulders-T-def} also explains the odd-size case.  When
$n$ is even, $\delta=0$, so the remainder correction used for the $i$th high
coefficient of $T$ is $[x^{i-1}]R_1$; the first correction is zero.  When $n$
is odd, $\delta=1$, and the correction is simply $[x^i]R_1$.  No padding or
separate odd-size algorithm is required.

\subsection{Complexity of balanced short division}

Let $S(n)$ denote the coefficient-operation cost of balanced short quotient
computation, let $F(n)$ denote the cost of a balanced full
quotient-and-remainder division of lengths $2n-1$ by $n$, and let $P(n)$
denote the cost of a middle product at scale $n$.  The recursive construction
gives
\begin{equation}\label{eq:mulders-recurrence}
 S(n)
 =S(\lfloor n/2\rfloor)
  +F(\lceil n/2\rceil)
  +P(\lceil n/2\rceil)
  +O(n).
\end{equation}
The $P$ term is written at the larger half-size merely to absorb the harmless
one-coefficient imbalance when $n$ is odd.

\begin{theorem}\label{thm:mulders-complexity}
Suppose for some $\alpha>1$ that
\[
 F(n)=\Theta(n^\alpha),
 \qquad
 P(n)=\Theta(n^\alpha).
\]
Then
\[
 S(n)=\Theta(n^\alpha).
\]
More precisely, if
\[
 F(n)\sim A n^\alpha,
 \qquad
 P(n)\sim B n^\alpha,
\]
then the balanced half-split recurrence has leading term
\begin{equation}\label{eq:mulders-leading-constant}
 S(n)\sim \frac{A+B}{2^\alpha-1}n^\alpha.
\end{equation}
\end{theorem}

\begin{proof}
Ignoring rounding, substitution of $S(n)\sim Kn^\alpha$ into
\eqref{eq:mulders-recurrence} gives
\[
 K=\frac{K+A+B}{2^\alpha},
\]
so
\[
 K=\frac{A+B}{2^\alpha-1}.
\]
Since $\alpha>1$, the $O(n)$ construction work is lower order.  The same
conclusion follows directly by summing the geometric recursion tree.
\end{proof}

As a check, consider a purely classical implementation.  A balanced full
division of lengths $2k-1$ by $k$ has leading cost $k^2$, while the particular
upper middle window used in Algorithm~\ref{alg:mulders-balanced} contains
\[
 1+2+\cdots+k=\frac{k(k+1)}2
\]
coefficient products.  Thus $A=1$, $B=1/2$ and $\alpha=2$ in
\eqref{eq:mulders-leading-constant}, giving
\[
 S(n)\sim\frac12 n^2.
\]
So even the balanced split removes one half of the leading classical
coefficient products compared with a full $n$-step long division.

In the algorithms already developed, a full division at this scale consists of
a quotient computation plus remainder recovery, so in the balanced
polynomial-exponent regimes one has
\[
 F(n)=O(\mathsf M(n)+\mathsf P(n)).
\]
Consequently Mulders short division has the same asymptotic exponent as the
underlying multiplication and middle-product algorithms.  The point of the
short algorithm is therefore a smaller constant and less discarded work, not
a new exponent.  Mulders also allows unequal split ratios; choosing and tuning
that ratio is a constant-factor question and is deliberately separate from the
balanced algorithm analyzed here.

For bit complexity, suppose input coefficients have at most $\tau$ bits and
quotient coefficients at most $\rho$ bits.  Every coefficient entering
$B_1Q_1$, the middle product $B_0Q_1$, the synthetic remainder $R_1$, and the
second dividend $T$ has
\[
 O(\tau+\rho+\log n)
\]
bits.  Hence the same recurrence proves the following.

\begin{theorem}\label{thm:mulders-bit-complexity}
If both the balanced full division and the relevant middle product at size $n$
have bit complexity
\[
 O\!\left(n^\alpha\MZ(\tau+\rho+\log n)\right),
 \qquad \alpha>1,
\]
then balanced Mulders short division has bit complexity
\[
 O\!\left(n^\alpha\MZ(\tau+\rho+\log n)\right).
\]
Using the general quotient-growth bound from the preceding section gives the
worst-case estimate
\[
 O\!\left(n^\alpha\MZ(n(\tau+1)+\log n)\right).
\]
\end{theorem}

\begin{proof}
The displayed coefficient-size bound follows because every synthetic
coefficient is either an input coefficient or a sum or difference of
coefficients of partial products of the divisor with the true quotient.
Applying the assumed bit bounds to each recursive node gives the same geometric
recurrence as \eqref{eq:mulders-recurrence}; $\alpha>1$ makes its sum the order
of the root contribution.  The final bound substitutes the quotient-growth
estimate already proved for classical and divide-and-conquer division.
\end{proof}

\subsection{Reduction of a short quotient to the balanced problem}

The balanced problem is enough whenever the quotient is no longer than the
divisor.  Let
\[
 d=\len(b),
 \qquad
 n=\len(a)-d+1,
 \qquad
 0<n\leq d.
\]
Put $s=d-n$ and define
\begin{equation}\label{eq:mulders-balanced-reduction}
 A=\left\lfloor a/x^s\right\rfloor,
 \qquad
 B=\left\lfloor b/x^s\right\rfloor,
\end{equation}
retaining only the resulting $2n-1$ and $n$ coefficients respectively.

\begin{theorem}\label{thm:mulders-balanced-reduction}
The ordinary quotient of $a$ by $b$ equals the balanced short quotient of
$A$ by $B$ in \eqref{eq:mulders-balanced-reduction}.  Thus, when $n\leq d$,
only the highest $2n-1$ coefficients of the dividend and the highest $n$
coefficients of the divisor are needed to compute the quotient.
\end{theorem}

\begin{proof}
Write
\[
 b=b_0+x^sB,
 \qquad \deg b_0<s,
\]
and suppose $a=bq+r$ with $\deg r<d-1$.  Since $\deg q<n$, both $b_0q$ and
$r$ have degree below
\[
 s+n-1=d-1.
\]
After dropping the lowest $s$ positions, their surviving contribution has
degree below $n-1$.  Therefore
\[
 A=Bq+R
\]
for some $R$ of degree below $n-1$.  Hence $q$ is exactly the ordinary quotient
of the balanced pair $(A,B)$.
\end{proof}

\begin{corollary}\label{cor:mulders-unbalanced-complexity}
If $n\leq d$, the quotient can be computed in
\[
 S(n)+O(n)
\]
coefficient operations, independently of the total divisor length apart from
locating the two high coefficient windows.  If a remainder is subsequently
wanted, forming $bq$ and subtracting it adds one product of lengths $d$ and
$n$.
\end{corollary}

When $n>d$, the balanced reduction no longer produces a length-$n$ divisor.
The companion implementation therefore returns to the general reversed
series-division algorithm of the preceding section in that regime.

\subsection{Bidirectional exact division}

Now suppose the desired division is exact:
\begin{equation}\label{eq:exact-division-identity}
 a=bq.
\end{equation}
Unlike ordinary quotient-with-remainder division, exactness exposes triangular
systems at \emph{both} ends of the product.  The low coefficients determine the
low part of $q$, while the high coefficients determine the high part
independently.  The remaining equations form one contiguous central window.

There is one preliminary issue.  If
\[
 v=\operatorname{ord}_x(b),
\]
then exact divisibility requires
\[
 a_0=\cdots=a_{v-1}=0.
\]
When this holds, write
\[
 a=x^vA,
 \qquad
 b=x^vB.
\]
Then $B(0)\ne0$ and, since $\ZZ[x]$ is an integral domain,
\begin{equation}\label{eq:valuation-exact-reduction}
 a=bq
 \quad\Longleftrightarrow\quad
 A=Bq.
\end{equation}
Thus the common power of $x$ may be removed before the bidirectional step.

Assume henceforth that $B(0)\ne0$.  Let
\[
 d=\len(B),
 \qquad
 n=\len(A)-d+1,
 \qquad
 \ell=\lceil n/2\rceil,
 \qquad
 h=\lfloor n/2\rfloor.
\]
Split
\[
 q=q_L+x^\ell q_H,
 \qquad
 \len(q_L)\leq\ell,
 \qquad
 \len(q_H)\leq h.
\]

\begin{lemma}[Independent low half]\label{lem:bidirectional-low}
If $A=Bq$, then
\[
 Bq_L\equiv A\pmod{x^\ell}.
\]
Consequently $q_L$ is determined uniquely by a length-$\ell$ truncated series
quotient from the low end.
\end{lemma}

\begin{proof}
The omitted term $x^\ell Bq_H$ is divisible by $x^\ell$, so reduction modulo
$x^\ell$ gives the congruence.  Since $B(0)\ne0$, the triangular series
quotient equations determine the coefficients of $q_L$ uniquely whenever they
are integral.
\end{proof}

\begin{lemma}[Independent high half]\label{lem:bidirectional-high}
If $A=Bq$, then the highest $h$ coefficients of $q$ are determined uniquely by
the highest $h$ coefficients of $A$ and the highest $\min(d,h)$ coefficients
of $B$.  They may be computed by reversing those coefficient windows and
performing a length-$h$ truncated series quotient.
\end{lemma}

\begin{proof}
Apply the reversal identity of
Theorem~\ref{thm:division-reversal} at precision $h$.  The top $h$ quotient
equations involve only the top $h$ coefficients of the dividend and the top
$\min(d,h)$ coefficients of the divisor.  After reversal they become the
first $h$ triangular series equations, whose solution is the reversal of the
highest $h$ quotient coefficients.
\end{proof}

The two lemmas determine all $n$ coefficients of a candidate quotient.  They
do not, however, prove exact divisibility.  The low computation enforces only
product degrees
\[
 0,1,\ldots,\ell-1,
\]
while the high computation enforces only degrees
\[
 d+\ell-1,d+\ell,\ldots,d+n-2.
\]
Exactly the central band
\begin{equation}\label{eq:bidirectional-central-window}
 \ell,\ell+1,\ldots,\ell+d-2
\end{equation}
remains unchecked.  Its length is $d-1$.

\begin{algorithm}[Bidirectional exact quotient]\label{alg:bidirectional-exact}
Input $a,b\in\ZZ[x]$ with $b\ne0$ and $\len(a)\geq\len(b)$; the
shorter-dividend case is trivial.
\begin{enumerate}
 \item Let $v\leftarrow\operatorname{ord}_x(b)$.  If any of
 $a_0,\ldots,a_{v-1}$ is nonzero, report that the division is not exact.
 Otherwise replace $(a,b)$ by $(A,B)=(a/x^v,b/x^v)$.
 \item Put
 \[
 n\leftarrow\len(A)-\len(B)+1,
 \quad
 \ell\leftarrow\lceil n/2\rceil,
 \quad
 h\leftarrow\lfloor n/2\rfloor.
 \]
 \item Compute $q_L$ from
 \[
 Bq_L\equiv A\pmod{x^\ell}
 \]
 by truncated series division.
 \item Reverse only the highest $h$ coefficients of $A$ and the highest
 $\min(d,h)$ coefficients of $B$.  Use a length-$h$ truncated series quotient
 and reverse its result to obtain the high quotient block $q_H$.
 \item Assemble
 \[
 q\leftarrow q_L+x^\ell q_H.
 \]
 \item Compute the $d-1$ coefficients of $Bq$ in the window
 \eqref{eq:bidirectional-central-window} by one middle-product operation.  If
 they do not equal the corresponding coefficients of $A$, report that the
 division is not exact.
 \item Return $q$.
\end{enumerate}
\end{algorithm}

\begin{theorem}\label{thm:bidirectional-exact-correctness}
Algorithm~\ref{alg:bidirectional-exact} returns $q$ if and only if $a=bq$ for
some $q\in\ZZ[x]$.
\end{theorem}

\begin{proof}
If $a=bq$, the valuation test is necessary and
\eqref{eq:valuation-exact-reduction} reduces the problem to $A=Bq$ with
$B(0)\ne0$.  Lemmas~\ref{lem:bidirectional-low} and
\ref{lem:bidirectional-high} then recover the true low and high blocks of
$q$, so the assembled candidate is the true quotient and the central window
necessarily agrees.

Conversely, suppose both half quotients are integral and the central window
passes.  The low series quotient proves equality of $A$ and $Bq$ in degrees
$0$ through $\ell-1$.  The reversed high quotient proves equality in degrees
$d+\ell-1$ through $d+n-2$.  The middle-product check proves equality in the
remaining degrees $\ell$ through $\ell+d-2$.  These three disjoint intervals
cover every possible coefficient of the product $Bq$, whose formal length is
$d+n-1=\len(A)$.  Hence $A=Bq$, and
\eqref{eq:valuation-exact-reduction} gives $a=bq$.
\end{proof}

\begin{remark}
The middle check is essential.  It is possible for the low and high triangular
systems to have integral solutions which assemble to a candidate quotient,
while one or more central product coefficients disagree with the dividend.
Changing a coefficient of the dividend strictly inside
\eqref{eq:bidirectional-central-window} leaves both half computations
unchanged.  The middle product is therefore not merely a consistency check on
roundoff or implementation details; it verifies equations which neither half
algorithm has examined.
\end{remark}

\subsection{Complexity of bidirectional exact division}

Let $D(k)$ denote the cost of a length-$k$ truncated series quotient and let
$W(d,n)$ denote the cost of computing the central $d-1$ coefficient window of
a product of lengths $d$ and $n$.  Then the bidirectional algorithm has
coefficient-operation complexity
\begin{equation}\label{eq:bidirectional-complexity}
 D(\lceil n/2\rceil)
 +D(\lfloor n/2\rfloor)
 +W(d,n)
 +O(n+d_0),
\end{equation}
where $d_0$ is the number of coefficients inspected while finding and checking
the initial valuation.  Apart from that scan, the high-half construction
copies and reverses only $O(n)$ coefficients, even when the divisor is much
longer than the quotient.

A full product always gives the required verification window, so
\[
 W(d,n)=O(\mathsf M(d,n)),
\]
where $\mathsf M(d,n)$ denotes the cost of multiplying operands of those two
lengths.  Direct middle-product algorithms can of course have a smaller
constant.

\begin{theorem}\label{thm:bidirectional-balanced-complexity}
Suppose $d=\Theta(n)=\Theta(N)$, and suppose series division and the required
product window have coefficient-operation complexity $\Theta(N^\alpha)$ for
some $\alpha>1$.  Then bidirectional exact division has complexity
\[
 \Theta(N^\alpha).
\]
If more precisely
\[
 D(N)\sim K N^\alpha,
 \qquad
 W(N,N)\sim C N^\alpha,
\]
then the leading smooth-model cost is
\begin{equation}\label{eq:bidirectional-leading-constant}
 \left(2^{1-\alpha}K+C\right)N^\alpha.
\end{equation}
\end{theorem}

\begin{proof}
The two half quotients have sizes $N/2+O(1)$, so together they contribute
\[
 2K(N/2)^\alpha+o(N^\alpha)
 =2^{1-\alpha}KN^\alpha+o(N^\alpha).
\]
The verification contributes $CN^\alpha+o(N^\alpha)$, while coefficient
assembly and reversals are linear.  Adding the terms proves both statements.
\end{proof}

For bit complexity, let $\tau$ bound the coefficients of $A$ and $B$, and let
$\rho$ bound the quotient coefficients.  The central product coefficients
have
\[
 O(\tau+\rho+\log\min(d,n))
\]
bits.  Hence, in a balanced regime where both series quotient and middle
window at scale $N$ cost
\[
 O\!\left(N^\alpha\MZ(\tau+\rho+\log N)\right),
\]
the complete bidirectional exact division has the same bit complexity.  The
worst-case substitution $\rho=O(N(\tau+1))$ is again available, giving
\[
 O\!\left(N^\alpha\MZ(N(\tau+1)+\log N)\right).
\]
For a highly unbalanced divisor and quotient, the two-parameter expression
\eqref{eq:bidirectional-complexity} is more informative: the central
verification window may dominate the two half quotients.

\subsection{Complexity summary}

Let $n$ denote quotient length and $d$ divisor length.  The main conclusions of
this section are
\[
\begin{array}{c|c}
\text{operation} & \text{coefficient-operation complexity}\\
\hline
\text{balanced Mulders short quotient} &
 S(n)=S(n/2)+F(n/2)+P(n/2)+O(n)\\
\text{Mulders quotient when }n\leq d & S(n)+O(n)\\
\text{bidirectional exact quotient} &
 D(\lceil n/2\rceil)+D(\lfloor n/2\rfloor)+W(d,n)+O(n+d_0).
\end{array}
\]
In the polynomial-exponent regimes used by the companion implementation, all
of these remain at the same asymptotic exponent as multiplication.  Their
purpose is to avoid arithmetic which is irrelevant to the requested quotient
or to exploit the additional information that the remainder is known to be
zero.



\section{Series inversion and Newton division}\label{sec:newton-division}

The division algorithms of the preceding sections obtain quotient coefficients
from triangular systems directly.  A different approach is to invert the
divisor as a formal power series and replace division by multiplication.  This
is particularly effective when multiplication and middle products are
subquadratic, and it also permits a preinverse to be reused for many dividends.

Throughout this section, for a power series $F$, write
\[
 \operatorname{low}_n(F)=F\bmod x^n,
\]
and, for $s,n\geq0$, write
\[
 \operatorname{mid}_{s,n}(F)
   =\sum_{i=0}^{n-1}[x^{s+i}]F\,x^i.
\]
Thus $\operatorname{mid}_{s,n}$ extracts a window and shifts it down to degree
zero.  Let $L(n)$ denote the cost of a low product of precision $n$, and let
$P(n)$ denote the cost of the balanced middle product
\[
 (2n-1)\times n\longrightarrow n.
\]
The arguments below remain valid with more precise two-parameter costs when the
operands are unbalanced.

The series reciprocal is available over $\ZZ[[x]]$ precisely when the constant
coefficient is a unit.  Hence in this section
\[
 B(x)=b_0+b_1x+b_2x^2+\cdots,
 \qquad b_0\in\{1,-1\}.
\]
For ordinary polynomial division this condition arises after reversal when the
leading coefficient of the divisor is $\pm1$.

Newton iteration for reciprocals is classical.  Its use together with short and
middle products is discussed by Brent and Zimmermann
\cite{BrentZimmermann2010}, and the middle-product formulation for power
series is developed by Hanrot, Quercia and Zimmermann
\cite{HanrotQuerciaZimmermann2004}.  The device of computing only a half
precision reciprocal and incorporating the dividend in the final correction is
due to Karp and Markstein \cite{KarpMarkstein1997}.

\subsection{Classical series inversion}

Let
\[
 G(x)=\sum_{i\geq0}g_i x^i=B(x)^{-1}.
\]
From $BG=1$ we obtain
\[
 b_0g_0=1,
\]
and, for $k\geq1$,
\begin{equation}\label{eq:inverse-coefficient-recurrence}
 b_0g_k+\sum_{j=1}^{k}b_jg_{k-j}=0,
\end{equation}
where coefficients $b_j$ beyond the length of $B$ are interpreted as zero.
Since $b_0^{-1}=b_0$, this gives a triangular recurrence.

\begin{algorithm}[Classical truncated reciprocal]\label{alg:inverse-classical}
Input $B\in\ZZ[x]$ with $b_0=\pm1$ and a precision $n$.
\begin{enumerate}
 \item Set $g_0\leftarrow b_0$.
 \item For $k=1,\ldots,n-1$, set
 \[
  g_k\leftarrow-b_0\sum_{j=1}^{k}b_jg_{k-j}.
 \]
 \item Return $G_n=\sum_{k=0}^{n-1}g_kx^k$.
\end{enumerate}
\end{algorithm}

\begin{theorem}\label{thm:inverse-classical-correctness}
Algorithm~\ref{alg:inverse-classical} returns the unique polynomial $G_n$ of
length at most $n$ satisfying
\[
 BG_n\equiv1\pmod{x^n}.
\]
If $B$ has at least $n$ coefficients, it uses exactly
\[
 \frac{n(n-1)}2
\]
coefficient multiplications in the recurrence, together with $O(n^2)$
coefficient additions.  Thus its coefficient-operation complexity is
$\Theta(n^2)$.
\end{theorem}

\begin{proof}
The equation for coefficient $k$ of $BG_n$ is exactly
\eqref{eq:inverse-coefficient-recurrence}.  At stage $k$, all coefficients
$g_0,\ldots,g_{k-1}$ are already known, and multiplication by $b_0=\pm1$ is
invertible in $\ZZ$.  The recurrence therefore determines the unique $g_k$
which makes coefficient $k$ vanish.  Induction on $k$ proves
$BG_n\equiv1\pmod{x^n}$ and uniqueness.

For $k\geq1$, the inner sum has $k$ products in the full-length case.  Hence
the number of coefficient products is
\[
 \sum_{k=1}^{n-1}k=\frac{n(n-1)}2.
\]
The additions have the same quadratic order.
\end{proof}

The inverse coefficients can grow rapidly even when the coefficients of $B$
are small.  The following elementary estimate is sufficient for subsequent bit
bounds.

\begin{lemma}\label{lem:inverse-coefficient-growth}
Suppose $\bits(b_j)\leq\tau$ for all $j$, with $b_0=\pm1$.  Then the
coefficients of $B^{-1}\bmod x^n$ satisfy
\[
 \bits(g_k)=O\bigl(k(\tau+1)\bigr),
 \qquad 0\leq k<n.
\]
More explicitly, if $S=\max(1,\max_{j\geq1}|b_j|)$, then
\[
 |g_k|\leq(1+S)^k.
\]
\end{lemma}

\begin{proof}
Let
\[
 H_k=\sum_{i=0}^{k}|g_i|.
\]
Since $|g_0|=1$, and the recurrence gives
\[
 |g_k|\leq S H_{k-1},
\]
we have
\[
 H_k\leq(1+S)H_{k-1}.
\]
Thus $H_k\leq(1+S)^k$, and hence the same bound holds for $|g_k|$.  Since
$S<2^\tau$ when $\tau>0$, the bit bound follows.
\end{proof}

Consequently, a direct bit-complexity bound for classical inversion is
\[
 O\!\left(
 n^2\MZ\bigl(n(\tau+1)\bigr)
 \right),
\]
which deliberately uses the largest coefficient size throughout.  A more
refined summation over $k$ can improve the constant but not the quadratic
coefficient-operation exponent.

\subsection{Newton reciprocal with middle and low products}

Suppose a reciprocal $G$ is already known modulo $x^m$:
\[
 BG\equiv1\pmod{x^m}.
\]
Let $n=m+h$ with $0\leq h\leq m$.  There is a unique polynomial $E$ of length
at most $h$ such that
\begin{equation}\label{eq:newton-error-block}
 BG\equiv1+x^mE\pmod{x^n}.
\end{equation}
Only the coefficient window $m,\ldots,n-1$ of $BG$ is needed to obtain $E$;
this is a middle product.  Define
\begin{equation}\label{eq:newton-reciprocal-correction}
 C=\operatorname{low}_h(GE),
 \qquad
 G'=G-x^mC.
\end{equation}
This is the truncated form of the familiar Newton identity
\[
 G' = G(2-BG).
\]
The unfactored form \eqref{eq:newton-reciprocal-correction} is preferable
computationally because it exposes that only the high error block and a low
correction product are required.

\begin{algorithm}[Newton reciprocal]\label{alg:newton-reciprocal}
Input $B\in\ZZ[x]$ with $b_0=\pm1$ and precision $n$.
\begin{enumerate}
 \item If $n$ is below a base-case threshold, use
 Algorithm~\ref{alg:inverse-classical}.
 \item Put
 \[
   m=\left\lceil\frac n2\right\rceil,
   \qquad h=n-m.
 \]
 \item Recursively compute $G=B^{-1}\bmod x^m$.
 \item Compute
 \[
   E=\operatorname{mid}_{m,h}(BG).
 \]
 \item Compute $C=\operatorname{low}_h(GE)$.
 \item Return $G-x^mC$.
\end{enumerate}
\end{algorithm}

\begin{theorem}\label{thm:newton-reciprocal-correctness}
Algorithm~\ref{alg:newton-reciprocal} returns
$B^{-1}\bmod x^n$.
\end{theorem}

\begin{proof}
By induction, $BG\equiv1\pmod{x^m}$, so
\eqref{eq:newton-error-block} holds for the $E$ extracted in step 4.  Since
$h\leq m$, reducing modulo $x^h$ gives
\[
 BG\equiv1\pmod{x^h}.
\]
Hence
\[
 B C
 \equiv BGE
 \equiv E
 \pmod{x^h}.
\]
Using \eqref{eq:newton-error-block},
\[
 B(G-x^mC)
 \equiv 1+x^mE-x^mBC
 \equiv1
 \pmod{x^{m+h}}.
\]
Since $m+h=n$, the returned polynomial is the reciprocal modulo $x^n$.
\end{proof}

For $n=2m$, step 4 is exactly a balanced middle product of shape
\[
 (2m-1)\times m\longrightarrow m,
\]
after the irrelevant constant coefficient of $B$ is skipped.  The odd case
differs by only one coefficient and has the same asymptotic cost.

\begin{theorem}\label{thm:newton-reciprocal-complexity}
Let $I(n)$ denote the coefficient-operation cost of Newton reciprocal.  Then
\begin{equation}\label{eq:newton-inverse-recurrence}
 I(n)
 =I(\lceil n/2\rceil)
  +P(\lfloor n/2\rfloor)
  +L(\lfloor n/2\rfloor)
  +O(n).
\end{equation}
In particular, if
\[
 L(n)=O(\mathsf M(n)),\qquad P(n)=O(\mathsf M(n))
\]
for a regular superlinear multiplication cost $\mathsf M$, then
\[
 I(n)=O(\mathsf M(n)).
\]
If more precisely
\[
 L(n)\sim Bn^\alpha,
 \qquad
 P(n)\sim Cn^\alpha,
 \qquad \alpha>1,
\]
then
\begin{equation}\label{eq:newton-inverse-leading-constant}
 I(n)\sim\frac{B+C}{2^\alpha-1}n^\alpha.
\end{equation}
\end{theorem}

\begin{proof}
The recurrence is the direct cost of steps 3--5; all coefficient assembly is
linear.  Under the stated regularity assumption, summing the costs over the
precision sequence $n,n/2,n/4,\ldots$ gives $O(\mathsf M(n))$.

For the smooth power-law model, write $I(n)\sim Kn^\alpha$.  Substitution into
\eqref{eq:newton-inverse-recurrence} gives
\[
 K=2^{-\alpha}K+2^{-\alpha}(B+C),
\]
whence
\[
 K=\frac{B+C}{2^\alpha-1}.
\]
\end{proof}

By Lemma~\ref{lem:inverse-coefficient-growth}, all coefficients at precision
$n$ have $O(n(\tau+1))$ bits.  Therefore, in a model where the length-$n$ low
and middle products on $w$-bit coefficients cost
\[
 O\bigl(n^\alpha\MZ(w)\bigr),
\]
Newton reciprocal has bit complexity
\begin{equation}\label{eq:newton-inverse-bit-complexity}
 O\!\left(
 n^\alpha\MZ\bigl(n(\tau+1)\bigr)
 \right).
\end{equation}
The same statement may of course be expressed directly in terms of the chosen
bit-complexity functions for low and middle products.

\subsection{Division by a reusable preinverse}

Suppose $G$ is already known with
\[
 BG\equiv1\pmod{x^n}.
\]
Then for any $A$,
\begin{equation}\label{eq:series-quotient-preinverse}
 \frac AB\bmod x^n
 =\operatorname{low}_n(AG).
\end{equation}
Thus, once a reciprocal has been precomputed, every further quotient requires
only one low product.

\begin{algorithm}[Series quotient from a preinverse]\label{alg:series-preinverse}
Input $A,B,G\in\ZZ[x]$ and $n$ with $BG\equiv1\pmod{x^n}$.
Return
\[
 Q=\operatorname{low}_n(AG).
\]
\end{algorithm}

\begin{theorem}\label{thm:series-preinverse-correctness}
Algorithm~\ref{alg:series-preinverse} returns the unique $Q$ of length at most
$n$ satisfying
\[
 BQ\equiv A\pmod{x^n}.
\]
After the preinverse has been computed, its cost is $L(n)+O(n)$.
\end{theorem}

\begin{proof}
Multiplying \eqref{eq:series-quotient-preinverse} by $B$ gives
\[
 BQ\equiv BAG\equiv A\pmod{x^n}.
\]
Uniqueness follows because the constant coefficient of $B$ is a unit.  The
only non-linear operation is the low product.
\end{proof}

For ordinary polynomial division, let $A$ have length $a$, let $B$ have length
$d$, and let
\[
 n=a-d+1
\]
be the quotient length.  Assume the leading coefficient of $B$ is $\pm1$.
As in Section~\ref{sec:classical-division}, reverse only the highest $n$
coefficients that can influence the quotient.  More precisely, put
\[
 r=\min(d,n),
\]
let $B^\star$ be the reversal of the highest $r$ coefficients of $B$, padded
conceptually by zeros to precision $n$, and precompute
\[
 G=(B^\star)^{-1}\bmod x^n.
\]
For each dividend, reverse its highest $n$ coefficients to obtain $A^\star$,
compute
\[
 Q^\star=\operatorname{low}_n(A^\star G),
\]
and reverse once more to obtain the ordinary quotient.

\begin{corollary}\label{cor:ordinary-preinverse}
A preinverse of a fixed unit-leading divisor to quotient precision $n$ costs
$I(n)+O(n)$.  Each subsequent length-$n$ quotient costs
\[
 L(n)+O(n).
\]
In the smooth model of Theorem~\ref{thm:newton-reciprocal-complexity}, the
precomputation has leading constant
\[
 \frac{B+C}{2^\alpha-1},
\]
while each quotient has leading constant $B$.
\end{corollary}

This is preferable when the same divisor is reused.  For a single division,
computing a reciprocal all the way to precision $n$ performs more work than is
necessary.  The Karp--Markstein construction avoids the final full-precision
reciprocal step.

\subsection{Karp--Markstein division}\label{subsec:karp-markstein}

Let
\[
 Q=\frac AB\bmod x^n,
 \qquad
 m=\left\lceil\frac n2\right\rceil,
 \qquad
 h=n-m.
\]
Compute only a half-precision reciprocal
\[
 G=B^{-1}\bmod x^m.
\]
The low half of the quotient is then
\begin{equation}\label{eq:km-q0}
 Q_0=\operatorname{low}_m(AG).
\end{equation}
Since $BQ_0\equiv A\pmod{x^m}$, the residual is divisible by $x^m$.  Define
its next $h$ coefficients by
\begin{equation}\label{eq:km-error}
 E=\operatorname{mid}_{m,h}(A-BQ_0).
\end{equation}
Then the remaining quotient block is
\begin{equation}\label{eq:km-q1}
 Q_1=\operatorname{low}_h(EG),
\end{equation}
and
\[
 Q=Q_0+x^mQ_1.
\]
The middle product in \eqref{eq:km-error} is the key economy: the coefficients
below $m$ are already known to cancel, and coefficients at degree $n$ and
above are irrelevant.

\begin{algorithm}[Karp--Markstein series quotient]\label{alg:karp-markstein}
Input $A,B\in\ZZ[x]$, with $b_0=\pm1$, and precision $n$.
\begin{enumerate}
 \item For small $n$, use classical series division.
 \item Put $m=\lceil n/2\rceil$ and $h=n-m$.
 \item Compute $G=B^{-1}\bmod x^m$ by Newton reciprocal.
 \item Compute $Q_0=\operatorname{low}_m(AG)$.
 \item Compute
 \[
  E=\operatorname{mid}_{m,h}(A-BQ_0).
 \]
 In practice the $A$ window is copied and the corresponding middle product
 of $B$ and $Q_0$ is subtracted.
 \item Compute $Q_1=\operatorname{low}_h(EG)$.
 \item Return $Q_0+x^mQ_1$.
\end{enumerate}
\end{algorithm}

\begin{theorem}\label{thm:karp-markstein-correctness}
Algorithm~\ref{alg:karp-markstein} returns the unique polynomial $Q$ of length
at most $n$ satisfying
\[
 BQ\equiv A\pmod{x^n}.
\]
\end{theorem}

\begin{proof}
By \eqref{eq:km-q0} and $BG\equiv1\pmod{x^m}$,
\[
 BQ_0\equiv A\pmod{x^m}.
\]
Thus, modulo $x^n$,
\[
 A-BQ_0\equiv x^mE.
\]
Because $h\leq m$, the half-precision reciprocal also satisfies
$BG\equiv1\pmod{x^h}$.  Therefore
\[
 BQ_1
 \equiv BEG
 \equiv E
 \pmod{x^h}.
\]
It follows that
\[
 B(Q_0+x^mQ_1)
 \equiv BQ_0+x^mE
 \equiv A
 \pmod{x^n}.
\]
Uniqueness again follows from the fact that $b_0$ is a unit.
\end{proof}

\begin{theorem}\label{thm:karp-markstein-complexity}
Let $K(n)$ denote the coefficient-operation cost of one Karp--Markstein series
quotient.  Then
\begin{equation}\label{eq:karp-markstein-cost}
 K(n)
 =I(\lceil n/2\rceil)
  +L(\lceil n/2\rceil)
  +P(\lfloor n/2\rfloor)
  +L(\lfloor n/2\rfloor)
  +O(n).
\end{equation}
Hence $K(n)=O(\mathsf M(n))$ whenever inverse, low product and middle product
are $O(\mathsf M(n))$.

If
\[
 L(n)\sim Bn^\alpha,
 \qquad P(n)\sim Cn^\alpha,
 \qquad \alpha>1,
\]
then
\begin{equation}\label{eq:karp-markstein-leading-constant}
 K(n)
 \sim
 \frac{1}{2^\alpha}
 \left(
   \frac{B+C}{2^\alpha-1}+2B+C
 \right)n^\alpha.
\end{equation}
\end{theorem}

\begin{proof}
The four non-linear operations in Algorithm~\ref{alg:karp-markstein} are,
in order, a half-precision inverse, a low product, a middle product, and a
second low product.  This gives \eqref{eq:karp-markstein-cost}.  The asymptotic
statements follow from Theorem~\ref{thm:newton-reciprocal-complexity} and
substitution of the smooth models at size $n/2+O(1)$.
\end{proof}

For comparison, computing a full precision reciprocal and then multiplying
once has leading smooth-model cost
\begin{equation}\label{eq:full-preinverse-oneoff-cost}
 \left(
  \frac{B+C}{2^\alpha-1}+B
 \right)n^\alpha.
\end{equation}
Karp--Markstein division avoids the last Newton reciprocal step and folds the
dividend into that precision-doubling stage.  The saving is therefore a
constant-factor one, while the asymptotic exponent is unchanged.

As a check, in a purely classical smooth model one may take
\[
 B=\frac12,
 \qquad C=1,
 \qquad \alpha=2,
\]
corresponding to a triangular low product and a square middle product.  Then
Newton reciprocal has leading cost $\tfrac12n^2$, a full reciprocal followed
by one low product has leading cost $n^2$, and the Karp--Markstein construction
has leading cost $\tfrac58n^2$.  Classical triangular division itself has
leading cost $\tfrac12n^2$, explaining why Newton methods are not expected to
win in the quadratic range.

\subsection{Ordinary Newton quotient by reversal}

The Karp--Markstein series quotient becomes an ordinary polynomial quotient by
exactly the reversal transformation of Section~\ref{sec:classical-division}.
Let $A$ and $B$ have lengths $a\geq d$, and let $n=a-d+1$.  Put
\[
 r=\min(d,n).
\]
Reverse the highest $n$ coefficients of $A$, and only the highest $r$
coefficients of $B$.  The latter are all that can enter the first $n$ equations
of the reversed power-series division.  Apply
Algorithm~\ref{alg:karp-markstein} to precision $n$, then reverse the result.

\begin{theorem}\label{thm:ordinary-newton-quotient}
If the leading coefficient of $B$ is $\pm1$, the above procedure returns the
ordinary polynomial quotient of $A$ by $B$.  Its coefficient-operation cost is
\[
 K(n)+O(n),
\]
where $K$ is given by \eqref{eq:karp-markstein-cost}.  The bound is independent
of the coefficients of $B$ below its highest $n$ positions except through the
case $d<n$, where the whole divisor is already shorter than the quotient.
\end{theorem}

\begin{proof}
The reversal identity from Section~\ref{sec:classical-division} converts the
highest-to-lowest quotient equations into the first $n$ triangular
power-series equations.  Those equations involve only the first $n$
coefficients of the reversed divisor, namely the highest $\min(d,n)$
coefficients of $B$.  Theorem~\ref{thm:karp-markstein-correctness} solves this
series quotient, and reversing restores the ordinary coefficient order.
The extra work consists only of the two length-$n$ reversals and the divisor
tail reversal.
\end{proof}

\subsection{Newton quotient and remainder recovery}

Once the ordinary quotient $Q$ is known, a full product $BQ$ is unnecessary if
the remainder alone is required.  If $d=\len(B)$, then
\[
 R=A-BQ,
 \qquad \deg R<d-1.
\]
Consequently
\begin{equation}\label{eq:newton-remainder-low}
 R
 =\operatorname{low}_{d-1}(A)
  -\operatorname{low}_{d-1}(BQ).
\end{equation}
All higher coefficients of $BQ$ are known, by correctness of the quotient, to
cancel the corresponding coefficients of $A$.

\begin{algorithm}[Newton quotient and remainder]\label{alg:newton-divrem}
Input $A,B\in\ZZ[x]$ with unit leading coefficient and $\len(A)\geq\len(B)$.
\begin{enumerate}
 \item Compute the quotient $Q$ by the ordinary Karp--Markstein/Newton
 quotient algorithm.
 \item Put $d=\len(B)$ and compute
 \[
   T=\operatorname{low}_{d-1}(BQ).
 \]
 \item Return
 \[
   \left(Q,\operatorname{low}_{d-1}(A)-T\right).
 \]
\end{enumerate}
\end{algorithm}

\begin{theorem}\label{thm:newton-divrem-correctness}
Algorithm~\ref{alg:newton-divrem} returns the ordinary quotient and remainder
of $A$ by $B$.
\end{theorem}

\begin{proof}
The quotient step gives the ordinary quotient $Q$, so $A-BQ$ has degree less
than $d-1$.  Therefore it equals its own truncation modulo $x^{d-1}$, and
\eqref{eq:newton-remainder-low} is exactly $A-BQ$.
\end{proof}

Let $L_{d,n}(d-1)$ denote the cost of the requested low window of the product
of a length-$d$ divisor and a length-$n$ quotient.  Then the complete division
cost is
\begin{equation}\label{eq:newton-divrem-complexity}
 K(n)+L_{d,n}(d-1)+O(d+n).
\end{equation}
For balanced $d,n=\Theta(N)$, if low products and the Newton quotient are
$O(\mathsf M(N))$, quotient-and-remainder division is also
$O(\mathsf M(N))$.  The low product in the remainder step is preferable to a
full multiplication whenever only the remainder is needed.

\subsection{Bit complexity and coefficient growth}

Let $\bits(b_i)\leq\tau$ and $\bits(a_i)\leq\sigma$.  By
Lemma~\ref{lem:inverse-coefficient-growth}, an inverse to precision $n$ has
coefficient size
\[
 O\bigl(n(\tau+1)\bigr).
\]
Since a series quotient is the low product of the dividend with the inverse,
its coefficients have size
\begin{equation}\label{eq:newton-quotient-coefficient-size}
 O\bigl(\sigma+n(\tau+1)+\log n\bigr).
\end{equation}
The same bound applies to the intermediate quotient blocks and residual
corrections in Karp--Markstein division.

Suppose length-$k$ low and middle products on coefficients of at most $w$ bits
cost
\[
 O\bigl(k^\alpha\MZ(w)\bigr)
\]
for some $\alpha>1$.  Then Newton reciprocal has the bit bound
\[
 O\!\left(
 n^\alpha\MZ\bigl(n(\tau+1)\bigr)
 \right),
\]
and Karp--Markstein series division has
\begin{equation}\label{eq:km-bit-complexity}
 O\!\left(
 n^\alpha
 \MZ\bigl(\sigma+n(\tau+1)+\log n\bigr)
 \right).
\end{equation}
For ordinary polynomial division the same estimate applies after reversal.
The remainder step multiplies divisor coefficients of size $\tau$ by quotient
coefficients of the size in \eqref{eq:newton-quotient-coefficient-size}; its
bit complexity is bounded by the corresponding low-product cost at precision
$d-1$.

These are worst-case bounds.  In many applications the quotient and reciprocal
coefficients are much smaller, and it is then more informative to replace
\eqref{eq:newton-quotient-coefficient-size} by an observed or separately proved
coefficient bound $\rho$.

\subsection{Complexity summary}

Let $n$ denote the requested reciprocal or quotient precision.  The algorithms
of this section have the following coefficient-operation costs:
\[
\begin{array}{c|c}
\text{operation} & \text{cost}\\
\hline
\text{classical reciprocal} & \Theta(n^2)\\
\text{Newton reciprocal} &
 I(n)=I(n/2)+P(n/2)+L(n/2)+O(n)\\
\text{quotient from preinverse} & L(n)+O(n)\\
\text{one-off Karp--Markstein quotient} &
 I(n/2)+2L(n/2)+P(n/2)+O(n)\\
\text{ordinary Newton quotient} & K(n)+O(n)\\
\text{Newton quotient and remainder} &
 K(n)+L_{d,n}(d-1)+O(d+n).
\end{array}
\]
Whenever low and middle products have the same asymptotic order as full
multiplication, all Newton-based rows except the classical base case are
$O(\mathsf M(n))$ in the balanced regime.  Their distinctions are therefore
primarily in leading constants and in whether the reciprocal can be reused.



\section{Pseudo-division over $\ZZ[x]$}\label{sec:pseudodivision}

Ordinary polynomial division in $\ZZ[x]$ can fail because the leading
coefficient of the divisor need not be a unit.  Pseudo-division clears the
necessary denominators in advance and remains entirely in $\ZZ[x]$.
It is a basic ingredient of polynomial remainder sequences and of several
integer-polynomial gcd algorithms; classical discussions include
Brown--Traub \cite{BrownTraub1971} and von zur Gathen--Gerhard
\cite{vonZurGathenGerhard2013}.

Let
\[
 A(x)=\sum_{i=0}^{m-1}a_i x^i,
 \qquad
 B(x)=\sum_{j=0}^{n-1}b_j x^j\ne0,
\]
with $\beta=b_{n-1}=\operatorname{lc}(B)$.  Put
\[
 d=\max(0,m-n+1).
\]
A \emph{pseudo-quotient} $Q$ and \emph{pseudo-remainder} $R$ are defined by
\begin{equation}\label{eq:pseudodiv-identity}
 \beta^d A=BQ+R,
 \qquad
 \deg R<\deg B.
\end{equation}
If $m<n$, this says simply $Q=0$ and $R=A$.  If $m\ge n$, then over
$\mathbb Q[x]$ the pair $(Q,R)$ is $\beta^d$ times the ordinary quotient and
remainder, and is therefore unique.  The point of pseudo-division is that
$Q,R\in\ZZ[x]$ can be constructed without ever forming a rational
coefficient.

When bit complexity is discussed, suppose
\[
 \bits(a_i)\le\sigma,
 \qquad
 \bits(b_j)\le\tau.
\]
We use $d$ for the quotient length and $n$ for the divisor length.  Thus the
length of the dividend is at most $n+d-1$.

\subsection{Classical pseudo-division}

The classical algorithm eliminates one leading term at a time.  It is useful
to state it in a form which also handles an accidental degree drop of more
than one.

\begin{algorithm}[Classical pseudo-division]\label{alg:pseudodiv-classical}
Input $A,B\in\ZZ[x]$ with $B\ne0$.
\begin{enumerate}
 \item Put $\beta\leftarrow\operatorname{lc}(B)$,
 $d\leftarrow\max(0,\len(A)-\len(B)+1)$,
 $Q\leftarrow0$, $R\leftarrow A$, and $t\leftarrow0$.
 \item While $R\ne0$ and $\deg R\ge\deg B$:
 \begin{enumerate}
  \item put $k\leftarrow\deg R-\deg B$ and
  $c\leftarrow\operatorname{lc}(R)$;
  \item set
  \[
   Q\leftarrow\beta Q+c x^k,
   \qquad
   R\leftarrow\beta R-cx^kB;
  \]
  \item set $t\leftarrow t+1$.
 \end{enumerate}
 \item If $t<d$, multiply both $Q$ and $R$ by $\beta^{d-t}$.
 \item Return $(Q,R)$.
\end{enumerate}
\end{algorithm}

The final scaling in Step 3 is needed only when one elimination cancels more
than the current leading coefficient.  It ensures that the conventional
power $\beta^d$ in \eqref{eq:pseudodiv-identity} is returned regardless of
such degree drops.

\begin{theorem}\label{thm:pseudodiv-classical-correctness}
Algorithm~\ref{alg:pseudodiv-classical} returns the unique pair $Q,R\in\ZZ[x]$
satisfying \eqref{eq:pseudodiv-identity}.
\end{theorem}

\begin{proof}
After $t$ eliminations maintain the invariant
\begin{equation}\label{eq:pseudodiv-loop-invariant}
 \beta^t A=BQ+R.
\end{equation}
It is true initially.  If $k=\deg R-\deg B$ and
$c=\operatorname{lc}(R)$, then the update gives
\[
 B(\beta Q+cx^k)+(\beta R-cx^kB)
 =\beta(BQ+R)=\beta^{t+1}A.
\]
Moreover, the leading coefficient of $\beta R$ is $\beta c$, which is
exactly cancelled by the leading coefficient of $cx^kB$.  Hence the degree
of $R$ strictly decreases at every iteration.  There can therefore be at
most $d$ iterations, and on termination $\deg R<\deg B$.

If $t<d$, multiplying both sides of \eqref{eq:pseudodiv-loop-invariant} by
$\beta^{d-t}$ gives \eqref{eq:pseudodiv-identity}.  Finally, uniqueness
follows from ordinary quotient-remainder uniqueness over $\mathbb Q[x]$.
\end{proof}

\begin{theorem}\label{thm:pseudodiv-classical-complexity}
Classical pseudo-division uses
\[
 O\bigl(d(n+d)\bigr)
\]
coefficient operations in the worst case.  This bound is attained in order
for dense inputs with no exceptional degree drops.  In the balanced regime
$n,d=\Theta(N)$ the cost is therefore $\Theta(N^2)$.
\end{theorem}

\begin{proof}
There are at most $d$ elimination steps.  At any step the current remainder
has length at most $n+d-1$, the quotient has length at most $d$, and the
shifted multiple $cx^kB$ has $n$ coefficients.  Scaling the current quotient
and remainder and subtracting the shifted divisor therefore costs
$O(n+d)$ coefficient operations per step.  This proves the upper bound.

For dense inputs whose remainder degree drops by exactly one at every step,
there are $d$ steps, each of which updates $\Theta(n)$ remainder
coefficients, while the accumulated quotient scalings contribute
$\Theta(d^2)$ operations.  Thus the order $d(n+d)$ is attained.
\end{proof}

Pseudo-division is fraction-free, but its coefficients need not remain small.
The following simple bound will be used for both the classical and fast
algorithms.

\begin{lemma}\label{lem:pseudodiv-coefficient-growth}
Let
\[
 H_A=\max(1,\max_i|a_i|),
 \qquad
 H_B=\max(1,\max_j|b_j|).
\]
During Algorithm~\ref{alg:pseudodiv-classical}, after $t$ eliminations,
\[
 \max\bigl(H(Q),H(R)\bigr)
 \le H_A(2H_B)^t,
\]
up to replacing the zero initial quotient height by $1$.  After the final
normalizing power of $\beta$, every output coefficient has
\begin{equation}\label{eq:pseudodiv-output-bits}
 O\bigl(\sigma+d(\tau+1)\bigr)
\end{equation}
bits.
\end{lemma}

\begin{proof}
Assume the displayed height bound at stage $t$.  Since
$|\beta|\le H_B$ and $|c|\le H(R)$,
\[
 H(\beta R-cx^kB)
 \le 2H_BH(R)
 \le H_A(2H_B)^{t+1}.
\]
Similarly,
\[
 H(\beta Q+cx^k)
 \le \max(H_BH(Q),H(R))
 \le H_A(2H_B)^{t+1}
\]
after harmlessly enlarging the constant in the first few stages.  Induction
proves the claim during the loop.

If only $t<d$ eliminations occur, the final multiplication by
$\beta^{d-t}$ contributes at most $H_B^{d-t}$.  Thus the final height is at
most $H_A2^dH_B^d$ up to an absolute constant.  Taking binary logarithms
gives \eqref{eq:pseudodiv-output-bits}.
\end{proof}

\begin{corollary}\label{cor:pseudodiv-classical-bit}
Put
\[
 W=O\bigl(\sigma+d(\tau+1)\bigr).
\]
A direct worst-case bit bound for classical pseudo-division is
\[
 O\bigl(d(n+d)\MZ(W)\bigr).
\]
For $n,d=\Theta(N)$ this is
\[
 O\!\left(N^2\MZ\bigl(\sigma+N(\tau+1)\bigr)\right).
\]
\end{corollary}

\begin{proof}
Theorem~\ref{thm:pseudodiv-classical-complexity} gives
$O(d(n+d))$ integer coefficient operations.  By
Lemma~\ref{lem:pseudodiv-coefficient-growth}, all intermediate integers have
$O(W)$ bits.  Charging every multiplication the upper bound $\MZ(W)$ proves
the result; additions are no more expensive asymptotically under the standing
assumption $\MZ(W)=\Omega(W)$.
\end{proof}

The linear growth in bit length for a \emph{single} pseudo-division should not
be confused with repeated pseudo-Euclidean division.  If pseudo-remainders are
fed back into subsequent pseudo-divisions without normalization, their already
large coefficients become the inputs to the next scaling step, and severe
intermediate expression swell can result.  This is one reason for primitive
and subresultant polynomial remainder sequences; see
Brown--Traub \cite{BrownTraub1971}.

\subsection{Scaled power-series reciprocals}

Fast pseudo-division needs a reciprocal without dividing by the nonunit
constant coefficient.  Let
\[
 F(x)=c+f_1x+f_2x^2+\cdots\in\ZZ[x],
 \qquad c\ne0.
\]
For $m\ge1$, define the \emph{scaled reciprocal} $I_m$ by
\begin{equation}\label{eq:scaled-inverse-definition}
 F I_m\equiv c^m\pmod{x^m},
 \qquad
 \deg I_m<m.
\end{equation}
Equivalently, over $\mathbb Q[[x]]$,
\[
 I_m=\operatorname{low}_m\!\left(\frac{c^m}{F}\right).
\]
The scale $c^m$ clears all denominators through precision $m$.  This power is
needed in general: for $F=c+x$, the coefficient of $x^{m-1}$ in
$c^{m-1}/F$ is $(-1)^{m-1}/c$.

\begin{lemma}\label{lem:scaled-inverse-integral}
For every $m\ge1$, the polynomial $I_m$ in
\eqref{eq:scaled-inverse-definition} belongs to $\ZZ[x]$ and is unique.  If
all coefficients of $F$ have at most $\tau$ bits, then
\begin{equation}\label{eq:scaled-inverse-bits}
 \bits([x^k]I_m)=O\bigl(m(\tau+1)\bigr),
 \qquad 0\le k<m.
\end{equation}
\end{lemma}

\begin{proof}
Write $F=c+U$, where $U(0)=0$.  In $\mathbb Q[[x]]$,
\[
 \frac{c^m}{F}
 =c^{m-1}\left(1+\frac{U}{c}\right)^{-1}
 =\sum_{t\ge0}(-1)^t c^{m-1-t}U^t.
\]
The coefficient of $x^k$, for $k<m$, receives contributions only from
$t\le k\le m-1$.  Hence every power $c^{m-1-t}$ has a nonnegative exponent,
so this coefficient is an integer.  This proves existence in $\ZZ[x]$.
Uniqueness follows because $c\ne0$, so $F$ is invertible in
$\mathbb Q[[x]]$.

For the size bound let
$H=\max(1,\max_i|f_i|,|c|)<2^{\tau}$, up to an immaterial endpoint
convention.  For $t\ge1$, the coefficient $[x^k]U^t$ is a sum over the
compositions of $k$ into $t$ positive parts.  There are
$\binom{k-1}{t-1}$ such compositions, and every product has magnitude at most
$H^t$.  Thus the contribution for fixed $t$ has magnitude at most
\[
 \binom{k-1}{t-1}H^{m-1}.
\]
Summing over $t$ gives at most $2^{k-1}H^{m-1}$ for $k>0$, while the constant
coefficient is $c^{m-1}$.  Taking logarithms proves
\eqref{eq:scaled-inverse-bits}.
\end{proof}

The scaled reciprocal admits exactly the same Newton error-squaring mechanism
as an ordinary reciprocal.

\begin{algorithm}[Scaled Newton reciprocal]\label{alg:scaled-newton-inverse}
Input $F\in\ZZ[x]$ with $c=F(0)\ne0$ and a power-of-two precision $p$.
\begin{enumerate}
 \item Set $m\leftarrow1$, $I_1\leftarrow1$, and $C\leftarrow c$.
 \item While $m<p$:
 \begin{enumerate}
  \item compute
  \[
   T\leftarrow\operatorname{low}_{2m}(FI_m);
  \]
  \item compute
  \[
   I_{2m}\leftarrow
   \operatorname{low}_{2m}\bigl(I_m(2C-T)\bigr);
  \]
  \item set $C\leftarrow C^2$, $m\leftarrow2m$, and
  $I_m\leftarrow I_{2m}$.
 \end{enumerate}
 \item Return $I_p$.
\end{enumerate}
\end{algorithm}

Here $C=c^m$ at the beginning of every loop iteration.  The companion
implementation recomputes this scalar power when required rather than storing
it across iterations; this does not change the polynomial arithmetic.

\begin{theorem}\label{thm:scaled-newton-correctness}
Algorithm~\ref{alg:scaled-newton-inverse} returns the scaled reciprocal $I_p$,
so
\[
 FI_p\equiv c^p\pmod{x^p}.
\]
If $L(s)$ denotes the coefficient-operation cost of a low product of precision
$s$, its recurrence is
\begin{equation}\label{eq:scaled-newton-recurrence}
 J(2m)=J(m)+2L(2m)+O(m).
\end{equation}
Consequently, under the usual geometric-growth hypothesis on $L$,
$J(p)=O(L(p))$.  More precisely, if
$L(s)\sim A s^\alpha$ with $\alpha>0$, then
\[
 J(p)\sim\frac{2A}{1-2^{-\alpha}}p^\alpha.
\]
\end{theorem}

\begin{proof}
Suppose
\[
 FI_m=c^m+x^mE.
\]
Then, before truncation,
\[
 F I_m(2c^m-FI_m)
 =(c^m+x^mE)(c^m-x^mE)
 =c^{2m}-x^{2m}E^2.
\]
Thus
\[
 FI_{2m}\equiv c^{2m}\pmod{x^{2m}},
\]
which doubles the precision.  Since $I_1=1$ satisfies
$FI_1\equiv c\pmod{x}$, induction proves correctness.

Each doubling performs two low products of precision $2m$, plus linear
coefficient additions and scalar work.  This gives
\eqref{eq:scaled-newton-recurrence}.  Summing the recurrence over geometrically
increasing precisions proves $O(L(p))$ under geometric growth.  If
$L(s)\sim As^\alpha$, the main terms form
\[
 2Ap^\alpha\sum_{j\ge0}2^{-j\alpha},
\]
which gives the stated constant.
\end{proof}

\begin{corollary}\label{cor:scaled-newton-bit}
Let $\mathcal L(s,w)$ denote the bit complexity of a low product of precision
$s$ whose intermediate coefficients have at most $w$ bits.  Then scaled
Newton inversion to precision $p$ has bit complexity
\[
 O\!\left(
 \sum_{m=1,2,4,\ldots<p}
 \mathcal L\bigl(2m,O(m(\tau+1))\bigr)
 \right),
\]
plus the cost of the scalar powers of $c$.  In particular, if
\[
 \mathcal L(s,w)=O\bigl(s^\alpha\MZ(w)\bigr)
\]
for some $\alpha>0$, then
\[
 J_{\rm bit}(p)
 =O\!\left(p^\alpha\MZ\bigl(p(\tau+1)\bigr)\right).
\]
\end{corollary}

\begin{proof}
Lemma~\ref{lem:scaled-inverse-integral} bounds the coefficients at precision
$m$ by $O(m(\tau+1))$ bits.  Substitute this into the two low products at each
doubling.  Monotonicity of $\MZ$ and the geometric factor $m^\alpha$ make the
sum dominated, up to a constant, by its largest precision.  Maintaining
$c^m$ by squaring adds one integer multiplication at each level; recomputing
it by binary powering changes only the displayed lower-order scalar term.
\end{proof}

\subsection{Fast pseudo-quotient by reversal}

Scaled inversion converts pseudo-division to truncated multiplication.  Assume
$m\ge n$ and put $d=m-n+1$.  Define fixed-length reversals
\[
 A^*=\operatorname{rev}_m(A),
 \qquad
 B^*=\operatorname{rev}_n(B),
 \qquad
 Q^*=\operatorname{rev}_d(Q).
\]
Reversing \eqref{eq:pseudodiv-identity} and observing that the reversed
remainder begins at degree at least $d$ gives
\begin{equation}\label{eq:pseudodiv-reversal-congruence}
 \beta^d A^*\equiv B^*Q^*\pmod{x^d}.
\end{equation}
The constant coefficient of $B^*$ is $\beta$, so its scaled reciprocal is
available.

Choose the least power of two $p\ge d$.  Then
\begin{equation}\label{eq:pseudodiv-p-range}
 d\le p<2d
\end{equation}
unless $d=0$.  Let $I_p$ satisfy
\[
 B^*I_p\equiv\beta^p\pmod{x^p}.
\]

\begin{lemma}\label{lem:fast-pseudoquotient}
Put
\[
 S=\operatorname{low}_d(A^*I_p).
\]
Then every coefficient of $S$ is divisible by $\beta^{p-d}$ and
\begin{equation}\label{eq:fast-pseudoquotient-scale}
 S=\beta^{p-d}Q^*.
\end{equation}
\end{lemma}

\begin{proof}
Since $d\le p$,
\[
 B^*S
 \equiv B^*A^*I_p
 \equiv \beta^p A^*
 \pmod{x^d}.
\]
Multiplying \eqref{eq:pseudodiv-reversal-congruence} by
$\beta^{p-d}$ gives
\[
 B^*(\beta^{p-d}Q^*)
 \equiv\beta^pA^*\pmod{x^d}.
\]
Subtracting,
\[
 B^*(S-\beta^{p-d}Q^*)\equiv0\pmod{x^d}.
\]
Over $\mathbb Q[[x]]$, the series $B^*$ is invertible because its constant
coefficient $\beta$ is nonzero.  Hence
\[
 S\equiv\beta^{p-d}Q^*\pmod{x^d}.
\]
Both sides have degree less than $d$, so this congruence is equality.  Exact
coefficientwise divisibility follows at once.
\end{proof}

Only the highest $d$ coefficients of $A$ can enter $S$, and only the highest
$\min(n,p)$ coefficients of $B$ are needed to construct $I_p$.  Thus the
reversals in the algorithm below may be restricted to those tails.

\begin{algorithm}[Fast pseudo-division]\label{alg:pseudodiv-fast}
Input $A,B\in\ZZ[x]$ with $B\ne0$.
\begin{enumerate}
 \item If $\len(A)<\len(B)$, return $(0,A)$.
 \item Put $d\leftarrow\len(A)-\len(B)+1$ and let $p$ be the least power of
 two with $p\ge d$.
 \item Reverse the highest $d$ coefficients of $A$ to obtain the first $d$
 coefficients of $A^*$, and reverse the highest $\min(n,p)$ coefficients of
 $B$ to obtain the required initial part of $B^*$.
 \item Compute the scaled reciprocal $I_p$ of $B^*$.
 \item Compute
 \[
  S\leftarrow\operatorname{low}_d(A^*I_p),
  \qquad
  Q^*\leftarrow S/\beta^{p-d},
 \]
 where the scalar division is exact.
 \item Set $Q\leftarrow\operatorname{rev}_d(Q^*)$.
 \item Compute only
 \[
  T\leftarrow\operatorname{low}_{n-1}(BQ).
 \]
 \item Return $Q$ and
 \[
  R=\operatorname{low}_{n-1}(\beta^dA)-T.
 \]
\end{enumerate}
\end{algorithm}

\begin{theorem}\label{thm:pseudodiv-fast-correctness}
Algorithm~\ref{alg:pseudodiv-fast} returns exactly the pseudo-quotient and
pseudo-remainder of \eqref{eq:pseudodiv-identity}.
\end{theorem}

\begin{proof}
Lemma~\ref{lem:fast-pseudoquotient} proves that Step 5 recovers $Q^*$ exactly,
and hence Step 6 recovers $Q$.  By definition of pseudo-division,
\[
 R=\beta^dA-BQ
\]
has degree less than $n-1=\deg B$.  Therefore its complete coefficient list is
its low window of length $n-1$, and
\[
 R
 =\operatorname{low}_{n-1}(\beta^dA)
  -\operatorname{low}_{n-1}(BQ).
\]
Steps 7 and 8 compute exactly this expression.
\end{proof}

The remainder formula is important computationally.  A full product $BQ$ has
length up to $n+d-1$, but pseudo-division needs only its lowest $n-1$
coefficients once the quotient is known.

\subsection{Complexity and coefficient growth of fast pseudo-division}

Let $J(p)$ denote the scaled-inverse cost, let $L(s)$ denote a balanced low
product of precision $s$, and let
$L_{u,v}(r)$ denote the cost of the first $r$ coefficients of a product of
operand lengths $u$ and $v$.  Integer scalar powers, scalar exact divisions,
and coefficientwise scalar multiplications are displayed separately when bit
complexity matters.

\begin{theorem}\label{thm:pseudodiv-fast-complexity}
For $d>0$, Algorithm~\ref{alg:pseudodiv-fast} has coefficient-operation cost
\begin{equation}\label{eq:pseudodiv-fast-cost}
 J(p)+L(d)+L_{d,n}(n-1)+O(d+n),
 \qquad d\le p<2d.
\end{equation}
If low products of comparable lengths have cost $O(\mathsf M(N))$, and
$n,d=\Theta(N)$, then fast pseudo-division has coefficient-operation cost
$O(\mathsf M(N))$.
\end{theorem}

\begin{proof}
The scaled reciprocal costs $J(p)$.  Step 5 uses one low product of precision
$d$.  Step 7 is one low product with operand lengths $d,n$ and output length
$n-1$.  Tail extraction, reversals, additions, and scalar coefficient
operations are linear in $d+n$.  Equation \eqref{eq:pseudodiv-p-range} and
Theorem~\ref{thm:scaled-newton-correctness} give the balanced conclusion.
\end{proof}

The scaled algorithm does not eliminate the intrinsic coefficient growth
caused by the factor $\beta^d$; it eliminates the quadratic number of
coefficient operations used to obtain those coefficients.

\begin{lemma}\label{lem:pseudodiv-fast-intermediate-growth}
Let the input coefficient bit sizes be bounded by $\sigma$ for $A$ and $\tau$
for $B$.  In Algorithm~\ref{alg:pseudodiv-fast}, put
\[
 W=O\bigl(\sigma+d(\tau+1)+\log(d+1)\bigr).
\]
Then the scaled inverse, scaled reversed quotient product, pseudo-quotient,
and pseudo-remainder all have coefficient bit size $O(W)$, up to replacing
$d$ by $p<2d$ in intermediate scaled quantities.
\end{lemma}

\begin{proof}
Lemma~\ref{lem:scaled-inverse-integral} gives
$O(p(\tau+1))$ bits for coefficients of $I_p$.  A coefficient of
$S=\operatorname{low}_d(A^*I_p)$ is a sum of at most $d$ products, so it has
\[
 O\bigl(\sigma+p(\tau+1)+\log(d+1)\bigr)
\]
bits.  Since $p<2d$, this is $O(W)$.  Exact division by
$\beta^{p-d}$ cannot increase absolute values when $|\beta|\ge1$, so the same
bound applies to $Q$.

Finally, $\beta^dA$ has coefficient size
$O(\sigma+d\tau)$, while each coefficient of $BQ$ is a sum of at most
$\min(d,n)$ products.  Hence the low remainder window also has $O(W)$-bit
coefficients.  The sharper output bound of
Lemma~\ref{lem:pseudodiv-coefficient-growth} applies to the final canonical
$Q,R$; the extra logarithm above is only a convenient bound for product
intermediates.
\end{proof}

\begin{theorem}\label{thm:pseudodiv-fast-bit-complexity}
Let $\mathcal L_{u,v}(r,w)$ bound the bit complexity of a low product of
operand lengths $u,v$, output length $r$, and coefficient size at most $w$.
Let $\mathsf D_{\ZZ}(w)$ bound exact division of a $w$-bit integer by an
integer of at most $w$ bits.  Then fast pseudo-division has bit complexity
bounded by
\[
 J_{\rm bit}(p)
 +\mathcal L_{d,d}(d,O(W))
 +\mathcal L_{d,n}(n-1,O(W))
 +O\bigl(d\mathsf D_{\ZZ}(W)+(d+n)\MZ(W)\bigr).
\]
In the balanced regime $n,d=\Theta(N)$, if
\[
 \mathcal L_{N,N}(N,W)=O\bigl(N^\alpha\MZ(W)\bigr)
\]
for some $\alpha>1$ and
$\mathsf D_{\ZZ}(W)=O(\MZ(W))$, this simplifies to
\[
 O\!\left(
 N^\alpha\MZ\bigl(\sigma+N(\tau+1)\bigr)
 \right).
\]
\end{theorem}

\begin{proof}
The first three terms are the scaled reciprocal and the two low products.
There are $d$ exact coefficient divisions by $\beta^{p-d}$.  Constructing
$\beta^d$, scaling the low $n-1$ coefficients of $A$, and the remaining
linear coefficient work are covered by the final term.  The growth bound is
Lemma~\ref{lem:pseudodiv-fast-intermediate-growth}.  Under the balanced
hypotheses with $\alpha>1$, the polynomial low products dominate the linear
number of scalar integer operations.
\end{proof}

\begin{remark}
The bit-size factor
$\sigma+N(\tau+1)$ is not an artefact of the fast algorithm.  The defining
identity itself contains the factor $\beta^d$, so a fraction-free
pseudo-quotient or pseudo-remainder can naturally have coefficients with
$\Theta(d\tau)$ additional bits.  Fast pseudo-division reduces the
\emph{number} of polynomial coefficient operations; controlling coefficient
growth across a whole remainder sequence requires additional normalization,
such as primitive or subresultant PRS methods.
\end{remark}

\subsection{Complexity summary}

For a divisor of length $n$ and quotient length $d$, the algorithms of this
section have the following costs:
\[
\begin{array}{c|c}
\text{operation} & \text{coefficient-operation cost}\\
\hline
\text{classical pseudo-division} & O(d(n+d))\\
\text{scaled reciprocal to precision }p &
 J(p)=J(p/2)+2L(p)+O(p)\\
\text{fast pseudo-division} &
 J(p)+L(d)+L_{d,n}(n-1)+O(d+n),\quad d\le p<2d.
\end{array}
\]
Thus classical pseudo-division is quadratic in the balanced degree parameter,
whereas the scaled-Newton algorithm inherits the asymptotic complexity of low
multiplication.  Both algorithms produce the same canonical pseudo-quotient
and pseudo-remainder and hence share the same worst-case output coefficient
growth.



\section{Subresultant normalization of polynomial remainder sequences}
\label{sec:subresultant-prs}

A single pseudo-division has the coefficient growth described in
Section~\ref{sec:pseudodivision}.  A different issue appears when
pseudo-remainders are fed repeatedly into a Euclidean algorithm.  The powers
of successive leading coefficients introduced by pseudo-division can then
accumulate from one remainder to the next.  Polynomial remainder sequences
(PRSs) control this growth by dividing each pseudo-remainder by a suitable
scalar before continuing.  The classical theory is due to Collins and to
Brown--Traub; the particularly simple normalization used below is Brown's
subresultant PRS formulation \cite{BrownTraub1971,Brown1978}.

Let $I$ be a unique factorization domain with quotient field $K$; the case of
interest here is $I=\ZZ$.  Let
\[
 F_1,F_2,\ldots,F_k\in I[x],
 \qquad
 n_i=\deg F_i,
 \qquad
 f_i=\operatorname{lc}(F_i),
\]
with
\[
 n_1\ge n_2>n_3>\cdots>n_k\ge0.
\]
For $1\le i<k$ put
\[
 \delta_i=n_i-n_{i+1}.
\]
A polynomial remainder sequence is obtained by choosing nonzero scalars
$\beta_i\in K$ such that, for $i\ge3$,
\begin{equation}\label{eq:prs-defining-relation}
 \beta_i F_i
 =\operatorname{prem}(F_{i-2},F_{i-1}).
\end{equation}
The scalar multiplier implicit in the pseudo-remainder on the right is
\begin{equation}\label{eq:prs-alpha}
 \alpha_i=f_{i-1}^{\,n_{i-2}-n_{i-1}+1}
          =f_{i-1}^{\,\delta_{i-2}+1}.
\end{equation}
Thus there is a polynomial $Q_i$ such that
\[
 \beta_iF_i=\alpha_iF_{i-2}-Q_iF_{i-1}.
\]
Different choices of the $\beta_i$ give different integral associates of the
same Euclidean remainders.

\subsection{Primitive normalization}

For a nonzero polynomial
\[
 P=\sum_j p_jx^j\in\ZZ[x]
\]
define its positive content by
\[
 \operatorname{cont}(P)=\gcd_j|p_j|,
\]
and its primitive part by
\[
 \operatorname{pp}(P)=P/\operatorname{cont}(P).
\]
A \emph{primitive PRS} replaces every nonzero pseudo-remainder $P$ by
$\operatorname{pp}(P)$ before the next division.

\begin{proposition}\label{prop:primitive-prs-maximal-scalar}
Let $P\in\ZZ[x]$ be nonzero.  If $c\in\ZZ\setminus\{0\}$ and
$P/c\in\ZZ[x]$, then $|c|$ divides $\operatorname{cont}(P)$.  Consequently
primitive-part normalization removes the largest possible positive integral
scalar from $P$.  In particular no nonunit integer divides every coefficient
of $\operatorname{pp}(P)$.
\end{proposition}

\begin{proof}
The condition $P/c\in\ZZ[x]$ says that $c$ divides every coefficient $p_j$.
Hence $|c|$ is a common positive divisor of those coefficients and therefore
divides their greatest common divisor.  Dividing by the content leaves gcd
one, so no further nonunit scalar division is possible in $\ZZ[x]$.
\end{proof}

Thus primitive normalization is locally coefficient-minimal among integral
scalar associates.  Its disadvantage is computational: determining the
content requires a gcd computation among the coefficients of every
pseudo-remainder.  If the coefficient ring is itself a polynomial ring, this
can be much more expensive than the pseudo-division whose growth one is
trying to control.  The subresultant PRS replaces these repeated content gcds
by explicit exact factors depending only on degrees and preceding leading
coefficients.

\subsection{Subresultants and the fundamental theorem}

For polynomials $A,B$ of degrees $m\ge n$, the $j$th subresultant
$\operatorname{Sres}_j(A,B)$, $0\le j<n$, is obtained from the Sylvester
matrix by deleting the appropriate final rows and columns and taking the
standard determinant polynomial.  Two facts about this construction are
important here.  First,
\begin{equation}\label{eq:subresultant-determinant-order}
 \deg\operatorname{Sres}_j(A,B)\le j,
\end{equation}
and every coefficient of $\operatorname{Sres}_j(A,B)$ is a determinant of
order
\begin{equation}\label{eq:subresultant-order}
 m+n-2j
\end{equation}
whose entries are coefficients of $A$ and $B$.  Second, the nonzero
subresultants occur in exactly the degree positions dictated by a PRS.

We record the latter statement in the form needed below.  This is the
fundamental theorem of subresultants of Brown--Traub
\cite{BrownTraub1971}; Brown's later formulation \cite{Brown1978} gives the
normalization formulas particularly transparently.

\begin{theorem}[Fundamental theorem of subresultants]
\label{thm:fundamental-subresultants}
Let $F_1,F_2,\ldots,F_k$ be a PRS over $K$ with degrees $n_i$ as above.  For
$i=3,\ldots,k$ there exist nonzero scalars $\gamma_i,\theta_i\in K$ such that
\begin{align}
 \operatorname{Sres}_{n_{i-1}-1}(F_1,F_2)&=\gamma_iF_i,
 \label{eq:fundamental-subres-upper}\\
 \operatorname{Sres}_{n_i}(F_1,F_2)&=\theta_iF_i,
 \label{eq:fundamental-subres-lower}
\end{align}
and
\[
 \operatorname{Sres}_j(F_1,F_2)=0
 \quad\text{for}\quad
 n_{i-1}-1>j>n_i.
\]
The same statement at the final degree says that all subresultants below
$n_k$ vanish.
\end{theorem}

\begin{proof}
We give the determinant argument in the form relevant to the normalization
below.  A pseudo-division step
\[
 \alpha_iF_{i-2}=Q_iF_{i-1}+\beta_iF_i
\]
replaces the coefficient rows arising from shifted copies of $F_{i-2}$ in a
Sylvester submatrix by fraction-free linear combinations with shifted copies
of $F_{i-1}$.  These row operations eliminate the leading block belonging to
$F_{i-2}$ and leave, after extraction of powers of
$\alpha_i,\beta_i$ and $f_{i-1}$, the corresponding Sylvester submatrix for
the pair $(F_{i-1},F_i)$.  Hence subresultants of $(F_{i-2},F_{i-1})$ below
$n_{i-1}$ are scalar multiples of subresultants of $(F_{i-1},F_i)$.

Iterating this reduction down the PRS gives two possible nonzero boundary
minors at each degree drop: the minor of index $n_{i-1}-1$ and the one of
index $n_i$.  Both are scalar multiples of $F_i$; all intermediate minors
contain an eliminated zero block and therefore vanish.  This yields
\eqref{eq:fundamental-subres-upper} and
\eqref{eq:fundamental-subres-lower}.  Keeping the extracted scalar factors
through the row eliminations gives the explicit similarity factors
$\gamma_i$ and $\theta_i$ of Brown--Traub.  Their formulas are used in
Lemma~\ref{lem:brown-similarity-recurrence} below.  A complete expansion of
the same determinant elimination is given in
\cite{BrownTraub1971,Brown1978}.
\end{proof}

The theorem says that every PRS remainder is, over $K$, an associate of a
subresultant.  The purpose of the subresultant PRS is to choose the scalars
$\beta_i$ so that the associates are the subresultants themselves.

For the next formulas define
\[
 H_i=\operatorname{Sres}_{n_i}(F_1,F_2),
 \qquad
 h_i=\operatorname{lc}(H_i).
\]
Thus, when the upper similarity factors have been normalized to one,
$F_i=\operatorname{Sres}_{n_{i-1}-1}(F_1,F_2)$ while $H_i$ is the adjacent
lower subresultant.

\begin{lemma}[Brown's similarity-factor recurrence]
\label{lem:brown-similarity-recurrence}
Suppose $\gamma_3=\cdots=\gamma_i=1$.  Then
\begin{equation}\label{eq:brown-beta-condition}
 \gamma_{i+1}=1
\end{equation}
if and only if
\begin{equation}\label{eq:brown-beta}
 \beta_{i+1}
 =(-1)^{\delta_{i-1}+1}
   f_{i-1}h_{i-1}^{\,\delta_{i-1}}.
\end{equation}
Moreover
\begin{equation}\label{eq:brown-h-ratio}
 \frac{h_i}{h_{i-1}}
 =(-1)^{\delta_{i-1}+1}
   \frac{f_{i-1}}{f_i}
   \frac{\alpha_{i+1}}{\beta_{i+1}},
\end{equation}
where
$\alpha_{i+1}=f_i^{\delta_{i-1}+1}$ is the pseudo-division multiplier.
Consequently, under \eqref{eq:brown-beta},
\begin{equation}\label{eq:brown-h-recurrence}
 h_i=f_i^{\,\delta_{i-1}}
     h_{i-1}^{\,1-\delta_{i-1}}.
\end{equation}
\end{lemma}

\begin{proof}
The explicit similarity factors in the fundamental theorem are products of
the pseudo-division multipliers $\alpha_j$, the chosen normalizing factors
$\beta_j$, and powers of the leading coefficients $f_j$.  Brown's key
observation is that their accumulated products disappear when consecutive
factors are divided.  With the preceding upper factors normalized to one,
the quotient of two consecutive upper similarity factors simplifies to
\begin{equation}\label{eq:brown-gamma-ratio}
 \frac{\gamma_{i+1}}{\gamma_i}
 =\frac{(-1)^{\delta_{i-1}+1}
         f_{i-1}h_{i-1}^{\,\delta_{i-1}}}
        {\beta_{i+1}}.
\end{equation}
This is obtained directly by cancelling the common factors in the explicit
Brown--Traub formulas for $\gamma_i$ and $\gamma_{i+1}$; it is equation
(14), in ratio form, in Brown's derivation \cite{Brown1978}.  Therefore
$\gamma_{i+1}=\gamma_i=1$ if and only if
\eqref{eq:brown-beta} holds.

The analogous cancellation between the lower similarity factor $\theta_i$
and the next upper factor $\gamma_{i+1}$ gives
\eqref{eq:brown-h-ratio}; equivalently this is Brown's equation (18).
Thus no factors from PRS stages earlier than $i-1$ survive in either update.

Finally substitute
$\alpha_{i+1}=f_i^{\delta_{i-1}+1}$ and
\eqref{eq:brown-beta} into \eqref{eq:brown-h-ratio}:
\[
 \frac{h_i}{h_{i-1}}
 =\frac{f_i^{\delta_{i-1}}}
        {h_{i-1}^{\delta_{i-1}}}.
\]
Multiplication by $h_{i-1}$ gives
\eqref{eq:brown-h-recurrence}.
\end{proof}

The cancellation proof is short because the complicated accumulated powers
in the general similarity factors disappear when consecutive factors are
compared.  This is the main simplification in Brown's formulation.

\subsection{Brown's subresultant PRS}

Assume now that $F_1,F_2\in\ZZ[x]$ are primitive.  Define
\begin{equation}\label{eq:brown-h-initial}
 h_2=f_2^{\delta_1}.
\end{equation}
Then use \eqref{eq:brown-h-recurrence} for subsequent $h_i$.

\begin{algorithm}[Brown subresultant PRS]
\label{alg:brown-subresultant-prs}
Input primitive $F_1,F_2\in\ZZ[x]$ with
$\deg F_1\ge\deg F_2$.
\begin{enumerate}
 \item Put $f_i\leftarrow\operatorname{lc}(F_i)$ and
 $\delta_i\leftarrow\deg F_i-\deg F_{i+1}$ whenever these quantities become
 defined.
 \item Compute
 \begin{equation}\label{eq:brown-first-remainder}
  F_3=(-1)^{\delta_1+1}
      \operatorname{prem}(F_1,F_2),
 \end{equation}
 and set $h_2=f_2^{\delta_1}$.
 \item For $i=4,5,\ldots$ while the pseudo-remainder is nonzero, compute
 \begin{equation}\label{eq:brown-prs-step}
  F_i=
  \frac{(-1)^{\delta_{i-2}+1}
        \operatorname{prem}(F_{i-2},F_{i-1})}
       {f_{i-2}h_{i-2}^{\delta_{i-2}}},
 \end{equation}
 where the scalar division is exact, and update
 \begin{equation}\label{eq:brown-prs-h-update}
  h_{i-1}
  =f_{i-1}^{\delta_{i-2}}
   h_{i-2}^{1-\delta_{i-2}}.
 \end{equation}
 \item Stop when the next pseudo-remainder is zero.
\end{enumerate}
\end{algorithm}

Formula \eqref{eq:brown-prs-h-update} may contain a negative exponent when a
degree drop exceeds one.  It is therefore a statement about an exact scalar
division, not about arithmetic in a Laurent ring.  The theorem below proves
that these divisions are exact in $\ZZ$.

\begin{theorem}\label{thm:brown-subresultant-prs-correct}
Algorithm~\ref{alg:brown-subresultant-prs} is well defined over $\ZZ$.  Every
scalar division in \eqref{eq:brown-prs-step} and
\eqref{eq:brown-prs-h-update} is exact, and for $i\ge3$,
\begin{equation}\label{eq:brown-prs-is-subresultant}
 F_i=\operatorname{Sres}_{n_{i-1}-1}(F_1,F_2)
\end{equation}
up to the fixed sign convention used in the definition of the
subresultants.  Moreover
\begin{equation}\label{eq:brown-h-is-subresultant-lc}
 h_i=\operatorname{lc}
 \operatorname{Sres}_{n_i}(F_1,F_2).
\end{equation}
\end{theorem}

\begin{proof}
For $i=3$, the Brown--Traub similarity factor in
Theorem~\ref{thm:fundamental-subresultants} becomes one after the sign choice
in \eqref{eq:brown-first-remainder}; hence $F_3$ is the first nonzero
subresultant in its degree block.  Equation \eqref{eq:brown-h-initial} is the
leading coefficient of the adjacent lower subresultant
$\operatorname{Sres}_{n_2}(F_1,F_2)$.

Assume inductively that all upper similarity factors through $\gamma_i$ are
one and that $h_{i-1}$ is the leading coefficient of the corresponding lower
subresultant.  Lemma~\ref{lem:brown-similarity-recurrence} says that choosing
\[
 \beta_{i+1}
 =(-1)^{\delta_{i-1}+1}
   f_{i-1}h_{i-1}^{\delta_{i-1}}
\]
makes $\gamma_{i+1}=1$.  Therefore the next quotient of the pseudo-remainder
is exactly the next subresultant.  Since a subresultant is a determinant with
integer entries, it lies in $\ZZ[x]$; hence the scalar division producing it
is exact.  The same lemma gives the update
\[
 h_i=f_i^{\delta_{i-1}}h_{i-1}^{1-\delta_{i-1}},
\]
and Theorem~\ref{thm:fundamental-subresultants} identifies this $h_i$ as the
leading coefficient of the adjacent integer subresultant.  Thus any division
implicit in the negative exponent is also exact.  Induction proves all
claims.
\end{proof}

This theorem gives the promised replacement for repeated content
computations.  Primitive PRS normalization divides by the full, and therefore
maximal, content of each individual pseudo-remainder.  Brown's factors need
not equal that full content.  Instead they remove the part forced by the
subresultant structure, and do so using only preceding degrees and leading
coefficients.  The resulting remainders are canonical determinants rather
than arbitrary integral associates.

\begin{proposition}[Uniqueness of the subresultant normalization]
\label{prop:subresultant-normalization-unique}
Fix the initial pair $F_1,F_2$, the pseudo-remainder convention and the signs
in Algorithm~\ref{alg:brown-subresultant-prs}.  Among PRSs of the form
\[
 F_i=\operatorname{prem}(F_{i-2},F_{i-1})/\beta_i,
\]
the requirement that every $F_i$ be the corresponding subresultant determines
$\beta_i$ uniquely, up to multiplication by a unit of the coefficient ring.
Over $\ZZ$, after the signs have been fixed, the factors are precisely those
of Brown's recurrence.
\end{proposition}

\begin{proof}
Once $F_{i-2}$ and $F_{i-1}$ have been fixed, the pseudo-remainder
$P=\operatorname{prem}(F_{i-2},F_{i-1})$ is fixed.  If the next member is
required to be a specified nonzero subresultant $S$, then
\[
 P=\beta_iS.
\]
Since $\ZZ[x]$ is an integral domain, the scalar $\beta_i$ is uniquely
determined, apart from the unit by which one chooses to normalize $S$.
Theorem~\ref{thm:brown-subresultant-prs-correct} shows that Brown's factors
produce exactly these $S$, so they are the unique factors with this property.
\end{proof}

This is the precise sense in which the Brown--Collins normalization is
canonical, and the sense in which it is often described as the best choice
within the family of PRS algorithms based on cheaply computed scalar factors
\cite{Collins1967,Brown1978}.  It should not be confused with maximal content
removal: accidental content may remain, and primitive normalization can divide
it out if one is willing to compute coefficient gcds.  Brown also describes
``improved PRS'' variants in which an additional divisor of the content, when
it is already known at negligible cost, is removed on top of the subresultant
normalization \cite{Brown1978}.

\subsection{Coefficient growth}

The determinant interpretation immediately explains why subresultant
normalization controls coefficient growth.

\begin{theorem}\label{thm:subresultant-coefficient-growth}
Let $A,B\in\ZZ[x]$ have degrees at most $N$ and coefficient bit sizes at most
$\tau$.  Every coefficient of every subresultant of $A$ and $B$, and hence
every coefficient in their subresultant PRS, has
\begin{equation}\label{eq:subresultant-bit-bound}
 O\bigl(N(\tau+\log(N+1))\bigr)
\end{equation}
bits.
\end{theorem}

\begin{proof}
Fix the $j$th subresultant.  By
\eqref{eq:subresultant-determinant-order}, each of its coefficients is a
determinant of order
\[
 r=\deg A+\deg B-2j\le2N.
\]
Every matrix entry is an input coefficient and therefore has absolute value
less than $2^\tau$.  Each row contains at most $N+1$ nonzero coefficient
entries, so its Euclidean norm is at most
\[
 \sqrt{N+1}\,2^\tau.
\]
Hadamard's determinant bound gives
\[
 |D|
 \le\bigl(\sqrt{N+1}\,2^\tau\bigr)^r.
\]
Taking binary logarithms and using $r\le2N$ gives
\[
 \bits(D)
 \le r\left(\tau+\frac12\log_2(N+1)\right)+O(1)
 =O\bigl(N(\tau+\log(N+1))\bigr).
\]
By Theorem~\ref{thm:brown-subresultant-prs-correct}, every PRS remainder
produced by Brown's normalization is one of these subresultants, so the same
bound applies to the whole sequence.
\end{proof}

The important contrast is with an unnormalized pseudo-PRS.  A pseudo-remainder
step multiplies the dividend by a power of the divisor's leading coefficient,
and the resulting enlarged coefficients become the input to the next step.
Those powers can therefore compound through the sequence.  Subresultant
normalization removes exactly the systematic factors required to keep each
remainder inside the determinant bound
\eqref{eq:subresultant-bit-bound}, without computing a coefficient gcd at
each stage.  This is the sense in which the Brown normalization controls
intermediate expression swell.

\begin{remark}
Primitive and subresultant PRSs optimize different costs.  Primitive
normalization gives the smallest integral scalar associate at every individual
step, but finding the scalar requires content computation.  The subresultant
PRS uses explicitly predictable exact scalars and guarantees determinant-size
coefficients.  It may leave accidental content which is not forced by the
subresultant identities.  Thus the Brown factors should not be described as
larger than the content or as producing universally smaller coefficients;
their advantage is that they remove the systematic content at negligible gcd
cost.
\end{remark}




\section{Greatest common divisors over $\ZZ[x]$}\label{sec:zz-poly-gcd}

The pseudo-division and polynomial-remainder-sequence machinery developed in
Sections~\ref{sec:pseudodivision} and the preceding section gives two direct
fraction-free algorithms for polynomial greatest common divisors.  The first
normalizes every pseudo-remainder by its full content.  The second uses Brown's
subresultant normalization and therefore avoids coefficient-gcd computations at
every remainder step.

Throughout this section the gcd of two polynomials in $\ZZ[x]$ is normalized to
have positive leading coefficient.  We use the conventions
\[
 \gcd(A,0)=\operatorname{sgn}(\operatorname{lc}(A))A,
 \qquad
 \gcd(0,0)=0,
\]
with the evident interpretation when $A$ is a nonzero constant.  For the zero
polynomial put
\[
 \operatorname{cont}(0)=0,
 \qquad
 \operatorname{pp}(0)=0.
\]
For a nonzero polynomial, content and primitive part are as in the preceding
section.

\subsection{Gauss's lemma and reduction to primitive inputs}

The basic structural fact is Gauss's lemma.  We include the proof because it is
what permits a gcd over $\ZZ[x]$ to be recovered from a fraction-free Euclidean
calculation on primitive polynomials.

\begin{lemma}[Gauss]\label{lem:gauss-primitive-product}
If $F,G\in\ZZ[x]$ are primitive, then $FG$ is primitive.
Consequently, for arbitrary nonzero $F,G\in\ZZ[x]$,
\begin{equation}\label{eq:content-multiplicative}
 \operatorname{cont}(FG)
 =\operatorname{cont}(F)\operatorname{cont}(G).
\end{equation}
\end{lemma}

\begin{proof}
Suppose first that $F$ and $G$ are primitive and that a prime $p$ divides every
coefficient of $FG$.  Reducing coefficients modulo $p$ gives
\[
 \overline F\,\overline G=0
 \quad\text{in}\quad \mathbb F_p[x].
\]
Since $F$ is primitive, not all of its coefficients are divisible by $p$, so
$\overline F\ne0$; similarly $\overline G\ne0$.  This contradicts the fact that
$\mathbb F_p[x]$ is an integral domain.  Hence no prime divides every
coefficient of $FG$, and $FG$ is primitive.

For general nonzero $F,G$, write
\[
 F=aF_0,\qquad G=bG_0,
\]
where $a=\operatorname{cont}(F)>0$, $b=\operatorname{cont}(G)>0$, and
$F_0,G_0$ are primitive.  Then $F_0G_0$ is primitive, so the content of
$FG=abF_0G_0$ is exactly $ab$.
\end{proof}

\begin{corollary}[Gauss divisibility lemma]\label{cor:gauss-divisibility}
Let $P,F\in\ZZ[x]$ with $P$ primitive.  If $P$ divides $F$ in $\mathbb Q[x]$,
then $P$ divides $F$ in $\ZZ[x]$.
\end{corollary}

\begin{proof}
Write $F=PQ$ with $Q\in\mathbb Q[x]$.  Choose coprime integers $a,b$, with
$b>0$, and a primitive $Q_0\in\ZZ[x]$ such that
\[
 Q=\frac ab Q_0.
\]
Then
\[
 bF=aPQ_0.
\]
By Lemma~\ref{lem:gauss-primitive-product}, $PQ_0$ is primitive.  Hence the gcd
of its coefficients is one.  Since every coefficient of $aPQ_0$ is divisible
by $b$, integrality of $F$ implies $b\mid a$.  Coprimality gives $b=1$, so
$Q\in\ZZ[x]$.
\end{proof}

The content and primitive parts of the inputs may now be separated completely.

\begin{theorem}[Content--primitive decomposition of the gcd]
\label{thm:gcd-content-primitive}
Let $A,B\in\ZZ[x]$, not both zero, and put
\[
 a=\operatorname{cont}(A),\qquad
 b=\operatorname{cont}(B),\qquad
 c=\gcd(a,b).
\]
Let
\[
 A_0=\operatorname{pp}(A),\qquad
 B_0=\operatorname{pp}(B),
\]
and let $G_0$ be the primitive polynomial with positive leading coefficient
which is associated in $\mathbb Q[x]$ to $\gcd_{\mathbb Q[x]}(A_0,B_0)$.
Then
\begin{equation}\label{eq:gcd-content-primitive}
 \gcd_{\ZZ[x]}(A,B)=cG_0.
\end{equation}
\end{theorem}

\begin{proof}
By Corollary~\ref{cor:gauss-divisibility}, the primitive polynomial $G_0$
divides both $A_0$ and $B_0$ in $\ZZ[x]$.  Hence $cG_0$ divides both $A$ and
$B$.

Conversely, let $D$ divide both $A$ and $B$ in $\ZZ[x]$.  Write
$D=dD_0$ with $d=\operatorname{cont}(D)>0$ and $D_0$ primitive.  By
\eqref{eq:content-multiplicative}, $d$ divides both $a$ and $b$, hence $d\mid c$.
The primitive factor $D_0$ divides both $A_0$ and $B_0$ in $\mathbb Q[x]$, so
it divides $G_0$ there; Corollary~\ref{cor:gauss-divisibility} promotes this to
divisibility in $\ZZ[x]$.  Therefore $D$ divides $cG_0$, up to the irrelevant
unit $-1$.  Positive leading coefficient fixes that unit.
\end{proof}

Thus the polynomial part of the problem reduces to computing the primitive gcd
of two primitive inputs.  This reduction is standard; see, for example,
\cite{vonZurGathenGerhard2013}.  The two algorithms below differ only in how
they normalize the pseudo-remainders.

\subsection{Primitive pseudo-Euclidean gcd}\label{subsec:primitive-gcd}

Let $U,V\in\ZZ[x]$ with $\deg U\ge\deg V\ge0$, and put
\[
 \delta=\deg U-\deg V+1.
\]
Pseudo-division gives
\begin{equation}\label{eq:gcd-pseudorem-identity}
 \operatorname{lc}(V)^\delta U=QV+R,
 \qquad \deg R<\deg V.
\end{equation}
The polynomial $R$ is the pseudo-remainder $\operatorname{prem}(U,V)$.

\begin{lemma}[Pseudo-remainders preserve the rational gcd]
\label{lem:pseudorem-preserves-gcd}
If $V\ne0$ and $R=\operatorname{prem}(U,V)$, then
\[
 \gcd_{\mathbb Q[x]}(U,V)
 \sim
 \gcd_{\mathbb Q[x]}(V,R),
\]
where $\sim$ means equality up to a nonzero rational scalar.  If $R\ne0$, the
same is true after replacing $R$ by $\operatorname{pp}(R)$.
\end{lemma}

\begin{proof}
The leading coefficient of $V$ is nonzero, and hence is a unit in
$\mathbb Q[x]$.  Equation~\eqref{eq:gcd-pseudorem-identity} therefore implies
that every common divisor of $U,V$ divides $R$.  Conversely, the same identity
solved for $U$ shows that every common divisor of $V,R$ divides $U$ in
$\mathbb Q[x]$.  Thus the two pairs have associated gcds.  Replacing a nonzero
$R$ by $\operatorname{pp}(R)$ only multiplies it by a nonzero rational scalar,
which cannot change a gcd in $\mathbb Q[x]$.
\end{proof}

\begin{algorithm}[Primitive pseudo-Euclidean gcd]
\label{alg:primitive-pseudo-gcd}
Input $A,B\in\ZZ[x]$.
\begin{enumerate}
 \item If $A=B=0$, return $0$.
 \item Compute
 \[
  c\leftarrow\gcd(\operatorname{cont}(A),\operatorname{cont}(B)),
  \qquad
  U\leftarrow\operatorname{pp}(A),
  \quad
  V\leftarrow\operatorname{pp}(B).
 \]
 \item If necessary interchange $U,V$ so that $\deg U\ge\deg V$.
 \item While $V\ne0$:
 \begin{enumerate}
  \item $R\leftarrow\operatorname{prem}(U,V)$;
  \item $P\leftarrow\operatorname{pp}(R)$;
  \item $(U,V)\leftarrow(V,P)$.
 \end{enumerate}
 \item Replace $U$ by $-U$ if its leading coefficient is negative.
 \item Return $cU$.
\end{enumerate}
\end{algorithm}

\begin{theorem}[Correctness of primitive pseudo-Euclidean gcd]
\label{thm:primitive-pseudo-gcd-correct}
Algorithm~\ref{alg:primitive-pseudo-gcd} terminates and returns the canonical
gcd of $A$ and $B$ in $\ZZ[x]$.
\end{theorem}

\begin{proof}
At every nonzero remainder step,
\[
 \deg\operatorname{pp}(R)=\deg R<\deg V,
\]
so the second degree decreases strictly.  The loop therefore terminates.

By Lemma~\ref{lem:pseudorem-preserves-gcd}, every replacement
$(U,V)\leftarrow(V,\operatorname{pp}(R))$ preserves the gcd in $\mathbb Q[x]$
up to a nonzero scalar.  When the loop stops, the pair is $(U,0)$, so $U$ is
associated in $\mathbb Q[x]$ to the rational gcd of the original primitive
parts.  Every nonzero polynomial inserted into the second position is
primitive, and the first position is a previous second position.  Hence the
final $U$ is primitive.  Fixing the sign therefore makes it exactly the
primitive polynomial $G_0$ of
Theorem~\ref{thm:gcd-content-primitive}.  Multiplication by the integer content
gcd $c$ then gives the canonical gcd of $A$ and $B$.
\end{proof}

The degree complexity is quadratic when pseudo-remainders are computed by the
classical algorithm.  The following count is useful because it remains valid
when degree drops are irregular.

\begin{theorem}[Coefficient complexity of primitive gcd]
\label{thm:primitive-gcd-coefficient-complexity}
Suppose the primitive inputs have degrees at most $N$.  With classical
pseudo-remainders, Algorithm~\ref{alg:primitive-pseudo-gcd} uses
$O(N^2)$ coefficient additions and multiplications, processes $O(N^2)$
coefficients in content computations, and performs $O(N^2)$ exact scalar
coefficient divisions when contents are removed.
\end{theorem}

\begin{proof}
Let the successive nonzero PRS degrees be
\[
 n_0>n_1>\cdots>n_s\ge0,
\]
and put
\[
 \Delta_i=n_{i-1}-n_i\ge1.
\]
At the division of a polynomial of degree $n_{i-1}$ by one of degree $n_i$,
the pseudo-quotient length is $d_i=\Delta_i+1$.  The classical pseudo-division
bound from Theorem~\ref{thm:pseudodiv-classical-complexity} is
\[
 O\bigl(d_i(n_i+d_i)\bigr)
 =O\bigl((\Delta_i+1)(n_{i-1}+1)\bigr).
\]
Now
\[
 \sum_i\Delta_i\le N,
 \qquad
 s\le N,
\]
so
\[
 \sum_i(\Delta_i+1)\le2N.
\]
Since $n_{i-1}+1\le N+1$, summing the division costs gives $O(N^2)$.

A content computation scans all coefficients of one pseudo-remainder.  There
are at most $N$ remainders and each has at most $N+1$ coefficients, so at most
$O(N^2)$ coefficients enter coefficient-gcd calculations.  Dividing each
pseudo-remainder by its content touches the same number of coefficients.
\end{proof}

Primitive normalization also gives the same determinant-size bound for the
\emph{stored} PRS members as subresultant normalization.

\begin{lemma}[Size of primitive PRS members]
\label{lem:primitive-prs-subresultant-size}
Let the primitive inputs have degree at most $N$ and coefficient bit size at
most $\tau$.  Every nonzero polynomial stored by the primitive PRS has
coefficients of
\begin{equation}\label{eq:primitive-prs-bit-bound}
 O\bigl(N(\tau+\log(N+1))\bigr)
\end{equation}
bits.
\end{lemma}

\begin{proof}
By the fundamental theorem of subresultants,
Theorem~\ref{thm:fundamental-subresultants}, every member of any PRS is a
nonzero rational scalar multiple of the appropriate subresultant.  If two
nonzero integer polynomials are rational scalar multiples of one another,
their primitive parts agree up to sign: writing each as its positive content
times a primitive polynomial and applying Gauss's lemma leaves only a rational
unit between two primitive integer polynomials, hence that unit is $\pm1$.
Therefore every primitive-PRS member is, up to sign, the primitive part of a
subresultant.  Dividing an integer polynomial by its positive content cannot
increase the absolute value of any coefficient.  The determinant bound
\eqref{eq:subresultant-bit-bound} now gives
\eqref{eq:primitive-prs-bit-bound}.
\end{proof}

The simple companion implementation first materializes a raw pseudo-remainder
and only then removes its content.  Consequently its transient integers can be
larger than the normalized bound above.  The next theorem records a direct
worst-case bit bound for this implementation style.

\begin{theorem}[Direct bit bound for primitive gcd]
\label{thm:primitive-gcd-bit-complexity}
Put
\[
 W=O\bigl(N(\tau+\log(N+1))\bigr),
 \qquad
 B=O\bigl(N(W+1)\bigr).
\]
Let $\mathsf G_{\ZZ}(t)$ and $\mathsf D_{\ZZ}(t)$ bound respectively the bit
complexity of gcd and exact division for $t$-bit integers.  A direct
implementation which forms each classical pseudo-remainder before taking its
content has bit complexity
\begin{equation}\label{eq:primitive-gcd-bit-complexity}
 O\!\left(
 N^2\bigl(\MZ(B)+\mathsf G_{\ZZ}(B)+\mathsf D_{\ZZ}(B)\bigr)
 \right).
\end{equation}
\end{theorem}

\begin{proof}
Lemma~\ref{lem:primitive-prs-subresultant-size} bounds the stored inputs to each
pseudo-division by $W$ bits.  At a step with pseudo-quotient length $d_i$,
Lemma~\ref{lem:pseudodiv-coefficient-growth} bounds the raw pseudo-remainder by
$O(d_i(W+1))\le B$ bits.  The coefficient-operation count over all
pseudo-divisions is $O(N^2)$ by
Theorem~\ref{thm:primitive-gcd-coefficient-complexity}; monotonicity allows
every integer multiplication to be charged $\MZ(B)$.

The same theorem shows that $O(N^2)$ coefficients are scanned while contents
are computed and then divided out.  Charging each coefficient-gcd operation by
$\mathsf G_{\ZZ}(B)$ and each exact scalar division by
$\mathsf D_{\ZZ}(B)$ proves the bound.  This is deliberately a worst-case
bound for the straightforward fraction-free implementation: most steps have
much smaller degree gaps, and hence much smaller transient integers.
\end{proof}

\subsection{Subresultant gcd}\label{subsec:subresultant-gcd}

Primitive PRS removes the largest possible scalar at each step, but pays for a
coefficient gcd every time.  Brown's subresultant PRS replaces those gcds by
the explicit factors proved exact in the preceding section.  The gcd algorithm
therefore has exactly the same outer structure; only the normalization of the
remainder sequence changes.

\begin{algorithm}[Subresultant gcd]
\label{alg:subresultant-gcd}
Input $A,B\in\ZZ[x]$.
\begin{enumerate}
 \item If $A=B=0$, return $0$.
 \item Compute
 \[
  c\leftarrow\gcd(\operatorname{cont}(A),\operatorname{cont}(B)),
  \qquad
  U\leftarrow\operatorname{pp}(A),
  \quad
  V\leftarrow\operatorname{pp}(B).
 \]
 \item If one of $U,V$ is zero, let $S$ be the other one.  Otherwise interchange
 $U,V$ if necessary so that $\deg U\ge\deg V$, and compute Brown's
 subresultant PRS from $U,V$ until its last nonzero member $S$.
 \item Put $G\leftarrow\operatorname{pp}(S)$ and replace $G$ by $-G$ if its
 leading coefficient is negative.
 \item Return $cG$.
\end{enumerate}
\end{algorithm}

\begin{theorem}[Last nonzero subresultant and the gcd]
\label{thm:last-subresultant-gcd}
Let $U,V\in\ZZ[x]$ be primitive and nonzero.  The last nonzero member $S$ of
Brown's subresultant PRS is associated in $\mathbb Q[x]$ to
$\gcd_{\mathbb Q[x]}(U,V)$.  Consequently
\[
 \operatorname{pp}(S)
\]
is, up to sign, the primitive gcd of $U,V$ in $\ZZ[x]$.
\end{theorem}

\begin{proof}
A Brown subresultant PRS is still a polynomial remainder sequence: every step
has the form
\[
 \beta_iF_i=\alpha_iF_{i-2}-Q_iF_{i-1},
\]
with nonzero integer scalars $\alpha_i,\beta_i$.  Over $\mathbb Q[x]$ these
scalars are units.  The same argument as in
Lemma~\ref{lem:pseudorem-preserves-gcd} therefore shows that each consecutive
pair has the same rational gcd, up to associates.  At termination the pair is
$(S,0)$, so $S$ is associated to the original rational gcd.

Theorem~\ref{thm:brown-subresultant-prs-correct} additionally says that $S$ is
an actual integer subresultant.  Taking its primitive part removes the remaining
integer scalar.  Gauss's divisibility lemma then identifies the resulting
primitive polynomial with the primitive gcd in $\ZZ[x]$.
\end{proof}

\begin{corollary}[Correctness of subresultant gcd]
\label{cor:subresultant-gcd-correct}
Algorithm~\ref{alg:subresultant-gcd} returns the canonical gcd of $A$ and $B$
in $\ZZ[x]$.
\end{corollary}

\begin{proof}
Theorem~\ref{thm:last-subresultant-gcd} computes the primitive polynomial
$G_0$ occurring in Theorem~\ref{thm:gcd-content-primitive}; the sign
normalization makes it canonical.  Multiplication by the input-content gcd $c$
then gives \eqref{eq:gcd-content-primitive}.
\end{proof}

\begin{theorem}[Coefficient complexity of subresultant gcd]
\label{thm:subresultant-gcd-coefficient-complexity}
For primitive inputs of degrees at most $N$, Brown's subresultant gcd with
classical pseudo-remainders uses $O(N^2)$ coefficient additions and
multiplications and $O(N^2)$ exact coefficient divisions by known Brown
normalization factors.  Apart from the two initial input contents, it requires
no coefficient-gcd computations.
\end{theorem}

\begin{proof}
The pseudo-division cost is summed exactly as in
Theorem~\ref{thm:primitive-gcd-coefficient-complexity}, giving $O(N^2)$
coefficient arithmetic.  Each normalized pseudo-remainder has at most $N+1$
coefficients, and there are at most $N$ PRS steps, so exact division by the
known Brown factor touches $O(N^2)$ coefficients in total.  The auxiliary
leading-coefficient powers and the $h_i$ updates require only $O(N)$ scalar
updates and do not change the quadratic coefficient count.  No content is
computed for an intermediate PRS member.
\end{proof}

By Theorem~\ref{thm:subresultant-coefficient-growth}, every normalized remainder
has coefficients of $W=O(N(\tau+\log(N+1)))$ bits.  As for the primitive PRS,
the straightforward companion implementation first forms a raw
pseudo-remainder before applying Brown's exact divisor, so a large degree gap
can create larger transient integers.

\begin{theorem}[Direct bit bound for subresultant gcd]
\label{thm:subresultant-gcd-bit-complexity}
Let
\[
 W=O\bigl(N(\tau+\log(N+1))\bigr),
 \qquad
 B=O\bigl(N(W+1)\bigr).
\]
With the notation of
Theorem~\ref{thm:primitive-gcd-bit-complexity}, the direct Brown-PRS gcd has bit
complexity
\begin{equation}\label{eq:subresultant-gcd-bit-complexity}
 O\!\left(
 N^2\bigl(\MZ(B)+\mathsf D_{\ZZ}(B)\bigr)
 +N\bigl(\mathsf G_{\ZZ}(\tau)+\mathsf D_{\ZZ}(\tau)\bigr)
 \right).
\end{equation}
\end{theorem}

\begin{proof}
The stored subresultants satisfy the $W$-bit determinant bound.  A classical
pseudo-division with quotient length $d_i$ produces a raw pseudo-remainder of
$O(d_i(W+1))\le B$ bits by
Lemma~\ref{lem:pseudodiv-coefficient-growth}.  The total coefficient arithmetic
is $O(N^2)$, so integer multiplications contribute
$O(N^2\MZ(B))$.

Brown's normalizing scalar divides every coefficient of the raw
pseudo-remainder exactly.  There are $O(N^2)$ such coefficient divisions in
total, each involving integers of at most $O(B)$ bits.  The powers used to form
the normalizing scalar have at most the same order of bit size and are absorbed
by this bound.  Only the two initial contents require coefficient gcds; forming the two
primitive input parts also requires exact scalar division of $O(N)$ input
coefficients.  Their coefficients have at most $\tau$ bits, giving the final
term.
\end{proof}

The value of the subresultant algorithm is now visible at both the algebraic
and bit-operation levels.  Primitive normalization computes the strongest
possible scalar reduction but invokes integer gcd repeatedly.  Brown
normalization leaves possible accidental content, yet it removes the entire
systematic pseudo-division factor by exact divisions whose divisors are known in
advance.  Both sequences have determinant-size normalized remainders; the
subresultant sequence reaches that scale without content gcds at every step.
Brown's discussion of this tradeoff is explicit in \cite{Brown1978}.

\subsection{Complexity summary}

For input degrees at most $N$, both fraction-free Euclidean gcd algorithms in
this section use $O(N^2)$ polynomial coefficient operations when classical
pseudo-remainders are used.  Their normalized PRS members have
\[
 O\bigl(N(\tau+\log(N+1))\bigr)
\]
bit coefficients.  The difference is normalization:
\[
\begin{array}{c|c|c}
 \text{algorithm} & \text{intermediate normalization} & \text{extra scalar work}\\
 \hline
 \text{primitive PRS} & \text{full content} & O(N^2)\ \text{integer gcds}\\
 \text{subresultant PRS} & \text{Brown factor} & O(N^2)\ \text{exact divisions}.
\end{array}
\]
The direct bit bounds above also account for the temporary coefficient growth
inside an unfused classical pseudo-remainder computation.  Faster asymptotic
gcd algorithms require a different organization of the Euclidean steps; the
fraction-free half-gcd method will be treated separately.




\section{Extended gcd and subresultant B\'ezout certificates}
\label{sec:zz-poly-xgcd}

Over a field, the extended Euclidean algorithm returns the monic gcd together
with B\'ezout cofactors.  Over $\ZZ[x]$ these two requirements are no longer
compatible in general.  The ring $\ZZ[x]$ is a unique factorization domain,
so gcds exist, but it is not a B\'ezout domain.  Consequently the polynomial
gcd need not belong to the ideal generated by the two inputs.

The extended algorithm used here therefore has a different contract.  Given
$A,B\in\ZZ[x]$, it returns polynomials
\[
 H,U,V\in\ZZ[x]
\]
such that
\begin{equation}\label{eq:xgcd-certificate}
 H=UA+VB,
\end{equation}
where $H$ is the last nonzero member of the Brown-normalized subresultant PRS
of the ordered pair $(A,B)$.  In particular, $H$ is associated over
$\mathbb Q[x]$ to the polynomial gcd, but it is not in general the canonical
gcd in $\ZZ[x]$.

\subsection{Why the gcd need not have integral B\'ezout cofactors}

The simplest example already occurs for a constant and a linear polynomial.

\begin{proposition}\label{prop:zpoly-not-bezout}
The ring $\ZZ[x]$ is not a B\'ezout domain.  In particular,
\[
 \gcd_{\ZZ[x]}(2,x)=1,
\]
but there do not exist $U,V\in\ZZ[x]$ such that
\[
 2U+xV=1.
\]
\end{proposition}

\begin{proof}
Any common divisor of $2$ and $x$ must divide the primitive polynomial $x$
and the nonzero constant $2$, so it is a unit.  Thus the gcd is $1$.
On the other hand, the constant term of $2U+xV$ is twice the constant term of
$U$, hence is even.  It cannot equal the constant term $1$ of the polynomial
$1$.  Therefore $1\notin(2,x)$.
\end{proof}

Thus a polynomial XGCD routine over $\ZZ[x]$ cannot promise
\[
 UA+VB=\gcd(A,B)
\]
for all inputs while insisting that $U,V$ remain integral.  It must either
work over $\mathbb Q[x]$, or return an integral multiple of the rational gcd.
The companion implementation takes the latter approach.

There is also a practical reason not to reduce the inputs to primitive parts
before running the extended PRS.  Primitive normalization divides a remainder
by its full content.  Although that quotient polynomial is integral, its
particular B\'ezout cofactors need not be divisible by the same content.
Brown normalization is different: the subresultant theory predicts a common
factor not only in the pseudo-remainder, but simultaneously in the canonical
B\'ezout cofactors.  This is the exact divisibility which makes a completely
integral extended PRS possible.

\subsection{Subresultant B\'ezout identities}

Let $A,B\in\ZZ[x]$ have degrees $m\ge n$, and let
$\operatorname{Sres}_j(A,B)$ denote the $j$th subresultant as in
Section~\ref{sec:subresultant-prs}.  The determinant construction provides
cofactors at the same time as the subresultant itself.

\begin{theorem}[Subresultant B\'ezout identity]
\label{thm:subresultant-bezout}
For every $0\le j<n$ there exist polynomials
$C_j,D_j\in\ZZ[x]$ such that
\begin{equation}\label{eq:subresultant-bezout}
 \operatorname{Sres}_j(A,B)=C_jA+D_jB,
\end{equation}
with degree bounds
\begin{equation}\label{eq:subresultant-bezout-degrees}
 \deg C_j\le n-j-1,
 \qquad
 \deg D_j\le m-j-1.
\end{equation}
Every coefficient of $C_j$ and $D_j$ is, up to sign, a minor of the same
Sylvester matrix used to define the subresultant, of order one smaller than
the corresponding subresultant determinant.  This determinantal form of the
subresultant B\'ezout identity is also developed in \cite{ElKahoui2003}.
\end{theorem}

\begin{proof}
Use the usual determinant matrix for
$\operatorname{Sres}_j(A,B)$, but leave its final column symbolic.  That
column consists of
\[
 x^{n-j-1}A,\ldots,A,
 x^{m-j-1}B,\ldots,B.
\]
Expanding the determinant along this column writes it as a linear combination
of these shifted copies of $A$ and $B$.  Collecting the coefficients of the
copies of $A$ gives a polynomial $C_j$ of degree at most $n-j-1$; collecting
the copies of $B$ gives $D_j$ of degree at most $m-j-1$.  The scalar
coefficients appearing in this Laplace expansion are cofactors of the final
column, hence signed minors of order one smaller than the subresultant
minor.  All matrix entries are integers, so $C_j,D_j\in\ZZ[x]$.
\end{proof}

For $j=0$, the left side is the resultant.  Thus, whenever $A$ and $B$ are
rationally coprime,
\begin{equation}\label{eq:resultant-bezout}
 \operatorname{Res}(A,B)=C_0A+D_0B
\end{equation}
for integral $C_0,D_0$.  This observation will be used below to compare the
terminal XGCD certificate with the resultant.

The same determinant proof supplies useful size information.

\begin{corollary}[Size of subresultant cofactors]
\label{cor:subresultant-cofactor-size}
Suppose $m,n\le N$ and every coefficient of $A,B$ has at most $\tau$ bits.
Then every coefficient of every $C_j,D_j$ in
\eqref{eq:subresultant-bezout} has
\begin{equation}\label{eq:subresultant-cofactor-bits}
 O\bigl(N(\tau+\log(N+1))\bigr)
\end{equation}
bits.
\end{corollary}

\begin{proof}
Each coefficient is a determinant of order at most $m+n-2j-1\le2N-1$.
Every row contains at most $N+1$ input coefficients of absolute value less
than $2^\tau$.  Hadamard's inequality therefore bounds the determinant by
\[
 \bigl(\sqrt{N+1}\,2^\tau\bigr)^{2N-1}.
\]
Taking binary logarithms gives the stated bound.
\end{proof}

\subsection{Propagating cofactors through a pseudo-division}

Suppose a PRS has reached consecutive members $F_{i-2},F_{i-1}$ and that
\begin{equation}\label{eq:xgcd-cofactor-invariant-old}
 F_r=U_rA+V_rB,
 \qquad r=i-2,i-1.
\end{equation}
Let
\[
 \alpha_i=\operatorname{lc}(F_{i-1})^{\deg F_{i-2}-\deg F_{i-1}+1}
\]
be the pseudo-division multiplier.  If $Q_i$ and $P_i$ are the
pseudo-quotient and pseudo-remainder, then
\begin{equation}\label{eq:xgcd-pseudodiv}
 P_i=\alpha_iF_{i-2}-Q_iF_{i-1}.
\end{equation}
Substitution of \eqref{eq:xgcd-cofactor-invariant-old} gives
\begin{align}
 \widetilde U_i&=\alpha_iU_{i-2}-Q_iU_{i-1},
 \label{eq:xgcd-raw-u}\\
 \widetilde V_i&=\alpha_iV_{i-2}-Q_iV_{i-1},
 \label{eq:xgcd-raw-v}
\end{align}
and hence
\begin{equation}\label{eq:xgcd-raw-certificate}
 P_i=\widetilde U_iA+\widetilde V_iB.
\end{equation}

For an arbitrary scalar normalization of $P_i$, integrality of the normalized
remainder does not by itself imply integrality of the corresponding
$\widetilde U_i,\widetilde V_i$.  Brown's factors have the stronger property
we need.

\begin{theorem}[Exact normalization of the extended subresultant PRS]
\label{thm:xgcd-exact-cofactor-normalization}
Let the remainders $F_i$ be Brown's subresultant PRS for $A,B$, using the
signed normalizing scalars $\beta_i$ from
\eqref{eq:prs-defining-relation}, so that
\[
 \beta_iF_i=P_i.
\]
Initialize
\[
 (U_1,V_1)=(1,0),
 \qquad
 (U_2,V_2)=(0,1).
\]
Then, at every subsequent step, every coefficient of
$\widetilde U_i$ and $\widetilde V_i$ in
\eqref{eq:xgcd-raw-u}--\eqref{eq:xgcd-raw-v} is divisible by $\beta_i$, and
\begin{equation}\label{eq:xgcd-normalized-cofactors}
 U_i=\widetilde U_i/\beta_i,
 \qquad
 V_i=\widetilde V_i/\beta_i
\end{equation}
lie in $\ZZ[x]$ and satisfy
\begin{equation}\label{eq:xgcd-cofactor-invariant}
 F_i=U_iA+V_iB.
\end{equation}
For $i\ge3$, these $U_i,V_i$ are the determinant cofactors belonging to the
same upper-boundary subresultant as $F_i$, up to the fixed sign convention.
\end{theorem}

\begin{proof}
The first step has only the unit sign normalization, so
\eqref{eq:xgcd-normalized-cofactors} is immediate.  For the general step,
consider the bordered Sylvester determinant used in the proof of
Theorem~\ref{thm:subresultant-bezout}.  Brown--Traub's fraction-free row
elimination, used in the proof of
Theorem~\ref{thm:brown-subresultant-prs-correct}, acts on the numerical
columns of this bordered matrix exactly as it acts on the ordinary
subresultant matrix.  The final symbolic column is subjected to the same row
operations, but is not otherwise special.

Consequently the similarity factor extracted from the determinant is the
same Brown factor $\beta_i$ simultaneously for three objects: the
subresultant itself, the coefficient multiplying $A$ in the expansion of the
bordered determinant, and the coefficient multiplying $B$.  Before this
factor is removed, those latter two objects are precisely
$\widetilde U_i$ and $\widetilde V_i$, because
\eqref{eq:xgcd-pseudodiv} applies the same fraction-free row operation to the
previous two B\'ezout identities.  Therefore $\beta_i$ divides every
coefficient of both polynomials exactly.  After division, the bordered
minors give the integral subresultant cofactors of
Theorem~\ref{thm:subresultant-bezout}.  Dividing
\eqref{eq:xgcd-raw-certificate} by $\beta_i$ then proves
\eqref{eq:xgcd-cofactor-invariant}.
\end{proof}

This theorem is stronger than exactness of Brown's remainder normalization.
It explains why the same division may safely be applied coefficientwise to
the two propagated cofactor polynomials.  This is the key fraction-free fact
behind integral XGCD.

If the $i$th PRS member is the upper subresultant of index
$j=n_{i-1}-1$, Theorem~\ref{thm:subresultant-bezout} also gives
\begin{equation}\label{eq:xgcd-cofactor-degree-bounds}
 \deg U_i\le n-n_{i-1},
 \qquad
 \deg V_i\le m-n_{i-1}.
\end{equation}
Thus the propagated cofactors always have $O(N)$ coefficients.

\subsection{The integral subresultant XGCD algorithm}

We can now state the algorithm in the form used by the companion
implementation.  Notice that it does not extract contents or primitive parts
from the inputs.

\begin{algorithm}[Integral subresultant XGCD]
\label{alg:subresultant-xgcd}
Input $A,B\in\ZZ[x]$.
\begin{enumerate}
 \item If one input is zero, return the other input together with the evident
 cofactor pair $(1,0)$ or $(0,1)$.
 \item If necessary interchange $A,B$ so that $\deg A\ge\deg B$, remembering
 to interchange the final cofactors back.
 \item Set
 \[
  F_1=A,\quad F_2=B,
  \qquad
  (U_1,V_1)=(1,0),\quad (U_2,V_2)=(0,1).
 \]
 \item At each Brown-PRS step compute the full pseudo-division
 \[
  P_i=\alpha_iF_{i-2}-Q_iF_{i-1}.
 \]
 If $P_i=0$, stop and return
 \[
  H=F_{i-1},\qquad U=U_{i-1},\qquad V=V_{i-1}.
 \]
 Otherwise compute
 \[
  \widetilde U_i=\alpha_iU_{i-2}-Q_iU_{i-1},
  \qquad
  \widetilde V_i=\alpha_iV_{i-2}-Q_iV_{i-1},
 \]
 and divide $P_i,\widetilde U_i,\widetilde V_i$ by the same signed Brown
 factor $\beta_i$.
 \item Continue with the next consecutive pair.
\end{enumerate}
\end{algorithm}

The quotient cannot be discarded here, in contrast with the gcd algorithms
of Section~\ref{sec:zz-poly-gcd}: it is needed explicitly in both cofactor
updates.

\begin{theorem}[Correctness of integral XGCD]
\label{thm:subresultant-xgcd-correct}
Algorithm~\ref{alg:subresultant-xgcd} terminates and returns integral
polynomials $H,U,V$ satisfying
\[
 H=UA+VB.
\]
If $G$ is the primitive polynomial with positive leading coefficient
associated to $\gcd_{\mathbb Q[x]}(A,B)$, then
\begin{equation}\label{eq:xgcd-h-is-multiple-gcd}
 H=cG
\end{equation}
for some nonzero integer $c$ whenever $A,B$ are both nonzero.
The integer $c$, and even the sign of $H$, depend on the ordered inputs and on
the PRS normalization; $H$ is not promised to be the canonical gcd in
$\ZZ[x]$.
\end{theorem}

\begin{proof}
Theorem~\ref{thm:xgcd-exact-cofactor-normalization} proves inductively that
every step is integral and preserves
$F_i=U_iA+V_iB$.  Degrees of nonzero PRS members strictly decrease, so the
algorithm terminates.  The returned identity is therefore
\eqref{eq:xgcd-cofactor-invariant} at the final nonzero member.

By Theorem~\ref{thm:last-subresultant-gcd}, that final member $H$ is associated
over $\mathbb Q[x]$ to the rational gcd.  Thus $H=qG$ for some
$q\in\mathbb Q^\times$.  Since $G$ is primitive and $H\in\ZZ[x]$, the same
content argument as in Gauss's lemma forces $q\in\ZZ$: write $q=a/b$ in
lowest terms; integrality of $(a/b)G$ implies that $b$ divides every
coefficient of $G$, hence $b=1$.  This proves
\eqref{eq:xgcd-h-is-multiple-gcd}.
\end{proof}

\begin{example}
For $A=2$ and $B=x$, the canonical gcd in $\ZZ[x]$ is $1$, but
Proposition~\ref{prop:zpoly-not-bezout} shows that no integral B\'ezout
certificate for $1$ exists.  Algorithm~\ref{alg:subresultant-xgcd} may return
\[
 H=2,
 \qquad U=1,
 \qquad V=0.
\]
This is a correct integral XGCD certificate, even though $H\ne\gcd(A,B)$.
\end{example}

\subsection{The certificate ideal and the resultant}

The output has a natural ideal-theoretic interpretation.  Let $G$ be the
primitive rational gcd from Theorem~\ref{thm:subresultant-xgcd-correct}, and
write
\[
 A=GA_0,
 \qquad
 B=GB_0,
\]
where $A_0,B_0\in\ZZ[x]$ by Gauss's divisibility lemma.  Define
\begin{equation}\label{eq:xgcd-certificate-ideal}
 \mathfrak c(A,B)
 =\{c\in\ZZ:cG\in(A,B)\}.
\end{equation}
This is an ideal of $\ZZ$, hence
\[
 \mathfrak c(A,B)=d\ZZ
\]
for a unique $d\ge0$.  When $A,B$ are nonzero, $d>0$ because the
subresultant XGCD itself supplies a nonzero element $c$ of this ideal.

\begin{proposition}\label{prop:xgcd-certificate-ideal}
If Algorithm~\ref{alg:subresultant-xgcd} returns $H=cG$, then
\[
 d\mid c.
\]
The multiplier $c$ need not equal $d$.  Thus the terminal Brown certificate
is an explicitly computable element of the certificate ideal, not in general
its least positive generator.
\end{proposition}

\begin{proof}
Equation~\eqref{eq:xgcd-certificate} shows that $cG\in(A,B)$, hence
$c\in\mathfrak c(A,B)=d\ZZ$.
\end{proof}

If $A$ and $B$ are rationally coprime, then $G=1$ and
\[
 \mathfrak c(A,B)=(A,B)\cap\ZZ.
\]
The resultant provides another canonical element of this same ideal.

\begin{corollary}[Resultant and the certificate ideal]
\label{cor:xgcd-resultant-ideal}
Suppose $\gcd_{\mathbb Q[x]}(A,B)=1$.  Then
\[
 0\ne\operatorname{Res}(A,B)\in(A,B)\cap\ZZ.
\]
If $(A,B)\cap\ZZ=d\ZZ$, then
\begin{equation}\label{eq:xgcd-d-divides-cert-resultant}
 d\mid H,
 \qquad
 d\mid\operatorname{Res}(A,B),
\end{equation}
where $H$ is the terminal integral XGCD certificate.
\end{corollary}

\begin{proof}
Rational coprimality makes the resultant nonzero.  The $j=0$ case of
Theorem~\ref{thm:subresultant-bezout} gives the integral identity
\eqref{eq:resultant-bezout}, so the resultant lies in $(A,B)\cap\ZZ$.
The XGCD identity places $H$ in the same ideal.  Since $d$ is its positive
generator, it divides both integers.
\end{proof}

Neither integer in \eqref{eq:xgcd-d-divides-cert-resultant} is asserted to be
$d$.  Determining the least positive generator of $(A,B)\cap\ZZ$ is an
integer-module problem rather than merely a PRS normalization problem.  It
can be approached by Hermite or Smith normal form of suitable coefficient
matrices; such infrastructure is deliberately outside the present
$\ZZ[x]$ layer.

\subsection{Complexity}

The extended algorithm performs the same PRS as the subresultant gcd, but it
also carries two cofactor sequences.  This does not change the quadratic
coefficient-operation order for the classical fraction-free implementation.

\begin{theorem}[Coefficient complexity of subresultant XGCD]
\label{thm:xgcd-coefficient-complexity}
Let $A,B$ have degrees at most $N$.  Using classical pseudo-division and
classical polynomial products in the cofactor updates,
Algorithm~\ref{alg:subresultant-xgcd} uses $O(N^2)$ coefficient additions,
multiplications and exact scalar divisions.
\end{theorem}

\begin{proof}
The underlying Brown PRS costs $O(N^2)$ coefficient operations by
Theorem~\ref{thm:subresultant-gcd-coefficient-complexity}.  At a step whose
degree drop is $\delta_i$, the pseudo-quotient has $O(\delta_i+1)$
coefficients.  By \eqref{eq:xgcd-cofactor-degree-bounds}, each previous
cofactor has $O(N)$ coefficients.  Thus multiplying the pseudo-quotient by
each of the two cofactors costs
\[
 O\bigl(N(\delta_i+1)\bigr)
\]
classical coefficient operations.  The degree drops sum to at most $N$, and
there are at most $N$ steps, so summing these products gives $O(N^2)$.
Scaling, subtraction and Brown normalization each touch only $O(N)$
cofactor coefficients per step and therefore also total $O(N^2)$.
\end{proof}

The normalized cofactors have the same determinant-size bit bound as the
subresultants themselves.

\begin{theorem}[Coefficient growth in subresultant XGCD]
\label{thm:xgcd-coefficient-growth}
If the input degrees are at most $N$ and the input coefficient bit sizes are
at most $\tau$, every coefficient of every normalized PRS member and of its
normalized B\'ezout cofactors has
\begin{equation}\label{eq:xgcd-normalized-bit-bound}
 W=O\bigl(N(\tau+\log(N+1))\bigr)
\end{equation}
bits.
\end{theorem}

\begin{proof}
The bound for PRS members is
Theorem~\ref{thm:subresultant-coefficient-growth}.  The bound for cofactors is
Corollary~\ref{cor:subresultant-cofactor-size}.  By
Theorem~\ref{thm:xgcd-exact-cofactor-normalization}, the propagated normalized
cofactors are precisely these determinant cofactors, up to the fixed signs.
\end{proof}

As in the preceding sections, a straightforward fraction-free implementation
materializes larger values before the exact Brown divisions.  Let
\[
 B_*=O\bigl(N(W+1)\bigr).
\]
The pseudo-quotients and raw pseudo-remainders have $O(B_*)$-bit coefficients
by the same pseudo-division growth estimate used in
Theorem~\ref{thm:subresultant-gcd-bit-complexity}.  The raw cofactor
numerators
\[
 \alpha_iU_{i-2}-Q_iU_{i-1},
 \qquad
 \alpha_iV_{i-2}-Q_iV_{i-1}
\]
also have $O(B_*)$-bit coefficients: multiplication by $\alpha_i$ adds at
most $O(NW)$ bits, while a coefficient of a quotient--cofactor product is a
sum of at most $N$ products of an $O(B_*)$-bit quotient coefficient and an
$O(W)$-bit cofactor coefficient.

\begin{corollary}[Direct bit bound for subresultant XGCD]
\label{cor:xgcd-bit-complexity}
With $W$ and $B_*$ as above, and with $\mathsf D_{\ZZ}(t)$ denoting exact
integer division cost, a direct classical implementation has the worst-case
bit bound
\begin{equation}\label{eq:xgcd-bit-complexity}
 O\!\left(
   N^2\bigl(\MZ(B_*)+\mathsf D_{\ZZ}(B_*)\bigr)
 \right).
\end{equation}
\end{corollary}

\begin{proof}
Theorem~\ref{thm:xgcd-coefficient-complexity} gives $O(N^2)$ coefficient
operations.  The preceding discussion bounds every transient integer charged
to a multiplication or exact Brown division by $O(B_*)$ bits.  Additions are
asymptotically no more expensive than multiplication under the standing
assumption $\MZ(t)=\Omega(t)$.
\end{proof}

The bound is deliberately for the simple fraction-free algorithm.  Faster
extended gcd algorithms use half-GCD transformations to reduce the number of
polynomial operations; that is the subject of the next section of these
notes.


\section{Fraction-free half-GCD}
\label{sec:zz-poly-hgcd}

The classical pseudo-Euclidean algorithms of the preceding sections discover
one quotient at a time.  Half-GCD (HGCD) algorithms instead use the high
coefficients of the current pair to predict a whole block of Euclidean
transformations, apply that block to the complete operands, perform one
crossing division, and recurse.  This divide-and-conquer organization goes
back, for polynomial gcds, to Moenck \cite{Moenck1973}; a particularly clean
correctness framework is due to Thull and Yap \cite{ThullYap1990}.  Modern
field versions are also described in \cite{vonZurGathenGerhard2013}.

The companion implementation uses the same degree divide-and-conquer idea but
keeps all polynomial coefficients integral.  Ordinary Euclidean divisions are
replaced by pseudo-divisions and the transformation matrices lie in
$\ZZ[x]^{2\times2}$.  This makes the algebraic degree complexity fast without
introducing $\mathbb Q[x]$, but coefficient control is more delicate than for
the Brown-normalized subresultant algorithms of
Sections~\ref{sec:subresultant-prs} and~\ref{sec:zz-poly-gcd}.

\subsection{Pseudo-Euclidean transformation matrices}

Let $C,D\in\ZZ[x]$ be nonzero with $\deg C\ge\deg D$.  Put
\[
 \delta=\deg C-\deg D+1,
 \qquad
 \beta=\operatorname{lc}(D).
\]
If pseudo-division gives
\begin{equation}\label{eq:hgcd-pseudodiv}
 \beta^\delta C=QD+R,
 \qquad \deg R<\deg D,
\end{equation}
then one pseudo-Euclidean step may be written
\begin{equation}\label{eq:hgcd-step-matrix}
 \begin{pmatrix}D\\R\end{pmatrix}
 =
 E(C,D)
 \begin{pmatrix}C\\D\end{pmatrix},
 \qquad
 E(C,D)=
 \begin{pmatrix}
  0&1\\
  \beta^\delta&-Q
 \end{pmatrix}.
\end{equation}
The determinant is
\[
 \det E(C,D)=-\beta^\delta\in\ZZ\setminus\{0\}.
\]
Thus $E(C,D)$ is invertible over $\mathbb Q[x]$ even though it is not in
general unimodular over $\ZZ[x]$.

\begin{definition}
An integral polynomial matrix
$M\in\ZZ[x]^{2\times2}$ will be called \emph{rationally unimodular} if
\[
 \det M\in\mathbb Q^*.
\]
Since $\det M\in\ZZ[x]$, this simply means that $\det M$ is a nonzero integer
constant.  Such a matrix is an element of $\operatorname{GL}_2(\mathbb Q[x])$.
\end{definition}

\begin{lemma}\label{lem:hgcd-matrix-preserves-gcd}
If $M\in\ZZ[x]^{2\times2}$ is rationally unimodular and
\[
 \begin{pmatrix}C\\D\end{pmatrix}
 =M\begin{pmatrix}A\\B\end{pmatrix},
\]
then $(A,B)$ and $(C,D)$ generate the same ideal of $\mathbb Q[x]$.  In
particular they have associated gcds in $\mathbb Q[x]$.
\end{lemma}

\begin{proof}
The inclusion $(C,D)\subseteq(A,B)$ over $\mathbb Q[x]$ follows from the
matrix equation.  Since $\det M\in\mathbb Q^*$, the inverse matrix
$M^{-1}$ also has entries in $\mathbb Q[x]$.  Applying it gives the reverse
inclusion.  The two ideals, and hence their monic generators, are equal.
\end{proof}

Every pseudo-Euclidean step matrix in \eqref{eq:hgcd-step-matrix} is therefore
rationally unimodular, and so is every product of such matrices.

\subsection{Primitive normalization of a transformation matrix}

Pseudo-division introduces large scalar powers of leading coefficients.  A
fraction-free matrix product consequently tends to acquire large common
integer factors.  The companion implementation removes these factors after
matrix compositions.

For
\[
 M=(m_{ij})\in\ZZ[x]^{2\times2},
\]
define its integer content by
\[
 \operatorname{cont}(M)
 =\gcd\{[x^k]m_{ij}:0\le i,j<2,\ k\ge0\},
\]
with positive gcd.  If this content is greater than $1$, replace $M$ by
\[
 \operatorname{pp}(M)=M/\operatorname{cont}(M).
\]

\begin{lemma}\label{lem:hgcd-matrix-content}
Let $M\in\ZZ[x]^{2\times2}$ be rationally unimodular, let
$c=\operatorname{cont}(M)>0$, and put $M'=M/c$.  Then:
\begin{enumerate}
 \item $M'\in\ZZ[x]^{2\times2}$ and is rationally unimodular;
 \item for every pair $A,B$,
 \[
  M'\begin{pmatrix}A\\B\end{pmatrix}
  =\frac1c
   M\begin{pmatrix}A\\B\end{pmatrix};
 \]
 hence the two output degrees are unchanged;
 \item $M'$ has no nonunit integer dividing every coefficient of all four
 entries.
\end{enumerate}
\end{lemma}

\begin{proof}
The definition of $c$ proves integrality and the final assertion.  Moreover
\[
 \det M'=\frac{\det M}{c^2}\in\mathbb Q^*,
\]
so $M'$ is invertible over $\mathbb Q[x]$.  The displayed action on a vector
is immediate, and multiplication by a nonzero rational scalar does not change
polynomial degrees.
\end{proof}

This normalization is mathematically harmless for degree reduction and gcd
computation.  It is not Brown's subresultant normalization: it removes only a
scalar common to \emph{all four} matrix entries, and no theorem here identifies
the normalized entries with canonical subresultant minors.

A related observation will be used below.  If a matrix $U$ has a desired
degree-reduction property on $(cA,cB)$ for some nonzero integer $c$, then it
has exactly the same degree-reduction property on $(A,B)$, because
\[
 U\begin{pmatrix}cA\\cB\end{pmatrix}
 =cU\begin{pmatrix}A\\B\end{pmatrix}.
\]
Thus a common scalar may be present or removed when choosing the next recursive
high-part transformation.

\subsection{High parts and the half-degree principle}

For $k\ge0$, write
\[
 \operatorname{hi}_k(F)=\left\lfloor\frac{F}{x^k}\right\rfloor,
\]
so that
\[
 F=F_0+x^k\operatorname{hi}_k(F),
 \qquad \deg F_0<k.
\]
For a polynomial matrix $M=(m_{ij})$, write
\[
 \deg M=\max_{i,j}\deg m_{ij},
\]
with the zero polynomial assigned degree $-\infty$.

\begin{lemma}[High-part transfer]\label{lem:hgcd-high-part-transfer}
Let $M\in\ZZ[x]^{2\times2}$ satisfy $\deg M\le s$.  Write
\[
 A=A_0+x^kA_1,
 \qquad
 B=B_0+x^kB_1,
 \qquad
 \deg A_0,\deg B_0<k.
\]
Then
\begin{equation}\label{eq:hgcd-high-part-transfer}
 M\begin{pmatrix}A\\B\end{pmatrix}
 =x^kM\begin{pmatrix}A_1\\B_1\end{pmatrix}
  +M\begin{pmatrix}A_0\\B_0\end{pmatrix},
\end{equation}
and every component of the second vector on the right has degree strictly
less than $k+s$.
\end{lemma}

\begin{proof}
The identity is linearity.  Every term in a component of
$M(A_0,B_0)^T$ is a product of a polynomial of degree at most $s$ with one of
degree less than $k$, and hence has degree less than $k+s$.
\end{proof}

The recursive algorithm also maintains a degree bound on its transformation
matrix.  This is the technical fact that makes high-half prediction rigorous.

\begin{definition}[Half-GCD contract]\label{def:hgcd-contract}
Let $A,B\in\ZZ[x]$ with
\[
 m=\deg A>\deg B\ge0,
 \qquad
 h=\left\lceil\frac m2\right\rceil.
\]
A matrix $M$ satisfies the HGCD contract for $(A,B)$ if it is integral and
rationally unimodular and, writing
\[
 \begin{pmatrix}C\\D\end{pmatrix}
 =M\begin{pmatrix}A\\B\end{pmatrix},
\]
one has
\begin{equation}\label{eq:hgcd-contract-degrees}
 \deg C\ge h>\deg D
\end{equation}
and
\begin{equation}\label{eq:hgcd-contract-matrix-degree}
 \deg M\le m-\deg C.
\end{equation}
The convention $\deg0=-\infty$ is used in
\eqref{eq:hgcd-contract-degrees}.
\end{definition}

The last inequality in particular gives
\[
 \deg M\le m-h=\left\lfloor\frac m2\right\rfloor.
\]

\subsection{The fraction-free half-GCD recursion}

The following is the mathematical form of the recursion used by the companion
implementation.  Matrix primitive-part operations mean the normalization of
Lemma~\ref{lem:hgcd-matrix-content}.

\begin{algorithm}[Fraction-free half-GCD]\label{alg:hgcd-pseudo}
Input nonzero $A,B\in\ZZ[x]$ with $m=\deg A>\deg B$.
\begin{enumerate}
 \item Put $h\leftarrow\lceil m/2\rceil$.  If $\deg B<h$, return the identity
 matrix.
 \item Set
 \[
  A_1\leftarrow\operatorname{hi}_h(A),
  \qquad
  B_1\leftarrow\operatorname{hi}_h(B),
 \]
 and recursively compute
 \[
  R\leftarrow\operatorname{HGCD}(A_1,B_1).
 \]
 \item Apply $R$ to the complete pair:
 \[
  \begin{pmatrix}C\\D\end{pmatrix}
  \leftarrow R\begin{pmatrix}A\\B\end{pmatrix}.
 \]
 If $D=0$ or $\deg D<h$, return $R$.
 \item Pseudo-divide $C$ by $D$:
 \[
  \operatorname{lc}(D)^\delta C=QD+E,
  \qquad
  \delta=\deg C-\deg D+1.
 \]
 Form
 \[
  T=
  \begin{pmatrix}
   0&1\\
   \operatorname{lc}(D)^\delta&-Q
  \end{pmatrix},
  \qquad
  S\leftarrow\operatorname{pp}(TR).
 \]
 If $E=0$ or $\deg E<h$, return $S$.
 \item Put
 \[
  k\leftarrow2h-\deg D,
 \]
 and recursively compute
 \[
  U\leftarrow
  \operatorname{HGCD}
  \bigl(\operatorname{hi}_k(D),\operatorname{hi}_k(E)\bigr).
 \]
 \item Return
 \[
  \operatorname{pp}(US).
 \]
\end{enumerate}
\end{algorithm}

In Step 5 the companion implementation uses the unnormalized pair $(D,E)$
from Step 4 even though $S$ may have had common matrix content removed.  This
is equivalent for degree purposes: $S(A,B)^T$ is a common nonzero scalar
multiple of $(D,E)^T$, and the observation following
Lemma~\ref{lem:hgcd-matrix-content} applies.

\begin{theorem}[Correctness of fraction-free HGCD]
\label{thm:hgcd-pseudo-correctness}
Algorithm~\ref{alg:hgcd-pseudo} terminates and returns a matrix satisfying the
HGCD contract of Definition~\ref{def:hgcd-contract}.
\end{theorem}

\begin{proof}
We use induction on $m=\deg A$.  Step 1 is immediate: the identity is
rationally unimodular, its matrix degree is zero, and
$\deg A\ge h>\deg B$.

Assume now $\deg B\ge h$.  Put
\[
 m_1=m-h,
 \qquad
 h_1=\left\lceil\frac{m_1}{2}\right\rceil.
\]
Then $\deg A_1=m_1$.  By induction, if
\[
 \begin{pmatrix}C_1\\D_1\end{pmatrix}
 =R\begin{pmatrix}A_1\\B_1\end{pmatrix},
\]
then
\[
 \deg C_1\ge h_1>\deg D_1,
 \qquad
 \deg R\le m_1-\deg C_1.
\]
Write $A=A_0+x^hA_1$ and $B=B_0+x^hB_1$.  By
Lemma~\ref{lem:hgcd-high-part-transfer}, applying $R$ to the full pair gives
\[
 C=x^hC_1+C_0,
 \qquad
 D=x^hD_1+D_0,
\]
where
\[
 \deg C_0,\deg D_0
 <h+m_1-\deg C_1.
\]
Since $2\deg C_1\ge m_1$, the low perturbation $C_0$ cannot cancel the
leading term of $x^hC_1$.  Hence
\begin{equation}\label{eq:hgcd-first-output-degree}
 \deg C=h+\deg C_1.
\end{equation}
Moreover
\[
 \deg D<h+h_1,
\]
because $\deg D_1<h_1$ and
$m_1-\deg C_1\le m_1-h_1\le h_1$.  Also
\[
 \deg R\le m_1-\deg C_1=m-\deg C.
\]
Thus if Step 3 returns, the complete HGCD contract holds.

Suppose instead that $d=\deg D\ge h$.  Equation
\eqref{eq:hgcd-first-output-degree} and the bound $\deg D<h+h_1$ show in
particular that $\deg C>d$, so Step 4 is a proper pseudo-Euclidean step.  The
matrix $T$ has nonzero constant determinant and hence is rationally
unimodular.  Before primitive normalization it maps $(A,B)$ to $(D,E)$.
Furthermore the pseudo-quotient has degree
\[
 \deg Q=\deg C-d.
\]
Using $\deg R\le m-\deg C$ gives
\[
 \deg(TR)
 \le \deg Q+\deg R
 \le (\deg C-d)+(m-\deg C)
 =m-d.
\]
Primitive matrix normalization cannot increase degrees, so
\begin{equation}\label{eq:hgcd-crossing-matrix-bound}
 \deg S\le m-d.
\end{equation}
If $E=0$ or $\deg E<h$, Step 4 therefore returns a matrix satisfying the
contract, with first transformed component of degree $d\ge h$.

It remains to consider Step 5.  We have
\[
 h\le d<h+h_1\le2h,
\]
so
\[
 k=2h-d
\]
is positive.  Put
\[
 r=d-h.
\]
Then
\begin{equation}\label{eq:hgcd-second-high-degree}
 \deg\operatorname{hi}_k(D)
 =d-k=2(d-h)=2r.
\end{equation}
Since this is even, the half-degree threshold in the second recursive call is
exactly $r$.  By induction, if that call returns $U$ and its first transformed
high component has degree $p$, then
\[
 p\ge r,
 \qquad
 \deg U\le2r-p,
\]
while its second transformed high component has degree less than $r$.
Applying $U$ to the complete pair $(D,E)$ and using
Lemma~\ref{lem:hgcd-high-part-transfer} at the shift $k$, the first component
therefore has degree exactly $k+p\ge k+r=h$, and the second has degree less
than $k+r=h$.  The same statement holds if $(D,E)$ is first divided by the
common scalar removed from $TR$.

Finally,
\[
 \deg U\le2r-p=d-(k+p).
\]
Combining this with \eqref{eq:hgcd-crossing-matrix-bound} gives
\[
 \deg(US)
 \le d-(k+p)+m-d
 =m-(k+p),
\]
which is precisely the matrix-degree part of the HGCD contract for the final
first component.  Primitive normalization preserves this inequality and
rational unimodularity.  Both recursive calls have first input degree strictly
smaller than $m$, so induction also proves termination.
\end{proof}

The proof shows why the somewhat unusual shift
$k=2h-\deg D$ is natural: it makes the degree of the second recursive first
operand exactly twice the remaining degree drop $\deg D-h$.

\subsection{Recovering the gcd}

Half-GCD is a block-reduction primitive.  To obtain a complete gcd, repeatedly
apply it and perform one crossing pseudo-division between blocks.

\begin{algorithm}[Fraction-free HGCD gcd]\label{alg:gcd-hgcd}
Input $A,B\in\ZZ[x]$.
\begin{enumerate}
 \item Compute
 \[
  c=\gcd(\operatorname{cont}(A),\operatorname{cont}(B)),
 \]
 and replace $A,B$ by their primitive parts.  Swap them if necessary so that
 $\deg A\ge\deg B$.
 \item While $B\ne0$:
 \begin{enumerate}
  \item for a small base case, or when $\deg A=\deg B$, perform one direct
  pseudo-division step;
  \item otherwise compute $M\leftarrow\operatorname{HGCD}(A,B)$ and replace
  $(A,B)^T$ by $M(A,B)^T$;
  \item if the new $B$ is nonzero, perform one crossing pseudo-division step.
 \end{enumerate}
 \item Replace the final nonzero $A$ by its primitive part, choose positive
 leading coefficient, multiply by $c$, and return it.
\end{enumerate}
\end{algorithm}

\begin{theorem}\label{thm:gcd-hgcd-correctness}
Algorithm~\ref{alg:gcd-hgcd} returns the canonical gcd in $\ZZ[x]$.
Whenever an HGCD block is used on a first operand of degree $m$, the degree of
the first operand at the beginning of the next outer iteration is strictly
less than $\lceil m/2\rceil$ after the crossing step.
\end{theorem}

\begin{proof}
Content removal and restoration are correct by
Theorem~\ref{thm:gcd-content-primitive}.  Every HGCD matrix is rationally
unimodular by Theorem~\ref{thm:hgcd-pseudo-correctness}, so
Lemma~\ref{lem:hgcd-matrix-preserves-gcd} shows that an HGCD block preserves
the rational gcd.  A pseudo-division step also preserves the rational gcd,
as proved in Lemma~\ref{lem:pseudorem-preserves-gcd}.  Hence every loop
iteration preserves the primitive rational gcd.

After an HGCD block, the second transformed component has degree less than
$h=\lceil m/2\rceil$.  The following crossing step makes that polynomial the
new first component.  Thus the first degree at the next outer iteration is
less than $h$.  Repeated blocks therefore terminate.  At termination the last
nonzero polynomial is a rational associate of the primitive gcd; taking its
primitive part, fixing the sign, and restoring the gcd of the input contents
gives the canonical gcd in $\ZZ[x]$.
\end{proof}

\subsection{Coefficient-operation complexity}

Let $\mathsf M(n)$ denote the coefficient-operation cost of a full product of
polynomials of length $O(n)$, and let $\mathsf P(n)$ denote the cost of a fast
pseudo-division on operands of degree $O(n)$.  Let $H(n)$ be the cost of one
HGCD call whose first input has degree at most $n$.

Both recursive calls in Algorithm~\ref{alg:hgcd-pseudo} have first input
degree at most $\lceil n/2\rceil$.  Outside those calls, the algorithm applies
and composes $2\times2$ polynomial matrices of degree $O(n)$, performs one
pseudo-division, and makes only linear additional polynomial additions and
scalar operations.  Matrix primitive normalization scans $O(n)$ integer
coefficients and performs $O(n)$ integer gcds and exact scalar divisions.
Thus, at the level of polynomial coefficient operations,
\begin{equation}\label{eq:hgcd-recurrence}
 H(n)
 =2H(\lceil n/2\rceil)
  +O\bigl(\mathsf M(n)+\mathsf P(n)+n\bigr),
\end{equation}
with the integer gcd work recorded separately.

\begin{theorem}\label{thm:hgcd-coefficient-complexity}
Assume $\mathsf P(n)=O(\mathsf M(n))$.
Then
\begin{equation}\label{eq:hgcd-sum-complexity}
 H(n)
 =O\!\left(
   \sum_{j=0}^{\lceil\log_2n\rceil}
   2^j\mathsf M\!\left(\left\lceil\frac{n}{2^j}\right\rceil\right)
  \right).
\end{equation}
In particular:
\begin{enumerate}
 \item if $\mathsf M(n)=\Theta(n^\alpha)$ with $\alpha>1$, then
 $H(n)=O(\mathsf M(n))$;
 \item if $\mathsf M(n)=O(n\log^\kappa n)$ for fixed $\kappa\ge0$, then
 \[
  H(n)=O(n\log^{\kappa+1}n)
      =O(\mathsf M(n)\log n).
 \]
\end{enumerate}
The complete gcd driver has the same asymptotic coefficient-operation bound,
because successive HGCD blocks reduce the active degree geometrically.
\end{theorem}

\begin{proof}
Unfolding \eqref{eq:hgcd-recurrence} by recursion level gives
\eqref{eq:hgcd-sum-complexity}.  If
$\mathsf M(n)=\Theta(n^\alpha)$ with $\alpha>1$, the contribution of level
$j$ is
\[
 2^j\Theta\bigl((n/2^j)^\alpha\bigr)
 =\Theta\bigl(n^\alpha2^{-j(\alpha-1)}\bigr),
\]
a convergent geometric series.  For
$\mathsf M(n)=O(n\log^\kappa n)$, every nonempty level costs
$O(n\log^\kappa n)$ and there are $O(\log n)$ levels.

By Theorem~\ref{thm:gcd-hgcd-correctness}, the sizes of successive outer
HGCD blocks are $n,n/2,n/4,\ldots$ up to rounding.  Summing the displayed
HGCD bounds over these geometrically decreasing sizes changes only the
constant in the power-growth case and preserves the stated
$O(n\log^{\kappa+1}n)$ bound in the quasi-linear case.
\end{proof}

For the multiplication algorithms already developed in this chapter,
Karatsuba and fixed-order Toom multiplication have exponent $\alpha>1$.
Thus fraction-free HGCD preserves their degree exponent.  With an
asymptotically quasi-linear multiplication algorithm, the usual extra
logarithmic HGCD factor appears.

\subsection{Coefficient growth and bit complexity}

The preceding theorem counts polynomial operations; over $\ZZ[x]$ it is not by
itself a bit-complexity theorem.  The reason is the same expression-swell
phenomenon that motivated subresultant normalization, but here it also affects
the transformation matrices.

The primitive matrix step is locally optimal among \emph{common scalar}
normalizations: after replacing $M$ by $M/\operatorname{cont}(M)$, no larger
nonunit integer can be divided simultaneously from all four entries while
remaining in $\ZZ[x]^{2\times2}$.  This can remove very large predictable
powers introduced by pseudo-division.  It does not, however, give the
canonical determinantal bounds of Theorem~\ref{thm:subresultant-coefficient-growth}
or Corollary~\ref{cor:subresultant-cofactor-size}.

To make the remaining cost explicit, let $W$ be a bound for the bit size of
all coefficients occurring at some size-$n$ HGCD node.  If
$\mathcal M(n,W)$ bounds a polynomial product at that size,
$\mathcal P(n,W)$ bounds fast pseudo-division, $\mathsf G_{\ZZ}(W)$ bounds an
integer gcd and $\mathsf D_{\ZZ}(W)$ an exact integer division, then the
nonrecursive bit cost at that node is bounded by
\begin{equation}\label{eq:hgcd-conditional-bit-cost}
 O\!\left(
   \mathcal M(n,W)+\mathcal P(n,W)
   +n\bigl(\mathsf G_{\ZZ}(W)+\mathsf D_{\ZZ}(W)\bigr)
  \right).
\end{equation}
The last term is the matrix-content computation and the four coefficientwise
exact divisions by the common content.

What is missing is a polynomial a priori bound for $W$ for this particular
normalization.  The single-step estimate of
Lemma~\ref{lem:pseudodiv-coefficient-growth} shows why this matters: a
pseudo-division with degree difference $O(n)$ can turn $w$-bit coefficients
into coefficients having $O(n(w+1))$ bits.  If such estimates are simply
iterated along a pseudo-remainder chain and no cancellation or common-content
removal is credited, only an exponential-in-degree coefficient bound results.
The matrix-content divisions observed in practice remove much of this growth,
but the present theory does not identify their outputs with bounded-size
subresultant minors and therefore does not turn
\eqref{eq:hgcd-conditional-bit-cost} into a useful quasi-linear bit bound.

\begin{remark}\label{rem:hgcd-experimental-status}
The fraction-free HGCD is therefore asymptotically attractive in the
\emph{polynomial degree} model and is a valid gcd algorithm, but its current
common-content normalization is weaker than the coefficient-control theory
available for fast subresultant algorithms.  This is why the companion
implementation exposes it as a separate reference/experimental gcd rather
than selecting it as the default over $\ZZ[x]$.  Brown's subresultant PRS
remains the path with a proved determinant-size coefficient bound.
\end{remark}



\section{Resultants over $\ZZ[x]$}\label{sec:resultants}

The resultant packages the question whether two polynomials have a common
root into a single scalar.  It also forms the zeroth member of the
subresultant family studied in Section~\ref{sec:subresultant-prs}.  This makes
Brown's polynomial remainder sequence a natural fraction-free resultant
algorithm: the same sequence used for gcd computation already contains the
resultant as its final lower-subresultant scalar.

Throughout this section let
\[
 A(x)=a_mx^m+\cdots+a_0,
 \qquad
 B(x)=b_nx^n+\cdots+b_0
\]
be nonzero polynomials over a field or over $\ZZ$, with $a_m b_n\ne0$.
The classical determinant and product formulas below are standard; see, for
example, \cite{Sturmfels2002,vonZurGathenGerhard2013}.  When one input is
the zero polynomial, whose degree is undefined, we adopt the companion
implementation's convention that the resultant is zero.

\subsection{The Sylvester determinant and product formula}

\begin{definition}[Sylvester matrix and resultant]
The \emph{Sylvester matrix} $\operatorname{Syl}(A,B)$ is the
$(m+n)\times(m+n)$ coefficient matrix whose first $n$ rows are the
coefficient vectors of
\[
 x^{n-1}A,\ldots,xA,A
\]
and whose final $m$ rows are the coefficient vectors of
\[
 x^{m-1}B,\ldots,xB,B,
\]
all written in the same descending monomial basis.  The resultant is
\begin{equation}\label{eq:resultant-sylvester}
 \operatorname{Res}(A,B)=\det\operatorname{Syl}(A,B).
\end{equation}
For two nonzero constants the matrix is empty and its determinant is $1$.
\end{definition}

The determinant has a useful linear-map interpretation.  If $K$ is a field,
then, up to the harmless choice of ordering of the two domain blocks,
the Sylvester matrix represents
\[
 \Phi:K[x]_{<n}\oplus K[x]_{<m}\longrightarrow K[x]_{<m+n},
 \qquad
 (U,V)\longmapsto UA+VB.
\]
Thus the determinant vanishes precisely when the two bounded families of
shifts are linearly dependent.

\begin{theorem}[Product formula]\label{thm:resultant-product-formula}
Let $L$ be a splitting field and write
\[
 A=a_m\prod_{i=1}^{m}(x-\alpha_i),
 \qquad
 B=b_n\prod_{j=1}^{n}(x-\beta_j),
\]
with roots repeated according to multiplicity.  Then
\begin{align}
 \operatorname{Res}(A,B)
 &=a_m^n\prod_{i=1}^{m}B(\alpha_i),
 \label{eq:resultant-product-a}\\
 &=a_m^n b_n^m
   \prod_{i=1}^{m}\prod_{j=1}^{n}(\alpha_i-\beta_j),
 \label{eq:resultant-product-roots}\\
 &=(-1)^{mn}b_n^m\prod_{j=1}^{n}A(\beta_j).
 \label{eq:resultant-product-b}
\end{align}
These expressions equal the Sylvester determinant
\eqref{eq:resultant-sylvester}.
\end{theorem}

\begin{proof}
Work first over the splitting field $L$.  In the quotient algebra
$L[x]/(A/a_m)$, multiplication by $B$ has determinant
\[
 \prod_{i=1}^{m}B(\alpha_i).
\]
Indeed, after extending scalars if necessary, multiplication by $x$ has the
roots $\alpha_i$ as eigenvalues, counted with multiplicity; applying $B$ to
that operator gives eigenvalues $B(\alpha_i)$.  Equivalently, this follows by
putting the companion matrix of the monic polynomial $A/a_m$ into triangular
form.

Now order the codomain basis as
\[
 x^{n-1}A,\ldots,xA,A,x^{m-1},\ldots,x,1.
\]
Relative to the descending monomial basis, its change-of-basis matrix is
block triangular with determinant $a_m^n$.  Order the first domain block in
the same way.  In this adapted basis the map $\Phi$ is block triangular: the
first diagonal block is the identity on the multiples of $A$, while the second
is multiplication by $B$ modulo $A$.  Hence
\[
 \det\operatorname{Syl}(A,B)
 =a_m^n\prod_iB(\alpha_i),
\]
which proves \eqref{eq:resultant-product-a}.  Factoring $B$ gives
\eqref{eq:resultant-product-roots}.  Interchanging the two roots in every
factor introduces $(-1)^{mn}$ and gives
\eqref{eq:resultant-product-b}.
\end{proof}

Several standard laws follow immediately.

\begin{proposition}[Basic resultant identities]\label{prop:resultant-basic}
For nonzero polynomials of degrees $m,n$ and nonzero scalars $c,d$,
\begin{align}
 \operatorname{Res}(B,A)
   &=(-1)^{mn}\operatorname{Res}(A,B),
 \label{eq:resultant-symmetry}\\
 \operatorname{Res}(cA,dB)
   &=c^n d^m\operatorname{Res}(A,B),
 \label{eq:resultant-scaling}\\
 \operatorname{Res}(A_1A_2,B)
   &=\operatorname{Res}(A_1,B)\operatorname{Res}(A_2,B),
 \label{eq:resultant-multiplicative-a}\\
 \operatorname{Res}(A,B_1B_2)
   &=\operatorname{Res}(A,B_1)\operatorname{Res}(A,B_2).
 \label{eq:resultant-multiplicative-b}
\end{align}
In particular, for a nonzero constant $c$,
\begin{equation}\label{eq:resultant-constant}
 \operatorname{Res}(c,B)=c^n,
 \qquad
 \operatorname{Res}(A,c)=c^m.
\end{equation}
\end{proposition}

\begin{proof}
Equation \eqref{eq:resultant-symmetry} follows from exchanging
$\alpha_i-\beta_j$ with $\beta_j-\alpha_i$ in
\eqref{eq:resultant-product-roots}.  Scaling either input changes only its
leading coefficient and gives \eqref{eq:resultant-scaling}.  The root
multisets of a product are the disjoint union of the root multisets of its
factors, which proves multiplicativity.  The constant formulas are the
special cases with one degree equal to zero.
\end{proof}

\begin{theorem}[Common-root criterion]\label{thm:resultant-zero-criterion}
For nonzero $A,B\in K[x]$,
\[
 \operatorname{Res}(A,B)=0
 \quad\Longleftrightarrow\quad
 \gcd(A,B)\ne1
 \quad\Longleftrightarrow\quad
 A\text{ and }B\text{ have a common root in }\overline K.
\]
\end{theorem}

\begin{proof}
By \eqref{eq:resultant-product-roots}, the resultant vanishes exactly when
$\alpha_i=\beta_j$ for some pair of roots.  This is equivalent to a common
linear factor over the algebraic closure, hence to a positive-degree gcd over
$K[x]$.
\end{proof}

Over $\ZZ[x]$, the determinant formula also shows directly that the resultant
is an integer.  The subresultant B\'ezout identity from
Theorem~\ref{thm:subresultant-bezout}, specialized to $j=0$, gives
\begin{equation}\label{eq:resultant-bezout-repeat}
 \operatorname{Res}(A,B)=UA+VB,
 \qquad U,V\in\ZZ[x],
\end{equation}
with
\[
 \deg U<n,
 \qquad
 \deg V<m.
\]
When the resultant is nonzero, this is the integral certificate used in
Corollary~\ref{cor:xgcd-resultant-ideal}.

\subsection{Euclidean and pseudo-remainder recurrences}

The product formula gives a useful recurrence under polynomial division.
Unlike the integer Euclidean algorithm, a change in degree contributes a
power of the divisor's leading coefficient.

\begin{lemma}[Resultant under ordinary remainder]\label{lem:resultant-remainder}
Let $K$ be a field, let $m\ge n>0$, and write
\[
 A=QB+R,
 \qquad
 r=\deg R<n,
\]
with $R\ne0$.  Then
\begin{equation}\label{eq:resultant-remainder}
 \operatorname{Res}(A,B)
 =(-1)^{mn}b_n^{m-r}\operatorname{Res}(B,R).
\end{equation}
If $R=0$, both sides are interpreted as zero.
\end{lemma}

\begin{proof}
At every root $\beta_j$ of $B$ one has $A(\beta_j)=R(\beta_j)$.  Applying
\eqref{eq:resultant-product-b} to $(A,B)$ and to $(R,B)$ gives
\[
 \operatorname{Res}(A,B)
 =(-1)^{mn}b_n^m\prod_jR(\beta_j)
\]
and
\[
 \operatorname{Res}(R,B)
 =(-1)^{rn}b_n^r\prod_jR(\beta_j).
\]
Eliminating the product and using
$\operatorname{Res}(R,B)=(-1)^{rn}\operatorname{Res}(B,R)$ yields
\eqref{eq:resultant-remainder}.  If $R=0$, then $B$ divides $A$, so the
resultant vanishes by Theorem~\ref{thm:resultant-zero-criterion}.
\end{proof}

For fraction-free arithmetic let
\[
 P=\operatorname{prem}(A,B).
\]
With $s=m-n+1$, pseudo-division satisfies
\[
 b_n^sA=QB+P.
\]
The ordinary remainder is therefore $R=P/b_n^s$ over the quotient field.

\begin{corollary}[Pseudo-remainder recurrence]
\label{cor:resultant-pseudoremainder}
Under the hypotheses of Lemma~\ref{lem:resultant-remainder}, let
$P=\operatorname{prem}(A,B)\ne0$ and put $r=\deg P$.  Then
\begin{equation}\label{eq:resultant-pseudoremainder}
 \operatorname{Res}(A,B)
 =(-1)^{mn}
  b_n^{\,m-r-n(m-n+1)}
  \operatorname{Res}(B,P).
\end{equation}
The exponent may be negative; the identity is an equality in the quotient
field, and the complete expression is nevertheless integral.
\end{corollary}

\begin{proof}
Since $P=b_n^sR$, scaling the second input gives
\[
 \operatorname{Res}(B,P)=b_n^{sn}\operatorname{Res}(B,R).
\]
Substitute this into \eqref{eq:resultant-remainder}.
\end{proof}

The negative powers in \eqref{eq:resultant-pseudoremainder} explain why a
naive resultant recurrence based directly on pseudo-remainders is awkward:
one repeatedly creates large powers only to divide them out again.  Brown's
subresultant normalization packages exactly these cancellations into known
exact factors.  The zeroth subresultant then emerges as an integer without
rational intermediate arithmetic.

\subsection{The terminal Brown invariant}

Recall the notation of Section~\ref{sec:subresultant-prs}.  Brown's PRS
produces members $F_i$ of degrees $n_i$ and auxiliary integers
\[
 h_i=\operatorname{lc}\operatorname{Sres}_{n_i}(F_1,F_2).
\]
The following observation turns the same PRS into a resultant algorithm.

\begin{theorem}[Terminal subresultant scalar]\label{thm:terminal-resultant}
Let $F_1,F_2\in\ZZ[x]$ be primitive and nonzero, with
$\deg F_1\ge\deg F_2$, and run Brown's subresultant PRS to its last nonzero
member $F_k$.
\begin{enumerate}
 \item If $\deg F_k>0$, then
 \[
  \operatorname{Res}(F_1,F_2)=0.
 \]
 \item If $\deg F_k=0$, then
 \begin{equation}\label{eq:terminal-h-resultant}
  h_k=\operatorname{Res}(F_1,F_2).
 \end{equation}
\end{enumerate}
The constant polynomial $F_k$ itself need not equal the resultant.
\end{theorem}

\begin{proof}
By Theorem~\ref{thm:brown-subresultant-prs-correct},
\[
 h_k=\operatorname{lc}\operatorname{Sres}_{n_k}(F_1,F_2).
\]
If $n_k>0$, the final clause of the fundamental theorem of subresultants says
that every subresultant of index below $n_k$ vanishes.  In particular
\[
 \operatorname{Sres}_0(F_1,F_2)=0.
\]
The zeroth subresultant is the Sylvester determinant, hence the resultant.

If $n_k=0$, the displayed identity becomes
\[
 h_k=\operatorname{lc}\operatorname{Sres}_0(F_1,F_2).
\]
Since the zeroth subresultant is a constant, its leading coefficient is the
constant itself, proving \eqref{eq:terminal-h-resultant}.
\end{proof}

\begin{example}[The last constant need not be the resultant]
Let
\[
 A=-2-x-5x^2,
 \qquad
 B=-3-x-5x^2.
\]
Both are primitive.  Brown's final polynomial subresultant is the constant
$5$, while its associated lower-subresultant scalar is
\[
 h=25.
\]
Since the final degree is zero, Theorem~\ref{thm:terminal-resultant} gives
\[
 \operatorname{Res}(A,B)=25.
\]
Thus returning the last PRS polynomial instead of the auxiliary $h$ would be
incorrect even on small primitive inputs.
\end{example}

The theorem also unifies the gcd, XGCD and resultant interpretations of the
same subresultant sequence.  A positive-degree terminal subresultant is a
rational associate of the gcd; a zero-degree terminal lower subresultant is
the resultant; and the bordered minors propagated in the extended sequence
supply the integral B\'ezout cofactors.

\subsection{Subresultant resultant algorithm over $\ZZ[x]$}

For arbitrary integer polynomials, contents and a possible input swap remain
to be restored.  Let
\[
 c_A=\operatorname{cont}(A),
 \qquad
 c_B=\operatorname{cont}(B),
 \qquad
 A_0=\operatorname{pp}(A),
 \qquad
 B_0=\operatorname{pp}(B).
\]
If $m=\deg A$ and $n=\deg B$, then
\begin{equation}\label{eq:resultant-content-law}
 \operatorname{Res}(A,B)
 =c_A^n c_B^m\operatorname{Res}(A_0,B_0).
\end{equation}

\begin{algorithm}[Subresultant resultant]
\label{alg:subresultant-resultant}
Input $A,B\in\ZZ[x]$.
\begin{enumerate}
 \item If either input is zero, return $0$.
 \item Put $m\leftarrow\deg A$, $n\leftarrow\deg B$.  If $m=n=0$, return
 $1$.  If $m=0$, return $A^n$; if $n=0$, return $B^m$.
 \item Compute $c_A,c_B,A_0,B_0$ as above.
 \item If $\deg A_0<\deg B_0$, interchange $A_0,B_0$ and remember that the
 pair was swapped.
 \item Run Brown's subresultant PRS on the ordered primitive pair, retaining
 its final nonzero member $S$ and the corresponding auxiliary scalar $h$.
 If $\deg S>0$, put $r\leftarrow0$; otherwise put $r\leftarrow h$.
 \item If the inputs were swapped and $mn$ is odd, replace $r$ by $-r$.
 \item Return
 \begin{equation}\label{eq:resultant-alg-return}
  c_A^n c_B^m r.
 \end{equation}
\end{enumerate}
\end{algorithm}

\begin{theorem}[Correctness of the subresultant resultant]
\label{thm:subresultant-resultant-correct}
Algorithm~\ref{alg:subresultant-resultant} returns
$\operatorname{Res}(A,B)$ with the Sylvester sign convention
\eqref{eq:resultant-sylvester}.
\end{theorem}

\begin{proof}
The constant cases are \eqref{eq:resultant-constant}, with the empty
Sylvester determinant equal to $1$ when both inputs are constants.  Assume
therefore that both degrees are positive.  Scaling law
\eqref{eq:resultant-scaling} gives
\eqref{eq:resultant-content-law}.

If the primitive inputs are already in nonincreasing degree order,
Theorem~\ref{thm:terminal-resultant} says that step 5 returns exactly
$\operatorname{Res}(A_0,B_0)$: a positive-degree last member is equivalent to
a common factor and hence zero resultant, while a constant last member makes
$h$ the zeroth subresultant.

If the primitive inputs were interchanged, equation
\eqref{eq:resultant-symmetry} shows that one must multiply the ordered
resultant by $(-1)^{mn}$.  Step 6 does exactly this.  Finally step 7 restores
the two contents by \eqref{eq:resultant-content-law}.
\end{proof}

The companion implementation uses this algorithm directly.  Its shared Brown
PRS routine returns the pair $(S,h)$ from step 5; the resultant routine tests
the degree of $S$ but returns $h$, not the constant coefficient of $S$.

\subsection{Coefficient growth and complexity}

The determinant definition gives the output-size bound before any algorithm
is chosen.

\begin{theorem}[Resultant coefficient bound]\label{thm:resultant-size}
Let $A,B\in\ZZ[x]$ have degrees at most $N$ and coefficients of at most
$\tau$ bits.  Then
\begin{equation}\label{eq:resultant-size}
 \bits\bigl(\operatorname{Res}(A,B)\bigr)
 =O\bigl(N(\tau+\log(N+1))\bigr).
\end{equation}
\end{theorem}

\begin{proof}
The Sylvester determinant has order $m+n\le2N$.  Each row contains at most
$N+1$ nonzero input coefficients, so its Euclidean norm is at most
\[
 \sqrt{N+1}\,2^\tau.
\]
Hadamard's inequality therefore gives
\[
 |\operatorname{Res}(A,B)|
 \le
 \bigl(\sqrt{N+1}\,2^\tau\bigr)^{m+n}.
\]
Taking binary logarithms proves \eqref{eq:resultant-size}.
\end{proof}

A direct determinant computation and the subresultant method have different
polynomial arithmetic costs.  Fraction-free Gaussian elimination, or Bareiss
elimination, on the $(m+n)\times(m+n)$ Sylvester matrix uses
$O(N^3)$ integer arithmetic operations.  Brown's PRS exploits the Toeplitz-like
polynomial structure of the same determinant and needs only quadratic
coefficient arithmetic in the classical implementation.

\begin{theorem}[Coefficient complexity of the subresultant resultant]
\label{thm:resultant-coefficient-complexity}
For primitive inputs of degrees at most $N$, Algorithm~\ref{alg:subresultant-resultant}
with classical pseudo-remainders uses $O(N^2)$ coefficient additions and
multiplications and $O(N^2)$ exact coefficient divisions by known Brown
factors.  For arbitrary inputs it additionally computes two contents and
performs $O(N)$ exact scalar divisions to form the primitive parts.
\end{theorem}

\begin{proof}
The PRS computation is the same one analyzed in
Theorem~\ref{thm:subresultant-gcd-coefficient-complexity}; stopping with the
terminal pair $(S,h)$ does not add any polynomial operation.  Hence its
coefficient arithmetic and exact Brown divisions are both $O(N^2)$.  Forming
the two primitive inputs touches only $O(N)$ coefficients.  The remaining
swap, sign and content-power operations are scalar integer operations.
\end{proof}

As in the gcd computation, stored Brown-normalized remainders satisfy the
subresultant determinant bound
\[
 W=O\bigl(N(\tau+\log(N+1))\bigr),
\]
but the straightforward pseudo-remainder kernel forms a larger raw remainder
before dividing by the Brown factor.  Put
\[
 B_*=O\bigl(N(W+1)\bigr).
\]
The direct bit bound therefore parallels
Theorem~\ref{thm:subresultant-gcd-bit-complexity}.

\begin{theorem}[Direct bit complexity of the subresultant resultant]
\label{thm:resultant-bit-complexity}
With the notation above, the classical Brown-PRS resultant has bit complexity
\begin{equation}\label{eq:resultant-bit-complexity}
 O\!\left(
 N^2\bigl(\MZ(B_*)+\mathsf D_{\ZZ}(B_*)\bigr)
 +N\bigl(\mathsf G_{\ZZ}(\tau)+\mathsf D_{\ZZ}(\tau)\bigr)
 +\log(N+1)\MZ(W)
 \right).
\end{equation}
The final term covers binary powering of the two contents and multiplication
of those powers into the primitive resultant.  In the balanced nontrivial
regime it is dominated by the displayed PRS bound.
\end{theorem}

\begin{proof}
The PRS contribution is identical to the proof of
Theorem~\ref{thm:subresultant-gcd-bit-complexity}: there are $O(N^2)$
coefficient arithmetic operations and exact normalizing divisions on
integers of at most $O(B_*)$ bits.  The two content computations and primitive
part divisions give the second term.  Each content has at most $\tau$ bits,
and its exponent is at most $N$; the powers therefore have
$O(N(\tau+1))\subseteq O(W)$ bits.  Binary exponentiation uses
$O(\log(N+1))$ integer multiplications, and the final products with the
primitive resultant also have $O(W)$-bit operands.  This gives the last term.
\end{proof}

The resultant section therefore closes the subresultant sequence developed in
the preceding sections.  The Sylvester determinant supplies the definition
and output-size bound; Brown normalization evaluates its zeroth subresultant
without constructing the matrix; the last positive-degree subresultant gives
the gcd information; and the extended minors supply the B\'ezout certificate.
The discriminant and squarefreeness criteria are immediate applications of
this resultant machinery and will be treated next.


\section{Discriminants and squarefreeness}\label{sec:discriminant-squarefree}

The derivative and resultant together detect repeated roots.  This section
extracts two standard consequences: the discriminant, which records the
pairwise root differences in one integer, and the squarefree part, which
removes repeated irreducible factors.  Throughout, squarefreeness of an
integer polynomial means squarefreeness after extension of scalars to
$\mathbb Q[x]$; integer content is therefore ignored.  This is the convention
used by the companion implementation.

Let
\[
 f(x)=a_nx^n+\cdots+a_0\in\ZZ[x],\qquad a_n\ne0.
\]
For $n\ge1$ write $f'$ for the formal derivative.  Differentiation requires
only the scalar products
\[
 [x^{i-1}]f'=i a_i,
\]
so it costs $O(n)$ coefficient operations and increases coefficient bit size
by at most $O(\log(n+1))$.

\subsection{The discriminant}

\begin{definition}[Discriminant]\label{def:discriminant}
For $n\ge1$, define
\begin{equation}\label{eq:discriminant-resultant}
 \operatorname{Disc}(f)
 =(-1)^{n(n-1)/2}\frac{\operatorname{Res}(f,f')}{a_n}.
\end{equation}
The companion implementation assigns discriminant $0$ to polynomials of
length at most one, including the zero polynomial and nonzero constants.
For a linear polynomial, \eqref{eq:discriminant-resultant} gives
$\operatorname{Disc}(f)=1$.
\end{definition}

The division by $a_n$ in \eqref{eq:discriminant-resultant} is exact over
$\ZZ$.  It is useful to prove this before appealing to roots, because it
justifies the fraction-free implementation directly.

\begin{theorem}[Integrality of the discriminant]\label{thm:discriminant-integral}
For every $f\in\ZZ[x]$ of positive degree,
\[
 a_n\mid\operatorname{Res}(f,f'),
\]
so \eqref{eq:discriminant-resultant} belongs to $\ZZ$.
Indeed, as a universal expression in the coefficients,
\[
 \frac{\operatorname{Res}(f,f')}{a_n}
 \in\ZZ[a_0,\ldots,a_n].
\]
\end{theorem}

\begin{proof}
Regard
\[
 R(a_0,\ldots,a_n)=\operatorname{Res}(f,f')
\]
as the determinant of the $(2n-1)\times(2n-1)$ Sylvester matrix over the
polynomial ring $\ZZ[a_0,\ldots,a_n]$.  Set $a_n=0$.  Then the formal degree
of $f$ drops below $n$ and the formal degree of $f'$ drops below $n-1$.
Consequently the highest-degree column of the Sylvester matrix is zero, so
$R(a_0,\ldots,a_{n-1},0)=0$.

Viewed as a polynomial in the single indeterminate $a_n$ over
$\ZZ[a_0,\ldots,a_{n-1}]$, the constant term of $R$ is therefore zero.
Hence $a_n$ divides $R$ in the coefficient polynomial ring, proving the
claim after specialization to arbitrary integer coefficients.
\end{proof}

The root formula explains the geometry behind the definition.

\begin{theorem}[Root-difference formula]\label{thm:discriminant-root-formula}
Let $L$ be a splitting field and write
\[
 f(x)=a_n\prod_{i=1}^n(x-\alpha_i),
\]
with roots repeated according to multiplicity.  Then
\begin{equation}\label{eq:discriminant-root-formula}
 \operatorname{Disc}(f)
 =a_n^{2n-2}\prod_{1\le i<j\le n}(\alpha_i-\alpha_j)^2.
\end{equation}
\end{theorem}

\begin{proof}
The product formula for the resultant gives
\[
 \operatorname{Res}(f,f')
 =a_n^{n-1}\prod_{i=1}^n f'(\alpha_i).
\]
For every $i$,
\[
 f'(\alpha_i)=a_n\prod_{j\ne i}(\alpha_i-\alpha_j).
\]
Thus
\[
 \operatorname{Res}(f,f')
 =a_n^{2n-1}
   \prod_i\prod_{j\ne i}(\alpha_i-\alpha_j).
\]
Each unordered pair $\{i,j\}$ contributes
\[
 (\alpha_i-\alpha_j)(\alpha_j-\alpha_i)
 =-(\alpha_i-\alpha_j)^2.
\]
There are $n(n-1)/2$ such pairs.  Substitution into
\eqref{eq:discriminant-resultant} proves the formula.
\end{proof}

\begin{example}[Quadratic discriminant]
For
\[
 f=ax^2+bx+c,
\]
formula \eqref{eq:discriminant-resultant} gives
\[
 \operatorname{Disc}(f)=b^2-4ac.
\]
For example, $3x^2+2x+1$ has discriminant $-8$.
\end{example}

\begin{corollary}[Repeated-root criterion]\label{cor:discriminant-repeated-root}
For $f\in\ZZ[x]$ of positive degree,
\[
 \operatorname{Disc}(f)=0
 \quad\Longleftrightarrow\quad
 f\text{ has a repeated root in }\overline{\mathbb Q}
 \quad\Longleftrightarrow\quad
 \deg\gcd_{\mathbb Q[x]}(f,f')>0.
\]
\end{corollary}

\begin{proof}
By \eqref{eq:discriminant-root-formula}, the discriminant vanishes exactly
when two roots coincide.  In characteristic zero this is equivalent to a
nonconstant common factor of $f$ and $f'$.
\end{proof}

Two useful identities also follow immediately from the root formula.

\begin{proposition}[Scaling and products]\label{prop:discriminant-laws}
Let $f,g\in\ZZ[x]$ have positive degrees $m,n$, and let $c\ne0$ be an
integer.  Then
\begin{align}
 \operatorname{Disc}(cf)
 &=c^{2m-2}\operatorname{Disc}(f),
 \label{eq:discriminant-scaling}\\
 \operatorname{Disc}(fg)
 &=\operatorname{Disc}(f)\operatorname{Disc}(g)
   \operatorname{Res}(f,g)^2.
 \label{eq:discriminant-product}
\end{align}
\end{proposition}

\begin{proof}
Scaling changes only the leading coefficient in
\eqref{eq:discriminant-root-formula}, giving
\eqref{eq:discriminant-scaling}.  For a product, split the pairwise root
differences into pairs internal to $f$, pairs internal to $g$, and cross
pairs.  The first two groups give the two discriminants.  The cross group,
including the required leading-coefficient powers, is exactly
$\operatorname{Res}(f,g)^2$ by
Theorem~\ref{thm:resultant-product-formula}.
\end{proof}

\begin{algorithm}[Discriminant over integer polynomials]\label{alg:zz-discriminant}
Input $f\in\ZZ[x]$.
\begin{enumerate}
 \item If $\len(f)\le1$, return $0$.
 \item Put $n\leftarrow\deg f$ and compute $d\leftarrow f'$.
 \item Compute $r\leftarrow\operatorname{Res}(f,d)$.
 \item Divide $r$ exactly by $\operatorname{lc}(f)$.
 \item If $n(n-1)/2$ is odd, replace $r$ by $-r$.
 \item Return $r$.
\end{enumerate}
\end{algorithm}

\begin{theorem}[Correctness of Algorithm~\ref{alg:zz-discriminant}]
Algorithm~\ref{alg:zz-discriminant} returns the convention of
Definition~\ref{def:discriminant}.
\end{theorem}

\begin{proof}
For positive degree this is exactly
\eqref{eq:discriminant-resultant}; the exactness of step 4 is
Theorem~\ref{thm:discriminant-integral}.  The low-degree branch is the stated
companion convention.
\end{proof}

The implementation tests the parity in step 5 through $n\bmod4$: the sign is
negative exactly for $n\equiv2,3\pmod4$.

\subsection{Squarefree polynomials over $\mathbb Q$}

For integer polynomials there are two possible notions of squarefreeness.  As
an element of $\ZZ[x]$, a polynomial with content divisible by $p^2$ is not
squarefree.  Polynomial algorithms usually want instead the factor structure
in $\mathbb Q[x]$, where every nonzero integer is a unit.  We use the latter
notion.

\begin{definition}[Rational squarefreeness]\label{def:rational-squarefree}
A nonzero polynomial $f\in\ZZ[x]$ is \emph{squarefree} if no nonconstant
square in $\mathbb Q[x]$ divides $f$.  Equivalently, in a factorization
\[
 f=c\prod_{i=1}^s p_i^{e_i},
 \qquad c\in\mathbb Q^\times,
\]
into pairwise nonassociate irreducibles $p_i\in\mathbb Q[x]$, every exponent
is $e_i=1$.  Nonzero constants are squarefree.  The zero polynomial is not
squarefree.
\end{definition}

Thus multiplying by a nonzero integer does not change squarefreeness.  This is
why, for example, $12(x+1)(x+2)$ is squarefree under the companion API.

\begin{theorem}[Derivative criterion]\label{thm:squarefree-derivative-criterion}
Let $f\in\ZZ[x]$ be nonzero.  Then
\begin{equation}\label{eq:squarefree-gcd-criterion}
 f\text{ is squarefree}
 \quad\Longleftrightarrow\quad
 \deg\gcd_{\mathbb Q[x]}(f,f')=0.
\end{equation}
For positive degree this is also equivalent to
$\operatorname{Disc}(f)\ne0$.
\end{theorem}

\begin{proof}
Over $\mathbb Q[x]$, write
\[
 f=c\prod_i p_i^{e_i}.
\]
Differentiation gives
\[
 f'=c\sum_i e_i p_i^{e_i-1}p_i'
       \prod_{j\ne i}p_j^{e_j}.
\]
Because the characteristic is zero, every irreducible $p_i$ is separable and
$p_i\nmid p_i'$.  Hence the exponent of $p_i$ in $\gcd(f,f')$ is exactly
$e_i-1$.  The gcd is constant precisely when every $e_i=1$.  The
discriminant equivalence is Corollary~\ref{cor:discriminant-repeated-root}.
\end{proof}

The canonical gcd over $\ZZ[x]$ may contain a nonunit integer content, but
that does not affect the degree test.  If $f=cF$ with $c=\operatorname{cont}f$
and $F=\operatorname{pp}(f)$, then
\[
 \gcd_{\ZZ[x]}(f,f')
 =c\,H
\]
for a primitive polynomial $H$ associated over $\mathbb Q[x]$ to
$\gcd(F,F')$.  Thus checking whether the integer-polynomial gcd has length at
most one implements \eqref{eq:squarefree-gcd-criterion} exactly.

\begin{algorithm}[Squarefree test]\label{alg:is-squarefree}
Input $f\in\ZZ[x]$.
\begin{enumerate}
 \item If $f=0$, return false.
 \item If $\deg f\le1$, return true.
 \item Compute $d\leftarrow f'$ and $g\leftarrow\gcd_{\ZZ[x]}(f,d)$.
 \item Return true if and only if $\deg g=0$.
\end{enumerate}
\end{algorithm}

\begin{theorem}[Correctness of Algorithm~\ref{alg:is-squarefree}]
Algorithm~\ref{alg:is-squarefree} returns true exactly for the squarefree
polynomials of Definition~\ref{def:rational-squarefree}.
\end{theorem}

\begin{proof}
The zero, constant and linear cases agree with the definition.  For degree at
least two, the integer gcd and rational gcd have the same polynomial degree,
so the result follows from Theorem~\ref{thm:squarefree-derivative-criterion}.
\end{proof}

Although the discriminant gives another squarefree test, the companion
implementation uses the gcd criterion.  This avoids constructing the possibly
large scalar resultant when only the existence of a repeated factor is
required.

\subsection{The squarefree part}

The quotient by the derivative gcd removes precisely the excess copies of
each irreducible factor.

\begin{definition}[Squarefree part]\label{def:squarefree-part}
For nonzero $f\in\ZZ[x]$, let
\[
 G=\gcd_{\ZZ[x]}(f,f')
\]
with the canonical positive-leading-coefficient convention.  Define
\begin{equation}\label{eq:squarefree-part}
 \operatorname{sqfpart}(f)
 =\operatorname{pos}\!\left(\frac{f}{G}\right),
\end{equation}
where $\operatorname{pos}(q)$ is $q$ or $-q$ according as its leading
coefficient is positive or negative.  Set
$\operatorname{sqfpart}(0)=0$.
\end{definition}

\begin{theorem}[Structure of the squarefree part]
\label{thm:squarefree-part-structure}
Let $f\ne0$ and factor its primitive part over $\mathbb Q[x]$ as
\[
 \operatorname{pp}(f)=u\prod_{i=1}^s p_i^{e_i},
\]
where the $p_i$ are pairwise nonassociate primitive irreducible integer
polynomials and the unit $u\in\{\pm1\}$ accounts for sign.  Then
\begin{equation}\label{eq:squarefree-part-factorization}
 \operatorname{sqfpart}(f)
 =\operatorname{pos}\!\left(u\prod_{i=1}^s p_i\right).
\end{equation}
In particular, $\operatorname{sqfpart}(f)$ is primitive, has positive leading
coefficient, is squarefree, and has exactly the distinct roots of $f$.
For a nonzero constant, $\operatorname{sqfpart}(f)=1$.
\end{theorem}

\begin{proof}
The proof of Theorem~\ref{thm:squarefree-derivative-criterion} shows that the
primitive polynomial part of $\gcd(f,f')$ contains $p_i^{e_i-1}$ for every
$i$.  Moreover, the common integer content of $f$ and $f'$ contains the full
positive content of $f$, so dividing by the canonical integer-polynomial gcd
removes the content as well.  Therefore $f/G$ is, up to sign, exactly the
primitive product $\prod_i p_i$.  Gauss's lemma makes this quotient primitive.
The final sign normalization gives
\eqref{eq:squarefree-part-factorization}.  The constant case is the empty
product.
\end{proof}

\begin{corollary}[Radical ideal interpretation]
For nonzero $f\in\ZZ[x]$,
\[
 \sqrt{(f)\mathbb Q[x]}
 =(\operatorname{sqfpart}(f))\mathbb Q[x].
\]
\end{corollary}

\begin{proof}
Both ideals are generated, up to a rational unit, by the product of the
distinct irreducible factors of $f$.
\end{proof}

\begin{algorithm}[Squarefree part over integer polynomials]\label{alg:squarefree-part}
Input $f\in\ZZ[x]$.
\begin{enumerate}
 \item If $f=0$, return $0$.
 \item Compute $d\leftarrow f'$ and $g\leftarrow\gcd_{\ZZ[x]}(f,d)$.
 \item Compute the exact quotient $q\leftarrow f/g$.
 \item If $\operatorname{lc}(q)<0$, replace $q$ by $-q$.
 \item Return $q$.
\end{enumerate}
\end{algorithm}

\begin{theorem}[Correctness of the squarefree-part algorithm]
Algorithm~\ref{alg:squarefree-part} returns
$\operatorname{sqfpart}(f)$ from Definition~\ref{def:squarefree-part}.
\end{theorem}

\begin{proof}
For $f=0$ this is the chosen convention.  Otherwise the quotient in step 3 is
exact by definition of the gcd, and Theorem~\ref{thm:squarefree-part-structure}
identifies it up to sign with the product of the distinct primitive factors.
Step 4 imposes the canonical positive-leading-coefficient normalization.
\end{proof}

\subsection{Coefficient growth and complexity}

Discriminants inherit both their size and their computational cost from the
resultant.  If $f$ has degree at most $N$ and input coefficients of at most
$\tau$ bits, then the coefficients of $f'$ have
$O(\tau+\log(N+1))$ bits.  The resultant bound from
Theorem~\ref{thm:resultant-size} immediately yields the following.

\begin{theorem}[Discriminant size]\label{thm:discriminant-size}
For $f\in\ZZ[x]$ of degree at most $N$ and coefficient bit size at most
$\tau$,
\[
 \bits(\operatorname{Disc}(f))
 =O\bigl(N(\tau+\log(N+1))\bigr).
\]
\end{theorem}

\begin{proof}
Apply Theorem~\ref{thm:resultant-size} to $(f,f')$.  Exact division by the
nonzero leading coefficient cannot increase the absolute value when
$|a_n|\ge1$.
\end{proof}

\begin{theorem}[Coefficient complexity]\label{thm:disc-squarefree-coefficient-complexity}
With classical Brown-PRS gcd and resultant algorithms, for input degree at
most $N$:
\begin{enumerate}
 \item the discriminant uses $O(N^2)$ polynomial coefficient operations;
 \item the squarefree test uses $O(N^2)$ polynomial coefficient operations;
 \item the squarefree part uses $O(N^2)$ polynomial coefficient operations
 in addition to one exact polynomial division.
\end{enumerate}
With a classical exact division, all three bounds are $O(N^2)$ coefficient
operations.
\end{theorem}

\begin{proof}
Differentiation is linear.  The first claim then follows from
Theorem~\ref{thm:resultant-coefficient-complexity}; the final division by the
leading coefficient is one scalar exact division.  The second follows from
Theorem~\ref{thm:subresultant-gcd-coefficient-complexity}.  The third performs
the same gcd and then an exact quotient of degree at most $N$.
\end{proof}

For bit complexity, the direct bounds from the preceding sections may be
reused with the derivative coefficient size
$\tau'=O(\tau+\log(N+1))$.  Thus the discriminant has the bit cost of the
Brown-PRS resultant on $(f,f')$, followed by one exact division of an
$O(N(\tau+\log N))$-bit integer by the leading coefficient.  The squarefree
test has the corresponding Brown-PRS gcd cost.  The squarefree part adds the
bit cost of one exact polynomial quotient.  As before, these direct bounds
include the raw pseudo-remainders formed before Brown normalization; the
stored normalized PRS coefficients themselves satisfy determinant-size
bounds.

The resulting hierarchy is therefore simple.  The discriminant is the scalar
certificate of repeated roots, the derivative gcd identifies their common
polynomial factor, and the squarefree part removes all excess multiplicities.
All three operations reuse the resultant, gcd and exact-division machinery
already developed in this chapter.


\begin{thebibliography}{99}

\bibitem{KaratsubaOfman1962}
A.~Karatsuba and Yu.~Ofman,
\emph{Multiplication of many-digital numbers by automatic computers},
Dokl. Akad. Nauk SSSR \textbf{145} (1962), 293--294.
\newline
\url{https://www.mathnet.ru/eng/dan26729}
(full-text PDF available from the page).

\bibitem{Toom1963}
A.~L.~Toom,
\emph{The complexity of a scheme of functional elements simulating the
multiplication of integers},
Dokl. Akad. Nauk SSSR \textbf{150} (1963), 496--498.
\newline
\url{https://www.mathnet.ru/eng/dan27978}
(full-text PDF available from the page).

\bibitem{Cook1966}
S.~A.~Cook,
\emph{On the Minimum Computation Time of Functions},
PhD thesis, Harvard University, 1966, Chapter III.
\newline
\url{https://www.cs.toronto.edu/~sacook/}
(scan of Chapter III available from the author's page).

\bibitem{BrentZimmermann2010}
R.~P.~Brent and P.~Zimmermann,
\emph{Modern Computer Arithmetic},
Cambridge University Press, 2010.
Preliminary version 0.5.9 freely available from
\newline
\url{https://maths-people.anu.edu.au/~brent/pub/pub226.html};
see also \url{https://arxiv.org/abs/1004.4710}.

\bibitem{Estrin1960}
G.~Estrin,
\emph{Organization of Computer Systems---The Fixed Plus Variable Structure
Computer},
in \emph{Proceedings of the Western Joint IRE--AIEE--ACM Computer Conference},
1960, 33--40.
\newline
\url{https://doi.org/10.1145/1460361.1460365}.

\bibitem{BrentKung1978}
R.~P.~Brent and H.~T.~Kung,
\emph{Fast Algorithms for Manipulating Formal Power Series},
Journal of the ACM \textbf{25} (1978), 581--595.
\newline
\url{https://doi.org/10.1145/322092.322099};
author PDF and bibliographic page:
\url{https://maths-people.anu.edu.au/~brent/pub/pub045.html}.

\bibitem{Mulders2000}
T.~Mulders,
\emph{On Short Multiplications and Divisions},
Applicable Algebra in Engineering, Communication and Computing \textbf{11}
(2000), 69--88.
\newline
\url{https://doi.org/10.1007/s002000000037};
preprint PDF:
\url{https://citeseerx.ist.psu.edu/document?doi=11d05ac21c5469d929ea983eda85aa90d0824669&repid=rep1&type=pdf}.

\bibitem{HartNovocin2011}
W.~B.~Hart and A.~Novocin,
\emph{Practical Divide-and-Conquer Algorithms for Polynomial Arithmetic},
in \emph{Computer Algebra in Scientific Computing, CASC 2011},
Lecture Notes in Computer Science \textbf{6885}, Springer, 2011, 200--214.
\newline
\url{https://doi.org/10.1007/978-3-642-23568-9_16};
accepted-version PDF:
\url{https://wrap.warwick.ac.uk/id/eprint/43601/1/WRAP_Hart_0584144-ma-270913-polycomp.pdf}.

\bibitem{KrandickJebelean1996}
W.~Krandick and T.~Jebelean,
\emph{Bidirectional Exact Integer Division},
J. Symbolic Comput. \textbf{21} (1996), 441--455.
\newline
\url{https://doi.org/10.1006/jsco.1996.0024};
early technical report RISC 94--50:
\url{ftp://ftp.risc.uni-linz.ac.at/pub/techreports/1994/94-50.ps.gz}.

\bibitem{HanrotZimmermann2004}
G.~Hanrot and P.~Zimmermann,
\emph{A long note on Mulders' short product},
J. Symbolic Comput. \textbf{37} (2004), 391--401.
\newline
\url{https://doi.org/10.1016/j.jsc.2003.03.001};
freely available PDF:
\url{https://www.csd.uwo.ca/~mmorenom/CS424/Ressources/Hanrot-Zimmermann-JSC-2004.pdf}.

\bibitem{Bodrato2007}
M.~Bodrato,
\emph{Towards Optimal Toom--Cook Multiplication for Univariate and
Multivariate Polynomials in Characteristic 2 and 0},
in \emph{Arithmetic of Finite Fields -- WAIFI 2007},
Lecture Notes in Computer Science \textbf{4547}, Springer, 2007, 116--133.
\newline
\url{https://doi.org/10.1007/978-3-540-73074-3_10};
author's paper index:
\url{http://bodrato.it/papers/#WAIFI2007}.


\bibitem{NissimSchwartzSpiizer2026}
R.~Nissim, O.~Schwartz and Y.~Spiizer,
\emph{The Division Barrier: Optimal Bounds and Structural Limits in
Toom--Cook Interpolation},
in \emph{Proceedings of the 2026 Annual ACM--SIAM Symposium on Discrete
Algorithms (SODA)}, SIAM, 2026, 5241--5254.
\newline
\url{https://doi.org/10.1137/1.9781611978971.190}.

\bibitem{Collins1967}
G.~E.~Collins,
\emph{Subresultants and Reduced Polynomial Remainder Sequences},
Journal of the ACM \textbf{14} (1967), 128--142.
\newline
\url{https://doi.org/10.1145/321371.321381};
publication page with abstract and resource links:
\url{https://research.ibm.com/publications/subresultants-and-reduced-polynomial-remainder-sequences}.

\bibitem{Brown1978}
W.~S.~Brown,
\emph{The Subresultant PRS Algorithm},
ACM Transactions on Mathematical Software \textbf{4} (1978), 237--249.
\newline
\url{https://doi.org/10.1145/355791.355795};
freely available PDF:
\url{https://people.eecs.berkeley.edu/~fateman/282/readings/brown.pdf}.

\bibitem{BrownTraub1971}
W.~S.~Brown and J.~F.~Traub,
\emph{On Euclid's Algorithm and the Theory of Subresultants},
Journal of the ACM \textbf{18} (1971), 505--514.
\newline
\url{https://doi.org/10.1145/321662.321665};
full scan from the Carnegie Mellon University archives:
\url{https://iiif.library.cmu.edu/file/Traub_box00027_fld00059_bdl0001_doc0001/Traub_box00027_fld00059_bdl0001_doc0001.pdf}.

\bibitem{KarpMarkstein1997}
A.~H.~Karp and P.~Markstein,
\emph{High-Precision Division and Square Root},
ACM Transactions on Mathematical Software \textbf{23} (1997), 561--589.
\newline
\url{https://doi.org/10.1145/279232.279237};
earlier Hewlett--Packard technical report HPL-93-42(R.1), with a freely
available copy indexed at
\url{https://cr.yp.to/bib/entries.html#1997/karp}.

\bibitem{HanrotQuerciaZimmermann2004}
G.~Hanrot, M.~Quercia and P.~Zimmermann,
\emph{The Middle Product Algorithm I: Speeding up the Division and Square Root
of Power Series},
Applicable Algebra in Engineering, Communication and Computing \textbf{14}
(2004), 415--438.
\newline
\url{https://doi.org/10.1007/s00200-003-0144-2};
freely available PDF:
\url{https://cr.yp.to/bib/2004/hanrot-middle.pdf}.

\bibitem{BostanLecerfSchost2003}
A.~Bostan, G.~Lecerf and E.~Schost,
\emph{Tellegen's Principle into Practice},
in \emph{Proceedings of ISSAC 2003}, ACM, 2003, 37--44.
\newline
\url{https://doi.org/10.1145/860854.860870};
author PDF:
\url{https://cs.uwaterloo.ca/~eschost/publications/p33-bostan.pdf}.

\bibitem{Harvey2012Middle}
D.~Harvey,
\emph{The Karatsuba Integer Middle Product},
Journal of Symbolic Computation \textbf{47} (2012), 954--967.
\newline
\url{https://doi.org/10.1016/j.jsc.2012.02.001};
author PDF:
\url{https://web.maths.unsw.edu.au/~davidharvey/research/mulmid/mulmid.pdf}.

\bibitem{Harvey2009}
D.~Harvey,
\emph{Faster polynomial multiplication via multipoint Kronecker substitution},
J. Symbolic Comput. \textbf{44} (2009), 1502--1510.
\newline
\url{https://arxiv.org/abs/0712.4046}.



\bibitem{ElKahoui2003}
M.~El Kahoui,
\emph{An elementary approach to subresultants theory},
Journal of Symbolic Computation \textbf{35} (2003), 281--292.
\newline
\url{https://doi.org/10.1016/S0747-7171(02)00135-9};
author-uploaded full text is available from
\url{https://www.researchgate.net/publication/220160919_An_elementary_approach_to_subresultants_theory}.


\bibitem{Sturmfels2002}
B.~Sturmfels,
\emph{Solving Systems of Polynomial Equations},
CBMS Regional Conference Series in Mathematics, vol.~97,
American Mathematical Society, 2002.
\newline
Freely available author PDF:
\url{https://math.berkeley.edu/~bernd/cbms.pdf}.

\bibitem{Moenck1973}
R.~T.~Moenck,
\emph{Fast Computation of GCDs},
in \emph{Proceedings of the Fifth Annual ACM Symposium on Theory of Computing
(STOC 1973)}, ACM, 1973, 142--151.
\newline
Bibliographic entry with scanned version:
\url{https://cr.yp.to/bib/entries.html#1973/moenck}.

\bibitem{ThullYap1990}
K.~Thull and C.~K.~Yap,
\emph{A Unified Approach to HGCD Algorithms for Polynomials and Integers},
manuscript, 1990.
\newline
Full-text copy:
\url{https://www.researchgate.net/publication/2307268_A_Unified_Approach_to_HGCD_Algorithms_for_polynomials_and_integers}.

\bibitem{vonZurGathenGerhard2013}
J.~von zur Gathen and J.~Gerhard,
\emph{Modern Computer Algebra}, third edition,
Cambridge University Press, 2013.
\newline
Book homepage:
\url{https://vonzurgathen.online/modern-computer-algebra/};
publisher excerpt:
\url{https://assets.cambridge.org/97811070/39032/excerpt/9781107039032_excerpt.pdf}.

\end{thebibliography}

\end{document}
